\documentclass[11pt]{amsart}

\usepackage[T1]{fontenc}
\usepackage[utf8]{inputenc}
\usepackage[margin=1.3in]{geometry}
\usepackage{amssymb,amsmath,amsfonts}
\usepackage{amsthm}
\usepackage{microtype}
\usepackage{nicefrac}
\usepackage[colorlinks=true,linkcolor=blue,citecolor=blue,urlcolor=blue]{hyperref}
\usepackage{enumitem}
\usepackage{mathtools}
\usepackage{seqsplit}

\setlength{\emergencystretch}{3em}

%%% Notation
\newcommand{\R}{\mathbb{R}}
\newcommand{\N}{\mathbb{N}}
\newcommand{\fint}{\mathop{\mathchoice%
  {\,\rlap{$\displaystyle-$}\!\int}%
  {\,\rlap{$\textstyle-$}\!\int}%
  {\,\rlap{$\scriptstyle-$}\!\int}%
  {\,\rlap{$\scriptscriptstyle-$}\!\int}}\nolimits}
\DeclareMathOperator*{\esssup}{ess\,sup}
\DeclareMathOperator*{\essinf}{ess\,inf}
\DeclareMathOperator{\dv}{div}
\DeclareMathOperator{\supp}{supp}
\DeclareMathOperator{\osc}{osc}
\DeclareMathOperator{\dist}{dist}
\newcommand{\eLpNormText}{\mathrm{eLpNorm}}
% Macros invented by translation agents — collected here so the document compiles
\newcommand{\Q}{\mathbb{Q}}
\newcommand{\Mst}{M_{\mathrm{st}}}
\newcommand{\Eext}{E_{\mathrm{ext}}}
\newcommand{\BMO}{\mathrm{BMO}}
\newcommand{\vol}{\mathrm{vol}}
\newcommand{\lam}{\lambda}
\newcommand{\Cpoinc}{C_{\mathrm{poinc}}}
\newcommand{\CJN}{C_{\mathrm{JN}}}
\newcommand{\loc}{\mathrm{loc}}
\newcommand{\tsupp}{\mathrm{tsupp}}
\newcommand{\unitBallDilate}{\mathrm{unitBallDilate}}
\newcommand{\unitBallShellFormula}{\mathrm{unitBallShellFormula}}
\newcommand{\smoothUnitBallExtensionApprox}{\mathrm{smoothUnitBallExtensionApprox}}
\newcommand{\smoothUnitBallExtensionGrad}{\mathrm{smoothUnitBallExtensionGrad}}
\newcommand{\tsupport}{\mathrm{tsupport}}
\newcommand{\sSup}{\mathrm{sSup}}
\newcommand{\sInf}{\mathrm{sInf}}
\newcommand{\IsSolution}{\mathrm{IsSolution}}
\newcommand{\IsSupersolution}{\mathrm{IsSupersolution}}
\newcommand{\CunitBallExtensionGrad}{C_{\mathrm{ext},\nabla}}
\newcommand{\CunitBallExtensionFun}{C_{\mathrm{ext}}}
\newcommand{\ENNReal}{\mathrm{ENNReal}}
\newcommand{\ENNRealofReal}{\mathrm{ENNReal.ofReal}}
\newcommand{\Fin}{\mathrm{Fin}}
\newcommand{\IsHomogeneousWeakSolution}{\mathrm{IsHomogeneousWeakSolution}}
\newcommand{\fderiv}{\mathrm{fderiv}}
\newcommand{\hw}{\mathrm{hw}}
\newcommand{\hwExt}{\mathrm{hwExt}}
\newcommand{\Lp}[1]{L^{#1}}
\newcommand{\Wkp}[2]{W^{#1,#2}}
\newcommand{\Wop}[1]{W_0^{1,#1}}
\newcommand{\Ball}[2]{B_{#1}(#2)}
\newcommand{\Bz}[1]{B_{#1}}
% Allow line breaks in lean names at underscores; use \small monospace
\newcommand{\leanname}[1]{%
  \begingroup
  \textnormal{\small\texttt{%
    \let\oldunderscore\_%
    \def\_{\oldunderscore\discretionary{}{}{}}%
    \seqsplit{#1}}}%
  \endgroup}
\emergencystretch=3em

%%% Theorem environments
\newtheorem{theorem}{Theorem}[section]
\newtheorem{proposition}[theorem]{Proposition}
\newtheorem{lemma}[theorem]{Lemma}
\newtheorem{corollary}[theorem]{Corollary}

\theoremstyle{definition}
\newtheorem{definition}[theorem]{Definition}

\theoremstyle{remark}
\newtheorem{remark}[theorem]{Remark}
\newtheorem{notation}[theorem]{Notation}

%%% Title
\title{Natural Language Translation of the De~Giorgi--Nash--Moser Lean Formalization}

\author{Machine-generated from Lean 4 source code}

\date{}

\begin{document}

\maketitle

\tableofcontents

%% ========================================================================
\section{Introduction}\label{sec:intro}
%% ========================================================================

This document is a complete natural-language translation of the Lean~4 formalization of the De~Giorgi--Nash--Moser regularity theory for uniformly elliptic divergence-form equations with bounded measurable coefficients. The formalization is completely sorry-free and built on top of Mathlib. The Lean source is available at \url{https://github.com/scottnarmstrong/DeGiorgi}.

The main results formalized are:
\begin{enumerate}[label=(\roman*)]
  \item \textbf{Local boundedness} of weak subsolutions via De~Giorgi iteration;
  \item \textbf{Weak Harnack inequality} for nonnegative weak supersolutions via Moser iteration;
  \item \textbf{Harnack inequality} for weak solutions;
  \item \textbf{Interior H\"older regularity} for weak solutions via Moser's method.
\end{enumerate}

Throughout, we work in dimension $d \ge 3$ with functions $u \colon \R^d \to \R$. We write $\Bz{r} = \{x \in \R^d : |x| < r\}$ for the open ball of radius~$r$ centered at the origin, and $\Ball{r}{x_0}$ for the ball centered at~$x_0$.

Each definition and theorem includes a reference to its Lean name in \texttt{typewriter font}, enabling cross-referencing with the source code.

\bigskip
\noindent\textbf{Notation.} We adopt the following conventions:
\begin{itemize}
  \item $\nabla u$ denotes the weak gradient of~$u$;
  \item $(u - k)_+ = \max(u - k, 0)$ denotes the positive part;
  \item $\|f\|_{\Lp{p}(\Omega)}$ denotes the $\Lp{p}$ norm over~$\Omega$;
  \item $\lambda, \Lambda > 0$ denote the ellipticity constants (lower and upper bounds);
  \item $|E|$ denotes the Lebesgue measure of a measurable set~$E$;
  \item $\fint_\Omega f\, dx = |\Omega|^{-1}\int_\Omega f\, dx$ denotes the average integral.
\end{itemize}



%% ========================================================================
\section{Sobolev Spaces and Weak Derivatives}\label{sec:sobolev}
%% ========================================================================

% Section 2: Sobolev Spaces and Weak Derivatives
% Translation of:
%   DeGiorgi/SobolevSpace.lean
%   DeGiorgi/SobolevSpace/WeakDerivatives.lean
%   DeGiorgi/SobolevSpace/Witnesses.lean
%   DeGiorgi/SobolevSpace/Approximation.lean
%   DeGiorgi/SobolevSpace/PositivePartPrelude.lean
%   DeGiorgi/WholeSpaceSobolev.lean

Throughout this section we fix a positive integer $d$ and write $E = \R^d$ (in
the Lean formalization, $E$ is the Euclidean space of dimension $d$, with the
canonical orthonormal basis $(e_i)_{i=1}^d$). All open sets $\Omega \subseteq E$
are equipped with Lebesgue measure. For a smooth scalar test function $\varphi
\in C_c^\infty(E)$, $i \in \{1,\ldots,d\}$ and $x \in E$ we abbreviate the
$i$-th partial derivative by $\partial_i \varphi(x) = (D\varphi(x))(e_i)$.

\subsection{Weak partial derivatives and weak gradients}

\begin{definition}[\leanname{DeGiorgi.HasWeakPartialDeriv}]\label{def:HasWeakPartialDeriv}
Let $\Omega \subseteq E$ be a set, $i \in \{1,\ldots,d\}$, and $f, g : E \to \R$.
We say that $g$ is the \emph{weak $i$-th partial derivative of $f$ on $\Omega$},
written $\partial_i^w f = g$ on $\Omega$, if for every test function
$\varphi \in C^\infty(E,\R)$ with compact support and $\supp \varphi \subseteq
\Omega$ we have
\begin{equation*}
  \int_\Omega f(x) \, \partial_i \varphi(x) \, dx
    = - \int_\Omega g(x) \, \varphi(x) \, dx.
\end{equation*}
\end{definition}

\begin{definition}[\leanname{DeGiorgi.HasWeakGrad}, \leanname{DeGiorgi.HasWeakDiv}]\label{def:HasWeakGrad}
A vector field $G : E \to E$ is a \emph{weak gradient} of $f$ on $\Omega$,
written $\nabla^w f = G$, if each component $G_i$ is the weak $i$-th partial
derivative of $f$ on $\Omega$.

A scalar function $g : E \to \R$ is a \emph{weak divergence} of a vector field
$F : E \to E$ on $\Omega$ if for every $\varphi \in C_c^\infty$ with $\supp
\varphi \subseteq \Omega$
\begin{equation*}
  \int_\Omega \sum_{i=1}^d F_i(x) \, \partial_i \varphi(x) \, dx
    = - \int_\Omega g(x)\, \varphi(x)\, dx.
\end{equation*}
\end{definition}

\begin{lemma}[Uniqueness; \leanname{DeGiorgi.HasWeakPartialDeriv.ae\_eq}]\label{lem:weak_partial_unique}
Let $\Omega \subseteq E$ be open, $i \in \{1,\ldots,d\}$ and $f, g_1, g_2 :
E \to \R$ with $g_1$ and $g_2$ locally integrable on $\Omega$. If both $g_1$
and $g_2$ are weak $i$-th partial derivatives of $f$ on $\Omega$, then
$g_1 = g_2$ almost everywhere on $\Omega$.
\end{lemma}

\begin{proof}
It suffices to show that $g_1 - g_2 = 0$ a.e.\ on $\Omega$. The function $g_1 -
g_2$ is locally integrable on $\Omega$, and for any $\varphi \in C_c^\infty(E)$
with $\supp \varphi \subseteq \Omega$ we have, subtracting the two test
identities,
\begin{equation*}
  \int_\Omega g_1(x)\,\varphi(x)\,dx = \int_\Omega g_2(x)\,\varphi(x)\,dx,
\end{equation*}
so $\int_\Omega (g_1 - g_2)(x)\,\varphi(x)\,dx = 0$. Since $\varphi$ vanishes
outside its support which lies in $\Omega$, this is the same as integrating over
the whole space $E$. By the standard fact that a locally integrable function
whose integral against every smooth compactly supported test function vanishes
is zero almost everywhere (\leanname{IsOpen.ae\_eq\_zero\_of\_integral\_contDiff\_smul\_eq\_zero}),
we conclude $g_1 = g_2$ a.e.\ on $\Omega$.
\end{proof}

\begin{lemma}[Classical $\Rightarrow$ weak derivative; \leanname{DeGiorgi.HasWeakPartialDeriv.of\_contDiff}]\label{lem:weak_of_contDiff}
Let $\Omega \subseteq E$ be open and $f \in C^1(E,\R)$. Then $\partial_i f$ is
the weak $i$-th partial derivative of $f$ on $\Omega$.
\end{lemma}

\begin{proof}
Take any $\varphi \in C_c^\infty(E)$ with $\supp\varphi \subseteq \Omega$. Both
$f \cdot \partial_i \varphi$ and $\partial_i f \cdot \varphi$ vanish outside
$\supp\varphi \subseteq \Omega$, so the corresponding integrals on $\Omega$ and
on $E$ coincide. The classical integration-by-parts formula
(\leanname{integral\_mul\_fderiv\_eq\_neg\_fderiv\_mul\_of\_integrable}) for
differentiable functions with compactly supported product gives
\begin{equation*}
  \int_E f(x)\,\partial_i \varphi(x)\,dx
    = - \int_E (\partial_i f)(x)\,\varphi(x)\,dx,
\end{equation*}
which is the weak-derivative identity on $\Omega$.
\end{proof}

\begin{lemma}[Product rule against a smooth factor; \leanname{DeGiorgi.HasWeakPartialDeriv.mul\_smooth}]\label{lem:weak_product_rule}
Let $\Omega \subseteq E$ be open, $i \in \{1,\ldots,d\}$, and let $f, g, \eta :
E \to \R$ with $\eta \in C^\infty$, $f$ and $g$ locally integrable on $\Omega$,
and $g$ a weak $i$-th partial derivative of $f$ on $\Omega$. Then
\begin{equation*}
  x \mapsto \eta(x)\,g(x) + (\partial_i \eta(x))\,f(x)
\end{equation*}
is the weak $i$-th partial derivative of $\eta f$ on $\Omega$.
\end{lemma}

\begin{proof}
Fix a test function $\varphi \in C_c^\infty(E)$ with $\supp\varphi \subseteq
\Omega$. Then $\eta \varphi$ is also smooth, compactly supported, with support
contained in $\Omega$. Applying the weak-derivative identity for $f$ to the test
function $\eta \varphi$ and using the smooth product rule
$\partial_i(\eta\varphi) = \eta\,\partial_i\varphi + \varphi\,\partial_i\eta$
gives
\begin{equation*}
  \int_\Omega f(x)\bigl(\eta(x)\,\partial_i\varphi(x) + \varphi(x)\,\partial_i\eta(x)\bigr)\,dx
    = -\int_\Omega g(x)\,\eta(x)\,\varphi(x)\,dx.
\end{equation*}
The local integrability of $f$ and $g$, together with continuity of $\eta$,
$\partial_i\eta$, $\varphi$, $\partial_i\varphi$ and compactness of supports,
ensures that all integrals are absolutely convergent and may be split. Splitting
the left-hand side and rearranging yields
\begin{align*}
  \int_\Omega \eta(x)\,f(x)\,\partial_i\varphi(x)\,dx
  &= -\int_\Omega g(x)\,\eta(x)\,\varphi(x)\,dx
     - \int_\Omega f(x)\,\varphi(x)\,\partial_i\eta(x)\,dx \\
  &= -\int_\Omega \bigl(\eta(x)g(x) + (\partial_i\eta(x))f(x)\bigr)\,\varphi(x)\,dx,
\end{align*}
which is the desired weak-derivative identity for $\eta f$.
\end{proof}

\begin{lemma}[Restriction to a sub-open set; \leanname{DeGiorgi.HasWeakPartialDeriv.restrict}]\label{lem:weak_restrict}
Let $\Omega' \subseteq \Omega \subseteq E$ with $\Omega'$ open. If $g$ is the
weak $i$-th partial derivative of $f$ on $\Omega$, then it is also the weak
$i$-th partial derivative of $f$ on $\Omega'$.
\end{lemma}

\begin{proof}
Let $\varphi \in C_c^\infty(E)$ with $\supp\varphi \subseteq \Omega'$. Then a
fortiori $\supp\varphi \subseteq \Omega$, so applying the hypothesis on
$\Omega$ gives the integral identity over $\Omega$. Outside $\Omega'$ the test
function $\varphi$ and the field $\partial_i\varphi$ both vanish (the latter
because $\varphi$ is locally constant equal to zero), so the integrals over
$\Omega$ and $\Omega'$ of $f\,\partial_i\varphi$ and $g\,\varphi$ agree. The
result follows.
\end{proof}

\begin{lemma}[$L^p$-closure of weak partial derivatives; \leanname{DeGiorgi.HasWeakPartialDeriv.of\_eLpNormApprox\_p}]\label{lem:weak_lp_closure}
Let $\Omega \subseteq E$ be open and $1 < p < \infty$. Suppose $f, g \in
\Lp{p}(\Omega)$ and there are sequences $(\psi_n)$ and $(g_{\psi_n})$ such that:
\begin{itemize}
  \item $g_{\psi_n}$ is the weak $i$-th partial derivative of $\psi_n$ on
        $\Omega$ for every $n$;
  \item $\psi_n - f$ and $g_{\psi_n} - g$ lie in $\Lp{p}(\Omega)$ and tend to
        $0$ in $\Lp{p}(\Omega)$ as $n \to \infty$.
\end{itemize}
Then $g$ is the weak $i$-th partial derivative of $f$ on $\Omega$.
\end{lemma}

\begin{proof}
Fix $\varphi \in C_c^\infty(E)$ with $\supp\varphi \subseteq \Omega$. Let $q$
denote the H{\"o}lder conjugate of $p$, so that $1 < q < \infty$ and $1/p + 1/q
= 1$. Both $\varphi$ and $\partial_i\varphi$ are continuous, compactly
supported, and hence belong to $\Lp{q}(\Omega)$.

By H{\"o}lder's inequality on $(\Omega, dx)$,
\begin{align*}
  \left| \int_\Omega \partial_i\varphi(x) (\psi_n(x) - f(x))\,dx \right|
  &\le \|\partial_i\varphi\|_{\Lp{q}(\Omega)}\,\|\psi_n - f\|_{\Lp{p}(\Omega)},\\
  \left| \int_\Omega \varphi(x) (g_{\psi_n}(x) - g(x))\,dx \right|
  &\le \|\varphi\|_{\Lp{q}(\Omega)}\,\|g_{\psi_n} - g\|_{\Lp{p}(\Omega)},
\end{align*}
both of which tend to zero by hypothesis. Hence
\begin{align*}
  \int_\Omega \psi_n(x)\,\partial_i\varphi(x)\,dx
    &\;\to\; \int_\Omega f(x)\,\partial_i\varphi(x)\,dx, \\
  \int_\Omega g_{\psi_n}(x)\,\varphi(x)\,dx
    &\;\to\; \int_\Omega g(x)\,\varphi(x)\,dx.
\end{align*}
By assumption, for every $n$
\begin{equation*}
  \int_\Omega \psi_n(x)\,\partial_i\varphi(x)\,dx
    = -\int_\Omega g_{\psi_n}(x)\,\varphi(x)\,dx,
\end{equation*}
and passing to the limit yields
\begin{equation*}
  \int_\Omega f(x)\,\partial_i\varphi(x)\,dx
    = -\int_\Omega g(x)\,\varphi(x)\,dx,
\end{equation*}
which is the weak-derivative identity for $f$ and $g$.
\end{proof}

\subsection{Sobolev membership and weak-gradient witnesses}

\begin{definition}[\leanname{DeGiorgi.MemW1p}]\label{def:MemW1p}
Let $p \in [1,\infty]$, $\Omega \subseteq E$ a set, and $f : E \to \R$. We say
that $f \in \Wkp{1}{p}(\Omega)$ if $f \in \Lp{p}(\Omega)$ and for every $i \in
\{1,\ldots,d\}$ there exists $g_i \in \Lp{p}(\Omega)$ such that $g_i$ is the
weak $i$-th partial derivative of $f$ on $\Omega$.
\end{definition}

\begin{definition}[\leanname{DeGiorgi.MemW1pWitness}]\label{def:MemW1pWitness}
A \emph{$\Wkp{1}{p}$-witness} for $f$ on $\Omega$ is a tuple consisting of a
proof that $f \in \Lp{p}(\Omega)$, a vector field $G : E \to E$ (the weak
gradient field) such that each component $x \mapsto G_i(x)$ lies in
$\Lp{p}(\Omega)$, together with a proof that $G$ is a weak gradient of $f$ on
$\Omega$ (Definition~\ref{def:HasWeakGrad}).
\end{definition}

The Hilbert-space case $p = 2$ is also denoted $H^1(\Omega) := W^{1,2}(\Omega)$
in the formalization (\leanname{DeGiorgi.MemH1}).

\begin{definition}[\leanname{DeGiorgi.MemW01p}]\label{def:MemW01p}
We say that $f \in \Wop{p}(\Omega)$ if $f \in \Wkp{1}{p}(\Omega)$ and there
exist a $\Wkp{1}{p}$-witness $(G)$ for $f$ on $\Omega$ and a sequence
$(\varphi_n)_{n\in\N}$ of smooth compactly supported functions with
$\supp\varphi_n \subseteq \Omega$ such that
\begin{equation*}
  \|\varphi_n - f\|_{\Lp{p}(\Omega)} \to 0
  \quad\text{and}\quad
  \|\partial_i \varphi_n - G_i\|_{\Lp{p}(\Omega)} \to 0
  \;\text{ for every } i.
\end{equation*}
We write $H_0^1(\Omega) := W_0^{1,2}(\Omega)$ (\leanname{DeGiorgi.MemH01}).
\end{definition}

\begin{remark}[Existence of an explicit witness; \leanname{DeGiorgi.MemW1p.someWitness}]
Given $f \in W^{1,p}(\Omega)$, axiom of choice produces a witness as in
Definition~\ref{def:MemW1pWitness}: choose for each coordinate $i$ a function
$g_i$ as in Definition~\ref{def:MemW1p}, and assemble them into the vector
field $G(x) := (g_1(x),\ldots,g_d(x))$.
\end{remark}

\begin{lemma}[Forgetting the witness]
\label{lem:forgetting-witness}
(\leanname{DeGiorgi.MemW1pWitness.memW1p},
\leanname{DeGiorgi.MemW01p.memW1p})\\
Every witness gives rise to membership: a $\Wkp{1}{p}$-witness yields
$f \in W^{1,p}(\Omega)$, and an element of $\Wop{p}(\Omega)$ is in particular
in $\Wkp{1}{p}(\Omega)$.
\end{lemma}

\begin{proof}
Direct from the definitions: take $g_i := G_i$.
\end{proof}

\begin{lemma}[Sum of witnesses; \leanname{DeGiorgi.MemW1pWitness.add}]\label{lem:witness_add}
For $u, v \in H^1(\Omega)$ with witnesses $G_u$ and $G_v$, the function $u + v$
admits the witness $G_u + G_v$.
\end{lemma}

\begin{proof}
Membership of $u+v$ in $\Lp{2}(\Omega)$ and of $G_u + G_v$ in $\Lp{2}(\Omega)$
in each coordinate follows from $\Lp{2}$ being a vector space. To verify the
weak-gradient identity, take $\varphi \in C_c^\infty(E)$ with $\supp\varphi
\subseteq \Omega$. The witnesses $G_u$ and $G_v$ give
\begin{align*}
  \int_\Omega u(x)\,\partial_i\varphi(x)\,dx
    &= -\int_\Omega (G_u)_i(x)\,\varphi(x)\,dx, \\
  \int_\Omega v(x)\,\partial_i\varphi(x)\,dx
    &= -\int_\Omega (G_v)_i(x)\,\varphi(x)\,dx.
\end{align*}
Adding these identities, using local integrability and the linearity of the
integral, yields the weak-gradient identity for $u + v$ with weak gradient
$G_u + G_v$.
\end{proof}

\begin{lemma}[Restriction of witnesses; \leanname{DeGiorgi.MemW1pWitness.restrict}]
Let $\Omega' \subseteq \Omega$ with $\Omega'$ open. Any
$\Wkp{1}{p}$-witness on $\Omega$ restricts to a $\Wkp{1}{p}$-witness on
$\Omega'$, with the same gradient field.
\end{lemma}

\begin{proof}
$\Lp{p}$ membership is monotone in the measure (here the volume measure
restricted to a smaller set), and Lemma~\ref{lem:weak_restrict} gives the weak
gradient identity on $\Omega'$.
\end{proof}

\begin{lemma}[Scalar multiple of a witness; \leanname{DeGiorgi.MemW1pWitness.smul}]
For $u \in H^1(\Omega)$ with witness $G$ and $c \in \R$, the function $cu$
admits the witness $cG$.
\end{lemma}

\begin{proof}
The $\Lp{2}$ memberships are immediate. For the test-function identity, fix
$\varphi$ as before; using linearity and the witness for $u$,
\begin{equation*}
  \int_\Omega (cu)(x)\partial_i\varphi(x)\,dx
    = c\int_\Omega u(x)\partial_i\varphi(x)\,dx
    = -c\int_\Omega G_i(x)\varphi(x)\,dx
    = -\int_\Omega (cG)_i(x)\varphi(x)\,dx. \qedhere
\end{equation*}
\end{proof}

\begin{proposition}[Multiplying by a bounded smooth scalar; \leanname{DeGiorgi.MemW1pWitness.mul\_smooth\_bounded\_p}]\label{prop:witness_mul_smooth_bounded}
Let $1 \le p \le \infty$, $\Omega$ open, $u$ a function with $\Wkp{1}{p}$-witness
$G$ on $\Omega$, and $\eta \in C^\infty(E)$ satisfying
$|\eta(x)| \le C_0$ and $\|D\eta(x)\| \le C_1$ for all $x$ and constants
$C_0, C_1 \ge 0$. Then $\eta u$ admits the $\Wkp{1}{p}$-witness whose gradient
field is
\begin{equation*}
  x \mapsto \eta(x)\,G(x) + \bigl(\partial_1\eta(x)\,u(x),\ldots,\partial_d\eta(x)\,u(x)\bigr).
\end{equation*}
\end{proposition}

\begin{proof}
The $\Lp{p}$ bound on $\eta u$ follows from the pointwise inequality
$|\eta(x)u(x)| \le C_0 |u(x)|$ and $u \in \Lp{p}(\Omega)$. For each component
of the displayed gradient field:
\begin{itemize}
  \item $|\eta(x)G_i(x)| \le C_0 |G_i(x)|$, and $G_i \in \Lp{p}(\Omega)$;
  \item $|\partial_i\eta(x)\,u(x)| \le \|D\eta(x)\|\,|u(x)| \le C_1 |u(x)|$,
        and $u \in \Lp{p}(\Omega)$.
\end{itemize}
By Minkowski, the sum lies in $\Lp{p}(\Omega)$. The weak-gradient identity is
the product rule (Lemma~\ref{lem:weak_product_rule}) applied componentwise to
the weak partial derivative $G_i$ of $u$ and the smooth factor $\eta$.
\end{proof}

\begin{lemma}[Smooth compactly supported functions are $\Wkp{1}{p}$]
\label{lem:witness_smooth_cpt}
(\leanname{DeGiorgi.MemW1pWitness.of\_contDiff\_hasCompactSupport})\\
Let $f \in C^\infty_c(E,\R)$. Then $f$ admits a canonical $\Wkp{1}{p}$-witness
on the whole space $E$, with weak gradient $\nabla^w f = \nabla f$ given by the
classical gradient field.
\end{lemma}

\begin{proof}
Smooth compactly supported functions are integrable to any power, so $f \in
\Lp{p}(E)$ and likewise each classical partial derivative. By
Lemma~\ref{lem:weak_of_contDiff} the classical partial derivatives are weak
partial derivatives.
\end{proof}

\begin{lemma}[Norm of weak gradient lies in $\Lp{p}$;
\leanname{DeGiorgi.MemW1pWitness.weakGrad\_memLp},
\leanname{DeGiorgi.MemW1pWitness.weakGrad\_norm\_memLp}]\label{lem:weak_grad_norm_memLp}
Let $G$ be the weak gradient field of a $\Wkp{1}{p}$-witness on $\Omega$. Then
$G \in \Lp{p}(\Omega; E)$, and in particular $\|G(\cdot)\| \in \Lp{p}(\Omega)$.
\end{lemma}

\begin{proof}
Each component $G_i$ lies in $\Lp{p}(\Omega)$ by definition of a witness. For
finite-dimensional Euclidean targets the $\Lp{p}$ membership of a vector field
is equivalent to the $\Lp{p}$ membership of each component
(\leanname{MemLp.of\_eval\_piLp}), so $G \in \Lp{p}(\Omega; E)$. Taking norms
preserves $\Lp{p}$ membership.
\end{proof}

\begin{proposition}[$H_0^1$ is a vector space; \leanname{DeGiorgi.MemW01p.add}, \leanname{DeGiorgi.MemW01p.smul}, \leanname{DeGiorgi.MemW01p.sub}]\label{prop:H01_vector_space}
The space $H_0^1(\Omega)$ is closed under addition, scalar multiplication, and
subtraction.
\end{proposition}

\begin{proof}
For addition: Lemma~\ref{lem:witness_add} provides a witness for the sum.
Approximating sequences add componentwise; by the triangle inequality for
$\Lp{2}$ on differences,
\begin{equation*}
  \|\varphi_n^u + \varphi_n^v - (u+v)\|_{\Lp{2}(\Omega)}
  \le \|\varphi_n^u - u\|_{\Lp{2}(\Omega)} + \|\varphi_n^v - v\|_{\Lp{2}(\Omega)}
  \;\to\; 0,
\end{equation*}
and similarly for the gradients (using the linearity of the classical
derivative on smooth functions).

For scalar multiplication by $c$: rescale the approximating sequence by $c$ and
use the homogeneity of the $\Lp{2}$ norm,
$\|c\varphi_n - c u\|_{\Lp{2}} = |c|\,\|\varphi_n - u\|_{\Lp{2}}$.

Subtraction follows by writing $u - v = u + (-1)v$.
\end{proof}

\begin{lemma}[Smooth compactly supported $\Rightarrow$ $\Wop{p}$;
\leanname{DeGiorgi.memW01p\_of\_contDiff\_hasCompactSupport},
\leanname{DeGiorgi.memW01p\_of\_contDiff\_hasCompactSupport\_subset}]\label{lem:memW01p_of_smooth_cpt}
If $f \in C_c^\infty(E)$ then $f \in \Wop{p}(E)$. More generally, if $\Omega
\subseteq E$ is open and $f \in C_c^\infty(E)$ with $\supp f \subseteq \Omega$,
then $f \in \Wop{p}(\Omega)$.
\end{lemma}

\begin{proof}
Take the constant approximating sequence $\varphi_n := f$. The $\Lp{p}$
differences are identically zero, so trivially converge to $0$, and
Lemma~\ref{lem:witness_smooth_cpt} provides the witness.
\end{proof}

\subsection{Cutoffs and zero-extension}

\begin{lemma}[Smooth cutoff between a compact set and an open neighborhood;
\leanname{DeGiorgi.exists\_smooth\_cutoff} (private)]\label{lem:smooth_cutoff}
Let $K \subseteq \Omega \subseteq E$ with $K$ compact and $\Omega$ open. Then
there exists $\eta \in C^\infty(E)$ with compact support such that
$0 \le \eta \le 1$, $\eta = 1$ on $K$, and $\supp\eta \subseteq \Omega$.
\end{lemma}

\begin{proof}
Since $K$ is compact and contained in the open set $\Omega$, there exists
$\delta > 0$ such that the closed $\delta$-thickening of $K$ is contained in
$\Omega$. The classical existence of a smooth bump (in Lean,
\leanname{exists\_contMDiff\_support\_eq\_eq\_one\_iff}) yields a smooth $\eta$
with $\eta(x) = 1$ on $K$, support equal to the open $\delta$-thickening of $K$,
and range in $[0,1]$. Its support is then contained in the closed
$\delta$-thickening of $K$ and hence in $\Omega$, and it is automatically
compactly supported.
\end{proof}

\begin{lemma}[Vanishing outside the support; \leanname{DeGiorgi.zero\_outside\_of\_tsupport\_subset}, \leanname{DeGiorgi.fderiv\_apply\_zero\_outside\_of\_tsupport\_subset}]\label{lem:zero_outside_tsupport}
If $\supp f \subseteq \Omega$ and $x \notin \Omega$, then $f(x) = 0$. If
furthermore $f \in C^\infty$ then $\partial_i f(x) = 0$ as well.
\end{lemma}

\begin{proof}
If $f(x) \ne 0$ then $x$ lies in the support of $f$, hence in $\Omega$,
contradicting $x \notin \Omega$. The same argument applied to $\partial_i f$
(whose support is contained in that of $f$) handles the second statement.
\end{proof}

\begin{theorem}[Zero-extension of a $\Wop{p}$ function; \leanname{DeGiorgi.zeroExtend\_memW1pWitness\_p}]\label{thm:zero_extend}
Let $\Omega \subseteq E$ be open, $1 < p < \infty$, and $v \in \Wop{p}(\Omega)$
with $\Wkp{1}{p}$-witness $G_v$; assume in addition that $v \in W_0^{1,p}(\Omega)$. Define $\tilde v := \mathbf{1}_\Omega \cdot v$
on $E$ and let $\tilde G$ be the vector field whose components are
$\mathbf{1}_\Omega \cdot (G_v)_i$. Then $\tilde G$ is a $\Wkp{1}{p}$-witness
for $\tilde v$ on the whole space $E$.
\end{theorem}

\begin{proof}
The function $v$ admits an approximating sequence $(\varphi_n)$ of smooth
compactly supported test functions with $\supp\varphi_n \subseteq \Omega$, and
$\varphi_n \to v$, $\partial_i\varphi_n \to (G_v)_i$ in $\Lp{p}(\Omega)$. By
Lemma~\ref{lem:weak_partial_unique} the gradient field $G_v$ obtained from the
arbitrary witness coincides almost everywhere on $\Omega$ with the one from the
$\Wop{p}$ data.

The $\Lp{p}$ membership of $\tilde v$ on $E$ is equivalent to the $\Lp{p}$
membership of $v$ on $\Omega$, via
\leanname{MeasureTheory.memLp\_indicator\_iff\_restrict}, and similarly for each
component of $\tilde G$.

To check that $\tilde G$ is a weak gradient on $E$, we apply
Lemma~\ref{lem:weak_lp_closure} (closure of weak partial derivatives in
$\Lp{p}$, on $\Omega = E$) with the same approximating sequence
$(\varphi_n)$ and the classical gradients $\partial_i\varphi_n$. The
key observations are:
\begin{itemize}
  \item Outside $\Omega$ both $\varphi_n$ and $\partial_i\varphi_n$ vanish (by
        Lemma~\ref{lem:zero_outside_tsupport}); hence
        $\varphi_n - \tilde v = \mathbf{1}_\Omega(\varphi_n - v)$ and
        $\partial_i\varphi_n - \tilde G_i =
         \mathbf{1}_\Omega(\partial_i\varphi_n - (G_v)_i)$ as functions on $E$.
  \item The $\Lp{p}$-norm of an indicator of a measurable set equals the
        $\Lp{p}$-norm computed on the restricted measure
        (\leanname{MeasureTheory.eLpNorm\_indicator\_eq\_eLpNorm\_restrict}),
        so the convergence on $\Omega$ implies convergence on $E$.
\end{itemize}
The hypotheses of Lemma~\ref{lem:weak_lp_closure} are then satisfied, and
$\tilde G_i$ is the weak $i$-th partial derivative of $\tilde v$ on $E$.
\end{proof}

\subsection{Convolution smoothing}

For $n \in \N$, define a $C^\infty$ bump function on $E$ centered at $0$ with
inner radius $r_{\text{in}}^{(n)} := \frac{1}{2(n+1)}$ and outer radius
$r_{\text{out}}^{(n)} := \frac{1}{n+1}$ (\leanname{DeGiorgi.shrinkingBump} (private)).
We denote the corresponding mass-one normalized bump by $\rho_n :=
(\text{shrinkingBump}_n).\text{normed}$, so $\int_E \rho_n = 1$, $\rho_n \ge 0$,
and $\supp \rho_n \subseteq \Bz{r_{\text{out}}^{(n)}}$. The radii satisfy
\begin{equation*}
  r_{\text{out}}^{(n)} = \frac{1}{n+1} \;\to\; 0
  \quad\text{as } n \to \infty
\end{equation*}
(\leanname{DeGiorgi.tendsto\_shrinkingBump\_rOut} (private)).

\begin{lemma}[Support of a convolution; \leanname{DeGiorgi.tsupport\_convolution\_subset\_cthickening} (private), \leanname{DeGiorgi.tsupport\_normed\_convolution\_subset\_cthickening} (private)]\label{lem:tsupport_conv}
If $\supp \kappa \subseteq \overline{\Bz{r}}$ and $u$ has compact support, then
\begin{equation*}
  \supp(\kappa * u) \;\subseteq\; \{x \in E : \dist(x, \supp u) \le r\}.
\end{equation*}
In particular, $\supp(\rho_n * u) \subseteq \{x : \dist(x,\supp u) \le
r_{\text{out}}^{(n)}\}$.
\end{lemma}

\begin{proof}
Pointwise, $(\kappa * u)(x) = \int \kappa(y) u(x - y)\,dy$ vanishes unless
there exists $y \in \supp\kappa$ and $z \in \supp u$ with $x = y + z$, hence
$\dist(x, z) = \|y\| \le r$. Closing under the topology, the closure is
contained in the closed $r$-thickening, which is closed.
\end{proof}

\begin{proposition}[$\Lp{p}$-contraction of normed convolution;
\leanname{DeGiorgi.eLpNorm\_normedConvolution\_le} (private)]\label{prop:Lp_contraction_conv}
Let $\rho \in C_c^\infty(E)$ with $\rho \ge 0$ and $\int_E \rho = 1$, and let
$1 \le p < \infty$. Then for every $f \in \Lp{p}(E)$,
\begin{equation*}
  \|\rho * f\|_{\Lp{p}(E)} \le \|f\|_{\Lp{p}(E)}.
\end{equation*}
\end{proposition}

\begin{proof}
Define the probability measure $\nu$ on $E$ by $d\nu = \rho\,d\,\mathrm{vol}$,
which is a probability measure since $\rho \ge 0$ and $\int \rho = 1$. For each
$x \in E$,
\begin{equation*}
  (\rho * f)(x) = \int_E f(x - y)\,d\nu(y).
\end{equation*}
By the integral form of Jensen's inequality applied to the convex function
$t \mapsto t^p$ on $[0,\infty)$ (\leanname{ConvexOn.map\_integral\_le}),
\begin{equation*}
  |(\rho * f)(x)|^p
    \le \left(\int_E |f(x-y)|\,d\nu(y)\right)^p
    \le \int_E |f(x-y)|^p \,d\nu(y)
    = (\rho * |f|^p)(x).
\end{equation*}
Integrating in $x$ and using Fubini together with $\int \rho = 1$,
\begin{equation*}
  \int_E |(\rho * f)(x)|^p\,dx \le \int_E (\rho * |f|^p)(x)\,dx
    = \left(\int_E \rho\right)\left(\int_E |f|^p\right)
    = \|f\|_{\Lp{p}}^p.
\end{equation*}
Taking $p$-th roots gives the claim.
\end{proof}

\begin{corollary}[Convolution is linear on differences;
\leanname{DeGiorgi.normedConvolution\_sub\_eq} (private),
\leanname{DeGiorgi.eLpNorm\_normedConvolution\_sub\_le} (private)]\label{cor:conv_sub}
Under the assumptions of Proposition~\ref{prop:Lp_contraction_conv} we also
have, for $f, g \in \Lp{p}(E)$,
\begin{equation*}
  \|\rho * f - \rho * g\|_{\Lp{p}(E)} = \|\rho * (f - g)\|_{\Lp{p}(E)}
   \le \|f - g\|_{\Lp{p}(E)}.
\end{equation*}
\end{corollary}

\begin{proof}
Linearity of $f \mapsto \rho * f$ on locally integrable functions yields the
equality, after which Proposition~\ref{prop:Lp_contraction_conv} applies to
$f - g$.
\end{proof}

\begin{lemma}[Convolution by a shrinking bump approximates a continuous function;
\leanname{DeGiorgi.tendsto\_eLpNorm\_normedConvolution\_sub\_of\_continuous} (private)]\label{lem:conv_approx_cont}
Let $g \in C_c(E)$ and $1 \le p < \infty$. Then
\begin{equation*}
  \|\rho_n * g - g\|_{\Lp{p}(E)} \;\to\; 0 \quad\text{as } n \to \infty.
\end{equation*}
\end{lemma}

\begin{proof}
Choose $S := \{x : \dist(x, \supp g) \le 1\}$ and $T := \{x : \dist(x, \supp g)
\le 2\}$, both compact, with $S \subseteq T$. The function $g$ is uniformly
continuous on the compact set $T$. Fix $\varepsilon > 0$ and let $\varepsilon_0
:= \varepsilon / (\mathrm{vol}(S)^{1/p} + 1)$. By uniform continuity, there
exists $\delta_0 > 0$ such that $|g(y) - g(x)| < \varepsilon_0$ whenever
$x, y \in T$ with $\dist(x, y) < \delta_0$. Set $\delta := \min(\delta_0, 1)$.

For $n$ large enough that $r_{\text{out}}^{(n)} < \delta$, we have $\supp(\rho_n
* g) \subseteq S$ (Lemma~\ref{lem:tsupport_conv}, since
$r_{\text{out}}^{(n)} \le 1$). Pointwise, for $x \in S$,
\begin{equation*}
  |(\rho_n * g)(x) - g(x)|
    = \left|\int_E \rho_n(y) (g(x - y) - g(x))\,dy\right|
    \le \varepsilon_0,
\end{equation*}
since $\rho_n$ is supported in $\Bz{r_{\text{out}}^{(n)}} \subseteq \Bz{\delta}$
and both $x, x - y \in T$ for $y \in \supp\rho_n$. For $x \notin S$ both
$(\rho_n * g)(x)$ and $g(x)$ vanish.

The bound $\|\rho_n * g - g\|_{\Lp{p}(E)} \le \varepsilon_0
\,\mathrm{vol}(S)^{1/p}$ then follows from
\leanname{eLpNorm\_sub\_le\_of\_dist\_bdd}, and the choice of $\varepsilon_0$
gives this is less than $\varepsilon$.
\end{proof}

\begin{theorem}[Convolution by a shrinking bump approximates an $\Lp{p}$ function;
\leanname{DeGiorgi.tendsto\_eLpNorm\_normedConvolution\_sub} (private)]\label{thm:conv_approx_Lp}
Let $1 \le p < \infty$ and $f \in \Lp{p}(E)$. Then
\begin{equation*}
  \|\rho_n * f - f\|_{\Lp{p}(E)} \;\to\; 0 \quad\text{as } n \to \infty.
\end{equation*}
\end{theorem}

\begin{proof}
By density of smooth compactly supported functions in $\Lp{p}$ (Mathlib's
\leanname{MemLp.exist\_eLpNorm\_sub\_le}), for any $\varepsilon > 0$ there
exists $g \in C_c^\infty(E)$ with $\|f - g\|_{\Lp{p}} < \varepsilon/3$. Then by
the triangle inequality
\begin{equation*}
  \|\rho_n * f - f\|_{\Lp{p}} \le
  \|\rho_n * f - \rho_n * g\|_{\Lp{p}} + \|\rho_n * g - g\|_{\Lp{p}} +
  \|g - f\|_{\Lp{p}}.
\end{equation*}
By Corollary~\ref{cor:conv_sub} the first term is at most $\|f - g\|_{\Lp{p}} <
\varepsilon/3$. By Lemma~\ref{lem:conv_approx_cont} the middle term tends to
$0$, so for $n$ large it is less than $\varepsilon/3$. The last term is also
less than $\varepsilon/3$. Hence $\|\rho_n * f - f\|_{\Lp{p}} < \varepsilon$
for all sufficiently large $n$, proving convergence to zero.
\end{proof}

\subsection{Uniqueness of the weak gradient and global smooth approximation}

\begin{theorem}[Uniqueness of the weak gradient;
\leanname{DeGiorgi.MemW1pWitness.ae\_eq\_p}, \leanname{DeGiorgi.MemW1pWitness.ae\_eq}]\label{thm:weak_grad_unique}
Let $\Omega \subseteq E$ be open, $1 \le p < \infty$, and $f$ a function with
two $\Wkp{1}{p}$-witnesses on $\Omega$ giving rise to weak gradient fields
$G_1, G_2$. Then $G_1 = G_2$ a.e.\ on $\Omega$. The same holds for $H^1$
witnesses ($p = 2$).
\end{theorem}

\begin{proof}
Apply Lemma~\ref{lem:weak_partial_unique} componentwise to $(G_1)_i$ and
$(G_2)_i$, then assemble: a.e.\ all components agree, so the vectors
themselves agree a.e.
\end{proof}

\begin{lemma}[Convolution differentiates against the weak partial derivative;
\leanname{DeGiorgi.convolution\_fderiv\_eq\_convolution\_weakPartial\_univ} (private)]\label{lem:conv_fderiv_weak_partial}
Let $u \in L^1_{loc}(E)$ and let $g$ be the weak $i$-th partial derivative of
$u$ on $E$. Then for every $\varphi \in C_c^\infty(E)$ and every $x \in E$,
\begin{equation*}
  \bigl((\partial_i\varphi) * u\bigr)(x) = (\varphi * g)(x).
\end{equation*}
\end{lemma}

\begin{proof}
For fixed $x$, define the test function $\psi(t) := \varphi(x - t)$, which is
smooth and compactly supported. A direct computation shows
$(\partial_i \psi)(t) = -(\partial_i\varphi)(x - t)$ (chain rule with the
sign-flipping translation). The weak-derivative identity for $u$ tested against
$\psi$ then reads
\begin{equation*}
  \int_E u(t)\,(\partial_i\psi)(t)\,dt = -\int_E g(t)\,\psi(t)\,dt,
\end{equation*}
i.e.,
\begin{equation*}
  -\int_E u(t)\,(\partial_i\varphi)(x - t)\,dt = -\int_E g(t)\,\varphi(x - t)\,dt.
\end{equation*}
Cancelling the minus signs and rewriting both sides as convolutions yields the
claim.
\end{proof}

\begin{theorem}[Smooth compactly supported approximation on $E$ for compactly
supported $\Wkp{1}{p}$ functions;
\leanname{DeGiorgi.exists\_smooth\_compactSupport\_W1p\_approx\_univ}]\label{thm:smooth_approx_univ}
Let $1 < p < \infty$, and let $u : E \to \R$ be compactly supported with a
$\Wkp{1}{p}$-witness on $E$ with weak gradient $G$. Then there exists a sequence
$(\varphi_n)$ of smooth compactly supported functions on $E$ with
\begin{align*}
  \supp \varphi_n &\subseteq \{x : \dist(x, \supp u) \le r_{\text{out}}^{(n)}\}, \\
  \|\varphi_n - u\|_{\Lp{p}(E)} &\to 0, \\
  \|\partial_i\varphi_n - G_i\|_{\Lp{p}(E)} &\to 0 \quad\text{for every } i.
\end{align*}
\end{theorem}

\begin{proof}
Define $\varphi_n := \rho_n * u$. Since $u$ is compactly supported and $\rho_n
\in C_c^\infty(E)$, $\varphi_n$ is smooth and has support contained in the
$r_{\text{out}}^{(n)}$-thickening of $\supp u$ (Lemma~\ref{lem:tsupport_conv}).
The $\Lp{p}$ approximation $\varphi_n \to u$ is exactly
Theorem~\ref{thm:conv_approx_Lp} applied to $u$.

For the gradients, the classical fact
\[
  \partial_i (\rho_n * u) = (\partial_i \rho_n) * u
\]
(Mathlib's
\leanname{HasCompactSupport.hasFDerivAt\_convolution\_left})
combined with Lemma~\ref{lem:conv_fderiv_weak_partial} (with $u$ on the right
and the smooth $\rho_n$ on the left) yields
\begin{equation*}
  \partial_i \varphi_n = (\partial_i\rho_n) * u = \rho_n * G_i,
\end{equation*}
since $G_i$ is the weak partial derivative of $u$ on $E$. Applying
Theorem~\ref{thm:conv_approx_Lp} to $G_i \in \Lp{p}(E)$ then gives
$\partial_i\varphi_n - G_i = \rho_n * G_i - G_i \to 0$ in $\Lp{p}(E)$.
\end{proof}

\begin{theorem}[Smooth approximation for a witness supported in a compact subset;
\leanname{DeGiorgi.exists\_smooth\_W1p\_approx\_of\_supportedWitness}]\label{thm:smooth_approx_local}
Let $\Omega$ be open, $1 < p < \infty$, and let $u$ have a $\Wkp{1}{p}$-witness
on $\Omega$ with gradient field $G$. Suppose $K \subseteq \Omega$ is compact
with $\supp u \subseteq K$ and $\supp G_i \subseteq K$ for every $i$. Then
there exists a sequence $(\varphi_n)$ of smooth compactly supported functions
with $\supp\varphi_n \subseteq \Omega$ such that
\begin{align*}
  \|\varphi_n - u\|_{\Lp{p}(\Omega)} &\to 0, \\
  \|\partial_i\varphi_n - G_i\|_{\Lp{p}(\Omega)} &\to 0 \quad\text{for every } i.
\end{align*}
\end{theorem}

\begin{proof}
Choose $\delta > 0$ so that the closed $\delta$-thickening $K' := \{x :
\dist(x,K) \le \delta\}$ is contained in $\Omega$
(\leanname{IsCompact.exists\_cthickening\_subset\_open}), and a smooth cutoff
$\chi$ as in Lemma~\ref{lem:smooth_cutoff} satisfying $\chi = 1$ on $K'$ and
$\supp \chi \subseteq \Omega$.

We promote the local witness on $\Omega$ to a witness on the whole space $E$
with the same gradient field. The $\Lp{p}$-membership extends to all of $E$
because $u$ and each $G_i$ vanish outside $K \subseteq \Omega$
(Lemma~\ref{lem:zero_outside_tsupport}). To verify the weak-gradient identity
on $E$, take any test function $\varphi \in C_c^\infty(E)$ and consider $\psi
:= \chi \varphi$. Then $\psi \in C_c^\infty$ and $\supp \psi \subseteq \Omega$,
so the witness on $\Omega$ gives
\begin{equation*}
  \int_\Omega u(x)\,\partial_i\psi(x)\,dx = -\int_\Omega G_i(x)\,\psi(x)\,dx.
\end{equation*}
On $K$ we have $\chi = 1$ and $\partial_i\chi = 0$ (because $\chi$ is constant
equal to $1$ on the larger neighborhood $K'$ of $K$); since $u$ and $G_i$ are
both supported in $K$, this gives pointwise on $E$
\begin{equation*}
  u(x)\,\partial_i\psi(x) = u(x)\,\partial_i\varphi(x), \qquad
  G_i(x)\,\psi(x) = G_i(x)\,\varphi(x),
\end{equation*}
and the integral identity transfers from $\psi$ to $\varphi$.

Now apply Theorem~\ref{thm:smooth_approx_univ} on the whole space to get a
sequence $(\widetilde\varphi_m)$ with
\begin{equation*}
  \supp \widetilde\varphi_m \subseteq \{x : \dist(x, \supp u) \le r_{\text{out}}^{(m)}\}.
\end{equation*}
For $m$ large enough that $r_{\text{out}}^{(m)} \le \delta$, this support is
contained in $K' \subseteq \Omega$. Setting $\varphi_n :=
\widetilde\varphi_{n+N}$ for an appropriate $N$ furnishes the desired sequence:
the support condition holds for every $n$, and $\Lp{p}(\Omega)$-convergence
follows from $\Lp{p}(E)$-convergence by monotonicity of the $\Lp{p}$-norm in
the underlying measure (\leanname{eLpNorm\_mono\_measure}).
\end{proof}

\begin{theorem}[Localization: compactly supported $\Wkp{1}{p}$ implies $\Wop{p}$;
\leanname{DeGiorgi.memW01p\_of\_memW1p\_of\_tsupport\_subset}]\label{thm:memW01p_of_compactly_supported}
Let $\Omega \subseteq E$ be open, $1 < p < \infty$, and $u \in W^{1,p}(\Omega)$
with compact support and $\supp u \subseteq \Omega$. Then $u \in \Wop{p}(\Omega)$.
\end{theorem}

\begin{proof}
Choose any $\Wkp{1}{p}$-witness for $u$ on $\Omega$, with gradient field $G$.
Pick $\delta > 0$ such that the closed $\delta$-thickening of $\supp u$ is
contained in $\Omega$, and let $K$ be that thickening. Apply
Lemma~\ref{lem:smooth_cutoff} to obtain a smooth cutoff $\eta$ equal to $1$ on
$K$ with $\supp \eta \subseteq \Omega$. Since $\eta$ is smooth and compactly
supported, both $|\eta|$ and $\|D\eta\|$ are bounded by some constants $C_0 :=
1$ and $C_1 \ge 0$ (the latter using continuity of $D\eta$ on its compact
support, \leanname{HasCompactSupport.exists\_bound\_of\_continuousOn}).

By Proposition~\ref{prop:witness_mul_smooth_bounded} the function $v := \eta u$
admits a $\Wkp{1}{p}$-witness. Since $\eta = 1$ on a neighborhood of $\supp u$
and $u = 0$ outside $\supp u$, we have $v = u$ pointwise. Moreover,
$\supp v \subseteq \supp \eta$ (a compact subset of $\Omega$), and the same is
true for each component of the new gradient, by
\leanname{DeGiorgi.tsupport\_mul\_smooth\_bounded\_p\_weakGrad\_component\_subset}
(private).

We are now in the situation of Theorem~\ref{thm:smooth_approx_local} with $K :=
\supp\eta$. The resulting smooth approximating sequence shows $u = v$ belongs
to $\Wop{p}(\Omega)$.
\end{proof}

\subsection{Whole-space Sobolev inequality}

\begin{definition}[Sobolev constant; \leanname{DeGiorgi.C\_gns}]
For $d \in \N_{\ge 1}$ and $p \in \R$, define
\begin{equation*}
  C_{\mathrm{gns}}(d, p) :=
  \mathrm{SNormLESNormFDerivOfEqConst_{\R}}\bigl(\mathrm{vol}_{\R^d}, p\bigr),
\end{equation*}
the constant supplied by Mathlib's general $W^{1,p} \hookrightarrow L^{p^*}$
embedding. By construction, $C_{\mathrm{gns}}(d,p) \ge 0$
(\leanname{DeGiorgi.C\_gns\_nonneg}).
\end{definition}

\begin{lemma}[Sobolev exponent;
\leanname{DeGiorgi.sobolev\_exp\_pos} (private),
\leanname{DeGiorgi.sobolev\_exponent\_relation} (private)]\label{lem:sobolev_exp}
Let $1 \le p < d$. Then $\frac{dp}{d - p} > 0$, and
\begin{equation*}
  \left(\frac{dp}{d - p}\right)^{-1} = \frac{1}{p} - \frac{1}{d}.
\end{equation*}
\end{lemma}

\begin{proof}
Both numerator $dp$ and denominator $d - p$ are positive (since $d \ge 1$ and
$p < d$). The exponent identity follows by direct algebraic manipulation
$\frac{d - p}{dp} = \frac{1}{p} - \frac{1}{d}$.
\end{proof}

\begin{theorem}[Sobolev inequality for smooth compactly supported functions;
\leanname{DeGiorgi.sobolev\_smooth}]\label{thm:sobolev_smooth}
Let $1 \le p < d$, $u \in C_c^1(E)$. Then
\begin{equation*}
  \|u\|_{L^{p^*}(E)} \le C_{\mathrm{gns}}(d, p)\,\|Du\|_{L^p(E)},
  \qquad p^* = \frac{dp}{d-p}.
\end{equation*}
\end{theorem}

\begin{proof}
This is a direct application of Mathlib's general result
\leanname{MeasureTheory.eLpNorm\_le\_eLpNorm\_fderiv\_of\_eq} with the exponent
relation supplied by Lemma~\ref{lem:sobolev_exp} (rewritten in terms of
$\mathrm{Real.toNNReal}$ to match the Mathlib API).
\end{proof}

\begin{lemma}[Operator norm of the differential equals partials norm;
\leanname{DeGiorgi.norm\_fderiv\_eq\_norm\_partials} (private)]\label{lem:norm_fderiv}
For any $f : E \to \R$ with derivative $Df(x)$ at $x$,
\begin{equation*}
  \|Df(x)\|
   = \left(\sum_{i=1}^d |\partial_i f(x)|^2\right)^{1/2}.
\end{equation*}
\end{lemma}

\begin{proof}
By Riesz representation in the Hilbert space $E$, $Df(x)$ corresponds to a
unique $v \in E$ with $Df(x) = \langle v, \cdot\rangle$, and $\|Df(x)\| = \|v\|$.
Evaluating at the standard basis vector $e_i$ gives $v_i = Df(x)(e_i) =
\partial_i f(x)$, so $\|v\|^2 = \sum_i |\partial_i f(x)|^2$.
\end{proof}

\begin{theorem}[Sobolev inequality from approximating sequences;
\leanname{DeGiorgi.sobolev\_of\_approx}]\label{thm:sobolev_of_approx}
Let $1 \le p < d$, $u : E \to \R$ a.e.\ strongly measurable, and $G : E \to E$
a vector field whose components are a.e.\ strongly measurable. Suppose there
is a sequence $(\varphi_n)$ of smooth compactly supported functions on $E$
with
\begin{equation*}
  \|\varphi_n - u\|_{\Lp{p}(E)} \to 0, \qquad
  \|\partial_i\varphi_n - G_i\|_{\Lp{p}(E)} \to 0 \;\text{ for every } i.
\end{equation*}
Then
\begin{equation*}
  \|u\|_{L^{p^*}(E)} \le C_{\mathrm{gns}}(d,p)\,\|\,\|G(\cdot)\|\,\|_{\Lp{p}(E)},
  \qquad p^* = \frac{dp}{d-p}.
\end{equation*}
\end{theorem}

\begin{proof}
Apply Theorem~\ref{thm:sobolev_smooth} to each $\varphi_n$:
\begin{equation*}
  \|\varphi_n\|_{L^{p^*}} \le C_{\mathrm{gns}}(d,p)\,\|D\varphi_n\|_{L^p}.
\end{equation*}
Define $\Phi_n(x) := (\partial_1\varphi_n(x),\ldots,\partial_d\varphi_n(x))$,
the vector of classical partial derivatives. By Lemma~\ref{lem:norm_fderiv},
$\|D\varphi_n(x)\| = \|\Phi_n(x)\|$, hence $\|D\varphi_n\|_{L^p} = \|\,\|\Phi_n\|\,\|_{L^p}$.

By the triangle inequality applied componentwise (and bounding the Euclidean
norm of a vector by the sum of the norms of its components, see
\leanname{DeGiorgi.gradVec\_eLpNorm\_le\_sum} (private)),
\begin{equation*}
  \|\Phi_n - G\|_{L^p(E; E)}
   \le \sum_{i=1}^d \|\partial_i\varphi_n - G_i\|_{\Lp{p}(E)} \;\to\; 0.
\end{equation*}
Hence $\|\,\|\Phi_n\|\,\|_{L^p} \to \|\,\|G\|\,\|_{L^p}$, by reverse triangle
inequality $|\,\|\Phi_n(x)\| - \|G(x)\|\,| \le \|\Phi_n(x) - G(x)\|$.

Convergence in $L^p$ implies convergence in measure
(\leanname{tendstoInMeasure\_of\_tendsto\_eLpNorm}), so by passing to a
subsequence we may assume $\varphi_{\sigma(n)} \to u$ a.e. By Fatou's lemma
applied to $L^{p^*}$ (\leanname{Lp.eLpNorm\_lim\_le\_liminf\_eLpNorm}),
\begin{equation*}
  \|u\|_{L^{p^*}}
   \le \liminf_{n\to\infty} \|\varphi_{\sigma(n)}\|_{L^{p^*}}
   \le \liminf_{n\to\infty} C_{\mathrm{gns}}(d,p)\,\|\,\|\Phi_{\sigma(n)}\|\,\|_{L^p}.
\end{equation*}
The right-hand side $\liminf$ is bounded above by $C_{\mathrm{gns}}(d,p)\,\|
\,\|G\|\,\|_{L^p}$, since $\|\,\|\Phi_{\sigma(n)}\|\,\|_{L^p} \le \|\,\|G\|\,\|_{L^p}
+ \|\Phi_{\sigma(n)} - G\|_{L^p}$ and the second term tends to zero, using
\leanname{ENNReal.liminf\_add\_of\_right\_tendsto\_zero}. Combining yields the
claim.
\end{proof}

\begin{corollary}[Sobolev inequality on $\Wop{p}(E)$;
\leanname{DeGiorgi.sobolev\_of\_memW01p\_univ}]\label{cor:sobolev_W01p}
Let $1 \le p < d$ and $u \in \Wop{p}(E)$. Then there exists a
$\Wkp{1}{p}$-witness for $u$ on $E$ such that
\begin{equation*}
  \|u\|_{L^{p^*}(E)} \le C_{\mathrm{gns}}(d,p)\,\|\,\|G(\cdot)\|\,\|_{\Lp{p}(E)},
  \qquad p^* = \frac{dp}{d - p},
\end{equation*}
where $G$ is the weak gradient field carried by the witness.
\end{corollary}

\begin{proof}
By Definition~\ref{def:MemW01p}, $u \in \Wop{p}(E)$ comes equipped with a
$\Wkp{1}{p}$-witness on $E$ and an approximating sequence $(\varphi_n)$ of
smooth compactly supported functions converging to $u$ and whose classical
gradients converge to the weak gradient $G$ in $\Lp{p}$. The hypotheses of
Theorem~\ref{thm:sobolev_of_approx} are satisfied, and its conclusion is
exactly the desired inequality.
\end{proof}

\subsection{Positive-part prelude}

\begin{lemma}[Indicator truncation preserves $L^2$;
\leanname{DeGiorgi.indicator\_component\_memLp}]\label{lem:indicator_component_memLp}
Let $\Omega \subseteq E$, $\sigma : E \to \R$ be a.e.\ measurable on $\Omega$
and $g \in \Lp{2}(\Omega)$. Then the function
\begin{equation*}
  x \mapsto \mathbf{1}_{\{\sigma > 0\}}(x)\,g(x)
\end{equation*}
also lies in $\Lp{2}(\Omega)$.
\end{lemma}

\begin{proof}
The set $\{\sigma > 0\}$ is measurable since $\sigma$ is a.e.\ measurable, so
the indicator multiplication preserves a.e.\ strong measurability. Pointwise,
$|\mathbf{1}_{\{\sigma > 0\}}(x)\,g(x)| \le |g(x)|$, and squaring, integrating
and applying the $\Lp{2}$ membership of $g$ gives the conclusion.
\end{proof}

\begin{lemma}[Truncated weak-gradient component in $L^2$;
\leanname{DeGiorgi.posPartGrad\_component\_memLp}]\label{lem:posPartGrad}
Let $u \in H^1(\Omega)$ with $H^1$-witness gradient field $G$. For each
$i \in \{1,\ldots,d\}$, the function
\begin{equation*}
  x \mapsto \mathbf{1}_{\{u > 0\}}(x)\,G_i(x)
\end{equation*}
belongs to $\Lp{2}(\Omega)$.
\end{lemma}

\begin{proof}
Apply Lemma~\ref{lem:indicator_component_memLp} with $\sigma := u$ and $g :=
G_i$, both of which lie in the relevant $\Lp{2}$ spaces (the second by the
witness assumption).
\end{proof}


%% ========================================================================
\section{Functional Inequalities}\label{sec:functional}
%% ========================================================================

% Section 3: Functional Inequalities
% Translation of DeGiorgi/{Poincare,SobolevPoincare,StampacchiaTruncation,Oscillation/*}.lean

This section gathers the functional analytic inequalities used throughout the De Giorgi
machinery: a representation--formula proof of the Poincar\'e inequality on the unit ball,
the Sobolev--Poincar\'e refinement obtained via extension and Gagliardo--Nirenberg--Sobolev,
the Stampacchia truncation lemma, and the BMO / John--Nirenberg / Campanato package
underlying the oscillation theory. Throughout the section, $E = \R^d$ with $d \ge 1$
(in most statements $d \ge 2$, supplied via a \texttt{NeZero} hypothesis), and $\Bz{r}$
denotes the open ball of radius $r$ centred at the origin.

\subsection{Poincar\'e inequality on the unit ball}

\begin{definition}[\leanname{C\_poinc\_val}]
For $d \in \N$, define the Poincar\'e constant
\[
C_{\mathrm{poinc}}(d) := 2^{d+1}\, d.
\]
The lemma \leanname{C\_poinc\_val\_pos} records that $C_{\mathrm{poinc}}(d) > 0$ for
$d > 0$.
\end{definition}

\begin{definition}[\leanname{smoothGradNorm}]
For a smooth scalar function $u : E \to \R$, the pointwise Euclidean gradient norm is
\[
\lvert \nabla u (x)\rvert
:= \Bigl(\sum_{i=1}^d \bigl(\partial_i u(x)\bigr)^{2}\Bigr)^{1/2}.
\]
\end{definition}

The lemma \leanname{norm\_fderiv\_eq\_smoothGradNorm} identifies this expression with the
operator norm of the Fr\'echet derivative $\lVert \mathrm{D}u(x)\rVert$, via the
Riesz representation $\mathrm{D}u(x)\mapsto \nabla u (x)$ in
$\R^d$. Auxiliary support lemmas
\leanname{smoothGradNorm\_eq\_zero\_of\_notMem\_tsupport} and
\leanname{support\_smoothGradNorm\_subset\_tsupport} record that the gradient vanishes
outside the closed support of $u$.

\begin{lemma}[\leanname{smoothCompactSupport\_L2\_bound\_on\_bounded\_ge\_two}]
Let $d \ge 2$ and let $\Omega \subseteq E$ be a bounded set. Then there exists
$C < \infty$ such that for every smooth, compactly supported $u$ with $\supp u
\subseteq \Omega$,
\[
\lVert u\rVert_{L^{2}(\Omega)}
\le C\, \lVert \lvert\nabla u\rvert\rVert_{L^{2}(\Omega)}.
\]
\end{lemma}
\begin{proof}
Let $R$ be such that $\Omega \subseteq \overline{B}(0,R)$ and $s = \overline{B}(0,R)$.
When $d = 2$, the Mathlib Sobolev embedding
\texttt{eLpNormLESNormFDerivOfLeConst} provides the constant
$C_{\mathrm{sob}} = $ \texttt{eLpNormLESNormFDerivOfLeConst}\,$\R\,\mathrm{vol}\,s\,1\,2$
giving
$\lVert u\rVert_{L^{2}(s)} \le C_{\mathrm{sob}}\,\mathrm{vol}(s)^{1/1 - 1/2}
\lVert \mathrm{D}u\rVert_{L^{1}(s)}^{}$ after raising the exponent. For $d \ge 3$ one
applies the same Mathlib bound at the conjugate exponent $2$. In both cases we set
$C := C_{\mathrm{sob}}\,\mathrm{vol}(s)^{1/1 - 1/2}$ and conclude using the inclusion
$\supp u \subseteq \Omega \subseteq s$.
\end{proof}

The proof of the Poincar\'e inequality is via the classical representation formula
through spherical means. We collect the auxiliary lemmas before the main statement.

\begin{lemma}[\leanname{sphere\_radial\_integral\_ball}]
For $R > 0$ and a measurable, integrable, radial--bounded function $f$ on $\Ball{R}{x}$,
\[
\int_{\Ball{R}{x}} f(y)\, dy
= \int_{0}^{R} \int_{\partial \Ball{r}{x}} f \, d\sigma_{r}\, dr.
\]
\end{lemma}

\begin{lemma}[\leanname{polar\_identity\_ball}, \leanname{star\_shaped\_polar\_identity}]
Polar coordinates expansion in $\Ball{R}{x}$: any $y \in \Ball{R}{x}$ can be written
$y = x + r\omega$ with $r \in (0,R)$, $\omega \in S^{d-1}$, and the volume element
factors as $dy = r^{d-1}\, dr\, d\sigma(\omega)$.
\end{lemma}

\begin{lemma}[\leanname{norm\_sub\_le\_integral\_fderiv\_along\_segment}]
For $u \in C^{1}$ and any segment $[x,y] \subseteq E$,
\[
\lvert u(y) - u(x)\rvert
\le \int_{0}^{1} \lvert \nabla u(x + t(y-x))\rvert\, \lvert y-x\rvert\, dt.
\]
\end{lemma}

\begin{lemma}[\leanname{integral\_norm\_rpow\_neg\_ball},
\leanname{riesz\_kernel\_integrable}, \leanname{riesz\_kernel\_L1\_bound}]
For $R > 0$ and $0 < \alpha < d$ the function $z\mapsto \lvert z\rvert^{\alpha-d}$ is
integrable on $\Ball{R}{0}$, and for $x \in \Ball{R}{0}$ the Riesz kernel satisfies
\[
\int_{\Ball{R}{0}} \lvert x - y\rvert^{1-d}\, dy
\le d\, \mathrm{vol}(\Bz{1})\, (2R).
\]
\end{lemma}

\begin{lemma}[Representation formula, \leanname{representation\_formula\_smooth}]
Let $u \in C^{\infty}$ and let $\bar u = \fint_{\Bz{1}} u$. For every $x \in \Bz{1}$,
\[
\lvert u(x) - \bar u\rvert
\le \frac{2^{d}}{d\,\mathrm{vol}(\Bz{1})}\,
\int_{\Bz{1}} \lvert \nabla u(y)\rvert\, \lvert x - y\rvert^{1-d}\, dy.
\]
\end{lemma}
\begin{proof}
For $\omega \in S^{d-1}$ and $r$ such that $x + r\omega \in \Bz{1}$ one writes
$u(x+r\omega) - u(x) = \int_{0}^{r} \partial_{\omega} u(x + s\omega)\, ds$ and integrates
in $r$ along radial segments inside $\Bz{1}$. Averaging this identity over $\omega$
weighted by the chord length, and using
\leanname{star\_shaped\_polar\_identity} to convert the spherical integral into a
weighted volume integral with weight $\lvert x - y\rvert^{1-d}$, gives the stated
pointwise estimate. The constant $2^{d}/(d\,\mathrm{vol}(\Bz{1}))$ comes from estimating
the maximal chord length of $\Bz{1}$ from any point $x \in \Bz{1}$ by $2$.
\end{proof}

\begin{lemma}[Weighted power mean for set integrals,
\leanname{weighted\_power\_mean\_setIntegral}]
Let $p \ge 1$, let $w \ge 0$ be a weight on a measurable set $S$ with $W := \int_S w\, dy$
finite, and let $f \ge 0$. Then
\[
\Bigl(\int_S f\,w\, dy\Bigr)^{p}
\le W^{p-1}\, \int_S f^{p}\,w\, dy.
\]
\end{lemma}
\begin{proof}
Normalise $w$ to a probability measure $d\nu = w/W\, dy$ and apply Jensen's inequality
to the convex function $t \mapsto t^{p}$, then rescale.
\end{proof}

\begin{theorem}[Poincar\'e on the unit ball for smooth functions,
\leanname{poincare\_smooth\_unitBall}]
Let $p \ge 1$ and $u \in C^{\infty}(E,\R)$. Then
\[
\Bigl\lVert u - \fint_{\Bz{1}} u \Bigr\rVert_{\Lp{p}(\Bz{1})}
\le C_{\mathrm{poinc}}(d)\, \bigl\lVert \lvert\nabla u\rvert \bigr\rVert_{\Lp{p}(\Bz{1})}.
\]
\end{theorem}
\begin{proof}
Set $B = \Bz{1}$, $\bar u = \fint_B u$, $f = u - \bar u$, $g = \lvert \nabla u\rvert$,
$K(z) = \lvert z\rvert^{1-d}$ and
\[
h(x) = \int_B g(y)\, K(x-y)\, dy,\qquad
M = d\,\mathrm{vol}(B)\cdot 2,\qquad
C_{\mathrm{rep}} = \frac{2^{d}}{d\,\mathrm{vol}(B)}.
\]
The representation formula yields $\lvert f(x)\rvert \le C_{\mathrm{rep}}\, h(x)$ pointwise on
$B$. Therefore $\lVert f\rVert_{\Lp{p}(B)} \le C_{\mathrm{rep}}\,\lVert h\rVert_{\Lp{p}(B)}$.

For the Young-type estimate $\lVert h\rVert_{\Lp{p}(B)} \le M\,\lVert g\rVert_{\Lp{p}(B)}$
we apply the weighted power mean lemma to the kernel weight $w(y) = K(x-y)$, using the
$L^{1}$ bound $\int_B K(x-y)\, dy \le M$ provided by \leanname{riesz\_kernel\_L1\_bound}.
This gives
\[
h(x)^{p}
\le M^{p-1}\, \int_B g(y)^{p}\, K(x-y)\, dy.
\]
Integrating in $x \in B$, using Fubini and the same kernel bound in the $x$ variable,
\[
\int_B h(x)^{p}\, dx
\le M^{p-1}\, \int_B g(y)^{p}\Bigl(\int_B K(x-y)\, dx\Bigr) dy
\le M^{p}\, \int_B g(y)^{p}\, dy.
\]
Taking $p$-th roots gives $\lVert h\rVert_{\Lp{p}(B)} \le M\,\lVert g\rVert_{\Lp{p}(B)}$.
Combining,
\[
\lVert f\rVert_{\Lp{p}(B)}
\le C_{\mathrm{rep}}\, M\, \lVert g\rVert_{\Lp{p}(B)}
= 2^{d+1}\, d\, \lVert g\rVert_{\Lp{p}(B)}
= C_{\mathrm{poinc}}(d)\,\lVert g\rVert_{\Lp{p}(B)}. \qedhere
\]
\end{proof}

A density argument lifts this to general $\Wkp{1}{p}(\Bz{1})$ functions.

\begin{theorem}[Poincar\'e on the unit ball for $W^{1,p}$,
\leanname{poincare\_unitBall\_W1p\_public}]
Let $1 < p < \infty$ and let $u \in \Wkp{1}{p}(\Bz{1})$ with weak gradient $G$. Then
\[
\Bigl\lVert u - \fint_{\Bz{1}} u \Bigr\rVert_{\Lp{p}(\Bz{1})}
\le C_{\mathrm{poinc}}(d)\, \lVert \lvert G\rvert\rVert_{\Lp{p}(\Bz{1})}.
\]
\end{theorem}
\begin{proof}
Approximate $u$ by smooth $\psi_n \to u$ in $\Wkp{1}{p}(\Bz{1})$. Apply the smooth
Poincar\'e inequality to each $\psi_n$, then pass to the limit using strong $L^{p}$
convergence of $\psi_n - \fint \psi_n$ to $u - \fint u$ and of
$\nabla \psi_n$ to $G$.
\end{proof}

\subsection{Sobolev--Poincar\'e on the unit ball}

\begin{definition}[\leanname{C\_extensionGrad}, \leanname{C\_sobolevPoincare}]
For $d \ge 1$ and $p > 1$,
\begin{align*}
C_{\mathrm{ext\!-\!grad}}(d,p)
  &:= 1 + C_{\mathrm{ubExt\,grad}}(d,p)\bigl(C_{\mathrm{poinc}}(d) + 1\bigr),\\
C_{\mathrm{sobPoinc}}(d,p)
  &:= C_{\mathrm{gns}}(d,p)\, C_{\mathrm{ext\!-\!grad}}(d,p),
\end{align*}
where $C_{\mathrm{gns}}$ is the whole-space Gagliardo--Nirenberg--Sobolev constant and
$C_{\mathrm{ubExt\,grad}}$ comes from the unit-ball extension.
\end{definition}

\begin{lemma}[Subtracting a constant preserves $W^{1,p}$,
\leanname{memW1pWitness\_sub\_const\_unitBall},
\leanname{memW1pWitness\_sub\_const\_unitBall\_weakGrad}]
If $u \in \Wkp{1}{p}(\Bz{1})$ has weak gradient $G$, then for any $c \in \R$ the function
$u - c$ is also in $\Wkp{1}{p}(\Bz{1})$ with the same weak gradient $G$.
\end{lemma}
\begin{proof}
$u - c \in L^{p}$ since $L^{p}$ is a vector space. Constants have zero weak gradient by
$\int_{\Bz{1}} c\, \partial_i \varphi = 0$ for compactly supported $\varphi$, hence
subtracting a constant from the test integral relation
$\int u\,\partial_i \varphi = -\int G_i \varphi$ leaves it unchanged.
\end{proof}

\begin{lemma}[Jensen for averages, \leanname{eLpNorm\_const\_average\_le}]
If $\mu$ is a finite measure, $1 \le p < \infty$ and $f \in L^{1}(\mu)$, then
\[
\Bigl\lVert x \mapsto \fint f\, d\mu \Bigr\rVert_{\Lp{p}(\mu)}
\le \lVert f \rVert_{\Lp{p}(\mu)}.
\]
\end{lemma}
\begin{proof}
Apply H\"older's inequality to bound the average $\lvert\fint f\rvert$ by the
$L^{p}$ norm of $f$ times $\mu(X)^{1-1/p}$, then integrate.
\end{proof}

\begin{theorem}[Sobolev--Poincar\'e on $\Bz{1}$, \leanname{sobolev\_poincare\_unitBall}]
Let $1 < p < d$ and $u \in \Wkp{1}{p}(\Bz{1})$ with weak gradient $G$, and write
$p^{*} = dp/(d-p)$. Then
\[
\Bigl\lVert u - \fint_{\Bz{1}} u \Bigr\rVert_{\Lp{p^{*}}(\Bz{1})}
\le C_{\mathrm{sobPoinc}}(d,p)\, \lVert\lvert G\rvert\rVert_{\Lp{p}(\Bz{1})}.
\]
\end{theorem}
\begin{proof}
Set $\bar u = \fint_{\Bz{1}} u$ and $v = u - \bar u$. By the previous lemma
$v \in \Wkp{1}{p}(\Bz{1})$ with weak gradient $G$. Let
$\widetilde v = \mathrm{ext}(v) \in \Wop{p}$ be the unit-ball extension, which has
compact support, agrees with $v$ on $\Bz{1}$ and satisfies the lintegral gradient bound
\[
\int \lvert \nabla \widetilde v\rvert^{p}
\le \int_{\Bz{1}} \lvert G\rvert^{p}
+ C_{\mathrm{ubExt\,grad}}(d,p)\,
\Bigl(\int_{\Bz{1}} \lvert v\rvert^{p} + \int_{\Bz{1}} \lvert G\rvert^{p}\Bigr).
\]
Combining this with Poincar\'e
$\lVert v\rVert_{\Lp{p}(\Bz{1})} \le C_{\mathrm{poinc}}(d)\,\lVert\lvert G\rvert\rVert_{\Lp{p}(\Bz{1})}$
and raising to the $p$-th power gives
\[
\lVert \lvert \nabla \widetilde v\rvert\rVert_{\Lp{p}(\R^{d})}
\le C_{\mathrm{ext\!-\!grad}}(d,p)\,\lVert\lvert G\rvert\rVert_{\Lp{p}(\Bz{1})}
\]
(this is the auxiliary lemma \leanname{extension\_gradient\_eLpNorm\_bound}, which uses
the algebraic inequalities $a^{p}+b^{p} \le (a+b)^{p}$ for $p\ge 1$ to absorb the
constants). The whole-space Gagliardo--Nirenberg--Sobolev inequality applied to
$\widetilde v$ gives
\[
\lVert \widetilde v\rVert_{\Lp{p^{*}}(\R^{d})}
\le C_{\mathrm{gns}}(d,p)\,
\lVert \lvert\nabla \widetilde v\rvert\rVert_{\Lp{p}(\R^{d})}.
\]
Restricting the left side to $\Bz{1}$, where $\widetilde v = v$, and chaining the two
displayed estimates with the Poincar\'e bound yields the conclusion with constant
$C_{\mathrm{sobPoinc}}(d,p) = C_{\mathrm{gns}}(d,p)\,C_{\mathrm{ext\!-\!grad}}(d,p)$.
\end{proof}

\begin{corollary}[Existential form, \leanname{sobolev\_poincare\_unitBall'}]
For $1 < p < d$ and $u$ in the predicate $\Wkp{1}{p}(\Bz{1})$ (existential form), there
exist a weak gradient $G$ and a constant $C = C_{\mathrm{sobPoinc}}(d,p)$ realising the
above Sobolev--Poincar\'e estimate.
\end{corollary}

\begin{corollary}[Smooth version, \leanname{sobolev\_poincare\_smooth\_unitBall}]
For $1 < p < d$ and $u \in C^{\infty}(E,\R)$,
\[
\Bigl\lVert u - \fint_{\Bz{1}} u \Bigr\rVert_{\Lp{p^{*}}(\Bz{1})}
\le C_{\mathrm{sobPoinc}}(d,p)\,\lVert\lvert\nabla u\rvert\rVert_{\Lp{p}(\Bz{1})}.
\]
\end{corollary}
\begin{proof}
Build the canonical $W^{1,p}$ witness for the smooth $u$ via
\leanname{smooth\_memW1pWitness\_unitBall} and apply the previous theorem; the weak
gradient agrees a.e.\ with the classical gradient on $\Bz{1}$.
\end{proof}

\subsection{Stampacchia truncation}

\begin{lemma}[\leanname{deriv\_ne\_zero\_implies\_isolated\_zero}]
Let $f : \R \to \R$ be differentiable at $x$ with $f(x)=0$ and $f'(x) = c \neq 0$. Then
there exists $\delta > 0$ such that $f(x+h) \neq 0$ for all $h$ with $0 < \lvert h\rvert < \delta$.
\end{lemma}
\begin{proof}
The little-$o$ definition of differentiability gives, for sufficiently small $h$,
$\lvert f(x+h) - c h\rvert \le (\lvert c\rvert/2)\, \lvert h\rvert$. With $f(x)=0$
the reverse triangle inequality yields $\lvert f(x+h)\rvert \ge (\lvert c\rvert/2)\,
\lvert h\rvert > 0$.
\end{proof}

\begin{corollary}[\leanname{HasDerivAt.deriv\_eq\_zero\_of\_zero\_and\_not\_isolated}]
If $f$ is differentiable at $x$ with $f(x) = 0$ and $x$ is a non-isolated zero, then
$f'(x) = 0$.
\end{corollary}

\begin{lemma}[Isolated points are countable, \leanname{countable\_isolated}]
For any $S \subseteq \R$, the set of isolated points of $S$ is countable.
\end{lemma}
\begin{proof}
Each isolated point $x$ admits a rational interval $(p,q) \ni x$ inside which $x$ is the
only point of $S$. Map $x \mapsto (p,q) \in \Q^{2}$; the map is well-defined and
injective on isolated points.
\end{proof}

\begin{lemma}[1D Stampacchia, a.e.\ form, \leanname{deriv\_eq\_zero\_ae\_on\_zeroSet}]
If $f : \R \to \R$ is differentiable a.e., then $f' = 0$ a.e.\ on $\{f = 0\}$.
\end{lemma}
\begin{proof}
The bad set $\{f = 0,\ f' \neq 0\}$ is contained in the union of the points where $f$
is not differentiable (a null set by hypothesis) and the set of isolated zeros of
$\{f = 0\}$, which is countable by the previous lemma. Both have measure zero.
\end{proof}

\begin{lemma}[Du Bois-Reymond fundamental lemma, \leanname{du\_bois\_reymond}]
Let $g \in L^{1}(a,b)$ satisfy $\int_a^b g\,\varphi' = 0$ for every test function
$\varphi \in C_c^{\infty}(a,b)$. Then $g$ is constant a.e.
\end{lemma}
\begin{proof}
Fix $\psi_0 \in C_c^{\infty}(a,b)$ with $\int \psi_0 = 1$ and let
$c = \int g\,\psi_0$. For any test $\eta$ write $\eta = \eta_0 + \alpha \psi_0$ where
$\alpha = \int\eta$ and $\int \eta_0 = 0$. Since $\eta_0$ has zero integral, it is the
derivative of a test function $\varphi$ (antiderivative lemma), so the hypothesis gives
$\int g \eta_0 = 0$, hence $\int g\,\eta = c\,\int\eta$. Equivalently
$\int (g - c)\eta = 0$ for all test $\eta$, so by the fundamental lemma of the calculus
of variations $g = c$ a.e.
\end{proof}

\begin{lemma}[$W^{1,1}$ AC representative, \leanname{w11\_ae\_eq\_ac\_representative}]
Let $u, g \in L^{1}(a,b)$ with the weak-derivative relation
$\int u\,\varphi' = -\int g\,\varphi$ for all $\varphi \in C_c^{\infty}(a,b)$. Then there
is a constant $C$ such that
\[
u(x) = C + \int_a^{x} g(t)\, dt \qquad \text{for a.e. } x \in (a,b).
\]
\end{lemma}
\begin{proof}
Let $F(x) = \int_a^{x} g$. By integration by parts for absolutely continuous functions,
$F$ has weak derivative $g$. Hence $u - F$ has zero weak derivative, and du Bois-Reymond
gives $u - F = C$ a.e.
\end{proof}

\begin{theorem}[1D Stampacchia for $W^{1,1}$, \leanname{stampacchia\_1d}]
Under the hypotheses of the previous lemma, $g(x) = 0$ for a.e.\ $x \in (a,b)$
with $u(x) = 0$.
\end{theorem}
\begin{proof}
Replace $u$ by its absolutely continuous representative $F(x) = C + \int_a^x g$. Then
$F' = g$ a.e.\ by the Lebesgue differentiation theorem for locally integrable
functions, and \leanname{deriv\_eq\_zero\_ae\_on\_zeroSet} applied to $F$ yields
$F' = 0$ a.e.\ on $\{F = 0\}$, i.e.\ $g = 0$ a.e.\ on $\{u = 0\}$.
\end{proof}

\begin{theorem}[$d$-dimensional Stampacchia, \leanname{weakGrad\_ae\_eq\_zero\_on\_zeroSet}]
Let $\Omega \subseteq E$ be open and $u \in \Wkp{1}{2}(\Omega)$ with weak gradient
$G : \Omega \to \R^{d}$. Then for each $i \in \{1,\dots,d\}$, $G_i(x) = 0$ for a.e.\
$x \in \Omega$ with $u(x) = 0$.
\end{theorem}
\begin{proof}
Set $h_i(x) = \mathbf{1}_{\{u = 0\}}(x)\, G_i(x)$. We show $h_i = 0$ a.e.\ on $\Omega$
by reducing to the variational fundamental lemma: for any test function
$\psi \in C_c^{\infty}(\Omega)$ the integral $\int \psi\, h_i$ vanishes. By Fubini's
theorem, the line restrictions of $u$ along the $i$-th coordinate axis lie in
$W^{1,1}$ for a.e.\ line, with weak derivative the line restriction of $G_i$. The 1D
Stampacchia theorem applied on each line gives $G_i = 0$ a.e.\ on
$\{u = 0\}$ along almost every line. Recombining via Fubini gives the result.
\end{proof}

\subsection{BMO and the John--Nirenberg geometry}

\begin{definition}[BMO seminorm, \leanname{bmoNorm}]
For $u : E \to \R$ and $U \subseteq E$,
\[
\lVert u\rVert_{\mathrm{BMO}(U)}
:= \sup\Bigl\{\, \fint_{\Ball{r}{x_0}} \bigl|\, u - \tfrac{1}{|B|}\!\int_{\Ball{r}{x_0}} u\,\bigr|
\;:\; \overline{\Ball{r}{x_0}} \subseteq U,\ r > 0\,\Bigr\}.
\]
\end{definition}

\begin{definition}[John--Nirenberg constants, \leanname{C\_JN}, \leanname{C\_iter}]
\[
C_{\mathrm{JN}}(d) := 16 \cdot 36^{d}, \qquad C_{\mathrm{iter}} := 1024.
\]
Both are strictly positive (\leanname{C\_JN\_pos}, \leanname{C\_iter\_pos}).
\end{definition}

\begin{definition}[\leanname{JNBall}]
For $x_0 \in E$, $R > 0$, a \emph{JN-ball} is a record consisting of a centre
$c \in \Ball{R}{x_0}$ and a radius $0 < \rho \le R$. Its \emph{carrier} is
$\Ball{\rho}{c}$ and its \emph{five-fold enlargement} is $\Ball{5\rho}{c}$.
\end{definition}

\begin{lemma}[Vitali $5r$-covering, \leanname{vitali\_covering\_lemma}]
Given a family of balls $(\Ball{r_i}{x_i})_{i \in \iota}$ with $r_i > 0$ uniformly bounded
above, there is a countable subfamily $S \subseteq \iota$ such that the carriers
$(\Ball{r_i}{x_i})_{i \in S}$ are pairwise disjoint, every original ball meets some
selected ball whose radius is at least half that of the original, and
$\bigcup_{i \in \iota} \Ball{r_i}{x_i} \subseteq \bigcup_{i \in S} \Ball{5 r_i}{x_i}$.
\end{lemma}

\begin{lemma}[Geometric absorption, \leanname{ball\_subset\_fivefold\_of\_inter\_nonempty}]
If $\Ball{\rho}{a} \subseteq \Ball{r}{c}$, $\Ball{r}{c} \cap \Ball{s}{d} \neq \emptyset$
and $r \le 2 s$, then $\Ball{\rho}{a} \subseteq \Ball{5 s}{d}$.
\end{lemma}

\begin{lemma}[Iteration of one-step decay,
\leanname{john\_nirenberg\_iteration\_decay\_iterate},
\leanname{john\_nirenberg\_iteration}]
Let $E_{\ell} \subseteq B$ be an antitone family of measurable subsets of a measurable
set $B$ with $\mathrm{vol}(B) < \infty$. Suppose there exist $A > 0$ and
$0 < \theta < 1$ with
\[
\mathrm{vol}(E_{\ell + A}) \le \theta\, \mathrm{vol}(E_{\ell})
\quad\text{for all } \ell > 0.
\]
Then for every $\ell > 0$,
\[
\mathrm{vol}(E_{\ell})
\le \frac{1}{\theta}\, \mathrm{vol}(B)\,
\exp\!\Bigl(-\ell\, \frac{-\log \theta}{A}\Bigr).
\]
\end{lemma}
\begin{proof}
Iterating one-step decay $n$ times gives
$\mathrm{vol}(E_{\ell_0 + n A}) \le \theta^{n}\,\mathrm{vol}(E_{\ell_0})$. Choose
$n = \lceil \ell/A\rceil - 1$, so $\ell_0 = \ell - n A \in (0, A]$, and bound
$\theta^{n} \le (1/\theta)\, \exp(-\ell\,(-\log\theta)/A)$ using
$n \ge \ell/A - 1$ and $\theta^{x} = \exp(x\log\theta)$.
\end{proof}

\begin{theorem}[John--Nirenberg from one-step decay, \leanname{john\_nirenberg\_iteration}]
Let $u : E \to \R$ be measurable and let $x_0, r$ with $r > 0$. If for some $A,\theta$
with $A > 0$, $0 < \theta < 1$ the sublevel sets of $|u - \fint_{\Ball{r}{x_0}} u|$
satisfy
\[
\mathrm{vol}\bigl\{ x \in \Ball{r}{x_0} : \lvert u - \tfrac{1}{|B|}\!\int_{\Ball{r}{x_0}} u\rvert > \ell + A\bigr\}
\le \theta\, \mathrm{vol}\bigl\{ x \in \Ball{r}{x_0} : \lvert u - \tfrac{1}{|B|}\!\int_{\Ball{r}{x_0}} u\rvert > \ell\bigr\}
\]
for all $\ell > 0$, then for every $t > 0$,
\[
\mathrm{vol}\bigl\{ x \in \Ball{r}{x_0} : \lvert u - \tfrac{1}{|B|}\!\int_{\Ball{r}{x_0}} u\rvert > t\bigr\}
\le \frac{1}{\theta}\, \mathrm{vol}(\Ball{r}{x_0})\,
\exp\!\Bigl(-t\,\frac{-\log\theta}{A}\Bigr).
\]
\end{theorem}

\begin{lemma}[Indicator BMO, \leanname{indicator\_bmo\_bound}]
If $u$ has BMO seminorm at most $M$ on every closed sub-ball of $\Ball{r}{x_0}$, then
the same bound holds for the function $\mathbf{1}_{\Ball{r}{x_0}}\, u$ on those sub-balls,
since the average over $\Ball{s}{z} \subseteq \Ball{r}{x_0}$ coincides with that of $u$.
\end{lemma}

\subsection{Simple iteration lemma}

\begin{lemma}[\leanname{simple\_iteration\_lemma}]
Let $A_{\mathrm{iter}} \ge 0$ and let $\rho : \R \to \R$ satisfy on $[1/2, 1)$:
\begin{enumerate}
\item $\rho \ge 0$;
\item $\sup_{1/2 \le t < 1}(1-t)^{2}\,\rho(t) < \infty$;
\item for all $1/2 \le s < t < 1$,
$\rho(s) \le \tfrac{1}{2}\rho(t) + A_{\mathrm{iter}}(t-s)^{-2}$.
\end{enumerate}
Then $\rho(1/2) \le C_{\mathrm{iter}}\, A_{\mathrm{iter}}$.
\end{lemma}
\begin{proof}
Let $M$ be a non-negative upper bound for $(1-t)^{2}\rho(t)$. With the choice
$t = (3s+1)/4$, one has $t-s = (1-s)/4$ and $1-t = 3(1-s)/4$, so multiplying the
contraction by $(1-s)^{2}$ and using $(1-s)^{2}\,(t-s)^{-2} = 16$ gives
\[
(1-s)^{2}\,\rho(s) \le \tfrac{8}{9}\,(1-t)^{2}\rho(t) + 16\, A_{\mathrm{iter}}
\le \tfrac{8}{9}\, M + 16\, A_{\mathrm{iter}}.
\]
Taking the supremum over $s$ shows $M \mapsto \tfrac{8}{9}M + 16 A_{\mathrm{iter}}$ also
bounds $(1-t)^{2}\rho(t)$. Iterating $n$ times gives the bound
$M_n = (8/9)^{n} M + (1-(8/9)^{n})\cdot 144\, A_{\mathrm{iter}}$, which tends to
$144\,A_{\mathrm{iter}}$. Evaluating at $s = 1/2$ where $(1-s)^{2} = 1/4$, we get
$\rho(1/2) \le 4 M_n \to 576\, A_{\mathrm{iter}} \le C_{\mathrm{iter}}\,A_{\mathrm{iter}}$.
\end{proof}

\subsection{Local John--Nirenberg theorem}

\begin{lemma}[Geometric decay step, \leanname{john\_nirenberg\_level\_set\_decay}]
Let $E_{\ell+A} \subseteq E_{\ell}$ be measurable sets with $\mathrm{vol}(E_{\ell}) < \infty$
and let $(B_i)$ be a countable family of pairwise disjoint JN-balls with carriers
$\bigcup_i B_i^{\mathrm{car}} \subseteq E_{\ell}$ and 5-fold enlargements covering
$E_{\ell}$. If
\[
\mathrm{vol}(E_{\ell+A} \cap B_i^{\mathrm{car}}) \le \tfrac{1}{2}\, \mathrm{vol}(B_i^{\mathrm{car}})
\quad \text{for every } i,
\]
then
\[
\mathrm{vol}(E_{\ell+A}) \le \Bigl(1 - \frac{1}{2 \cdot 5^{d}}\Bigr)\, \mathrm{vol}(E_{\ell}).
\]
\end{lemma}
\begin{proof}
Let $U = \bigcup_i B_i^{\mathrm{car}}$. By disjointness and the 5-fold cover,
$\mathrm{vol}(U) \ge 5^{-d}\, \mathrm{vol}(E_{\ell})$. The hypothesis on each carrier
sums to $\mathrm{vol}(E_{\ell+A}\cap U) \le \tfrac{1}{2}\,\mathrm{vol}(U)$. Splitting
$E_{\ell+A} = (E_{\ell+A}\cap U) \cup (E_{\ell+A}\setminus U)$ and using
$E_{\ell+A}\setminus U \subseteq E_{\ell}\setminus U$ and the lower bound on
$\mathrm{vol}(U)$ gives the stated estimate.
\end{proof}

\begin{theorem}[Iterated decay from a base level, \leanname{john\_nirenberg\_from\_base},
\leanname{level\_set\_family\_from\_base}]
If a measurable antitone family $E_{\ell} \subseteq B$ satisfies the one-step decay
$\mathrm{vol}(E_{\ell+A}) \le \theta\,\mathrm{vol}(E_{\ell})$ only for $\ell \ge A$, then
for every $t > 0$,
\[
\mathrm{vol}(E_{t}) \le \frac{1}{\theta^{2}}\, \mathrm{vol}(B)\,
\exp\!\Bigl(-t\,\frac{-\log\theta}{A}\Bigr).
\]
\end{theorem}
\begin{proof}
For $t > A$ apply \leanname{john\_nirenberg\_iteration} to the shifted family
$F_s = E_{s+A}$, picking up an extra factor $1/\theta$ from the shift; the bound
$\theta^{n} \le (1/\theta) \exp(\dots)$ contributes another $1/\theta$, totalling
$1/\theta^{2}$. For $t \le A$ the trivial inclusion
$E_{t} \subseteq B$ already lies below the right-hand side because the exponential factor
is bounded below by $\theta$.
\end{proof}

\begin{theorem}[Local John--Nirenberg, \leanname{john\_nirenberg\_local}]
Let $R > 0$, $M > 0$, and $u$ measurable on $\Ball{6R}{x_0}$ with the BMO bound
\[
\fint_{\Ball{s}{z}} \Bigl|\, u - \fint_{\Ball{s}{z}} u\, \Bigr|\, dy \le M
\]
for every closed sub-ball $\overline{\Ball{s}{z}} \subseteq \Ball{6R}{x_0}$. Then for every
$t > 0$,
\[
\mathrm{vol}\Bigl\{ x \in \Ball{R}{x_0} : \Bigl|\,u(x) - \fint_{\Ball{R}{x_0}} u\,\Bigr| > t\Bigr\}
\le C_{\mathrm{JN}}(d)\, \mathrm{vol}(\Ball{R}{x_0})\,
\exp\!\Bigl(-\frac{t}{C_{\mathrm{JN}}(d)\,M}\Bigr).
\]
\end{theorem}
\begin{proof}
Set $w(x) = |u(x) - \fint_{\Ball{R}{x_0}} u|$, $\theta = 1/2$, $A = 8\cdot 25^{d}\, M$.
For each $\ell \ge A$, apply a Calder\'on--Zygmund stopping-time / Vitali selection
inside $\Ball{R}{x_0}$ on the level set $\{w > \ell\}$ to extract a disjoint family of
JN-balls $(B_i)$ whose 5-fold enlargements cover the level set (up to a null set), with
the property that on each parent ball the average of $w$ does not exceed $\ell$. The
BMO hypothesis applied on each $B_i$ then yields
\[
\mathrm{vol}(\{w > \ell + A\}\cap B_i^{\mathrm{car}})
\le \tfrac{1}{2}\, \mathrm{vol}(B_i^{\mathrm{car}}),
\]
because raising $w$ by $A$ over the parent average exceeds the BMO control of the
oscillation. Lemma \leanname{john\_nirenberg\_level\_set\_decay} converts this into the
geometric decay $\mathrm{vol}(\{w > \ell + A\}) \le \theta\, \mathrm{vol}(\{w > \ell\})$ for
all $\ell \ge A$. Theorem \leanname{john\_nirenberg\_from\_base} then yields the
exponential bound, where $-\log\theta/A = \log 2 / (8\cdot 25^{d}\, M)$ and the
constants combine into $C_{\mathrm{JN}}(d) = 16\cdot 36^{d}$.
\end{proof}

\subsection{Campanato characterisation of H\"older spaces}

\begin{definition}[\leanname{CampanatoBall}, \leanname{campanatoBallValue},
\leanname{campanatoSeminorm}]
For $x_0 \in E$ and $R > 0$, a \emph{Campanato ball} is a pair $(c,\rho)$ with
$0 < \rho \le R$ and $\Ball{\rho}{c} \subseteq \Ball{R}{x_0}$. For $\alpha > 0$ define
\[
N_{\alpha}(u; c,\rho)
:= \rho^{-\alpha}\, \fint_{\Ball{\rho}{c}}
\Bigl|\, u(x) - \fint_{\Ball{\rho}{c}} u\, \Bigr|\, dx,
\]
and the Campanato seminorm
\[
[u]_{\mathcal{L}^{\alpha}(\Ball{R}{x_0})}
:= \sup_{(c,\rho)\,\text{Campanato ball}} N_{\alpha}(u;c,\rho).
\]
\end{definition}

\begin{definition}[\leanname{HasCampanatoBound}]
We write $\mathrm{HasCampanatoBound}(u; x_0, R, \alpha, C)$ if $u$ is locally integrable on every
admissible sub-ball and $N_{\alpha}(u; c, \rho) \le C$ for every Campanato ball
$(c,\rho) \subseteq \Ball{R}{x_0}$.
\end{definition}

\begin{definition}[Dyadic ball averages and limits,
\leanname{dyadicBallAverage}, \leanname{dyadicBallAverageLimit}]
For $\rho > 0$ set
\begin{gather*}
a_n(x) := \fint_{\Ball{\rho/2^{n}}{x}} u\, dz, \\
\widetilde u(x) := \lim_{n\to\infty} a_n(x)
\end{gather*}
when the limit exists; $\widetilde u$ is the candidate H\"older representative.
\end{definition}

\begin{definition}[\leanname{C\_campanato\_holder}]
\[
C_{\mathrm{Camp\to H\ddot ol}}(d,\alpha) := \frac{2^{d+3}}{1 - 2^{-\alpha}}.
\]
\end{definition}

The chain of lemmas
\leanname{abs\_ballAverage\_half\_sub\_ballAverage\_le},
\leanname{ballAverage\_sub\_ballAverage\_sameRadius\_le\_campanato},
\leanname{dyadicBallAverage\_step\_le\_campanato},
\leanname{dyadicBallAverage\_step\_le\_geometric},
\leanname{dyadicBallAverage\_cauchy},
\leanname{tendsto\_dyadicBallAverageLimit},
\leanname{dyadicBallAverage\_tail\_le}
establishes that the dyadic averages form a
Cauchy sequence with geometric ratio $2^{-\alpha}$, and converge to a limit
$\widetilde u$ with quantitative tail control. Lemmas
\leanname{dyadicBallAverageLimit\_holder\_small} and
\leanname{dyadicBallAverageLimit\_holder\_large}
give pointwise H\"older estimates on the
limit in the regimes $|x-y| \le \rho$ and $|x-y| > \rho$ respectively.

\begin{theorem}[Campanato implies H\"older, \leanname{campanato\_implies\_holder}]
Let $0 < \alpha \le 1$, $R > 0$ and $\mathrm{HasCampanatoBound}(u; x_0, R, \alpha, C_{\mathrm{camp}})$.
Then there is a function $v : E \to \R$ that agrees with $u$ a.e.\ on $\Ball{R}{x_0}$ and
satisfies, for all $x, y \in \Ball{R/2}{x_0}$,
\[
\bigl|v(x) - v(y)\bigr|
\le C_{\mathrm{Camp\to H\ddot ol}}(d,\alpha)\, C_{\mathrm{camp}}\, \lvert x - y\rvert^{\alpha}.
\]
\end{theorem}
\begin{proof}
Take $\rho = R/2$ and define $v(x) := \widetilde u(x)$ for $x \in \Ball{R/2}{x_0}$ and
$v(x) := u(x)$ otherwise. By
\leanname{ae\_eq\_dyadicBallAverageLimit\_on\_halfBall}, $\widetilde u = u$ a.e.\ on
$\Ball{R/2}{x_0}$, so $v = u$ a.e.\ on $\Ball{R}{x_0}$. For $x, y \in \Ball{R/2}{x_0}$,
split into the cases $|x - y| \le \rho$ and $|x - y| > \rho$, applying the corresponding
H\"older estimate for the dyadic limit. Both cases yield the same bound with constant
$C_{\mathrm{Camp\to H\ddot ol}}(d,\alpha)\, C_{\mathrm{camp}}$.
\end{proof}

\begin{lemma}[H\"older implies Campanato bound, \leanname{holder\_hasCampanatoBound}]
If $u$ is $\alpha$-H\"older on $\Ball{R}{x_0}$ with constant $C_{\mathrm{H\ddot ol}}$, then
$\mathrm{HasCampanatoBound}(u; x_0, R, \alpha, C_{\mathrm{H\ddot ol}}\cdot 2^{\alpha})$.
\end{lemma}
\begin{proof}
On any Campanato ball $\Ball{\rho}{y}$, for any $x, z \in \Ball{\rho}{y}$ we have
$|x - z| \le 2\rho$, so the H\"older bound gives
\[
|u(x) - u(z)| \le C_{\mathrm{H\ddot ol}}\,\lvert x - z\rvert^{\alpha}
\le C_{\mathrm{H\ddot ol}}\,(2\rho)^{\alpha}.
\]
Averaging over $z \in \Ball{\rho}{y}$ via $\overline u = \fint_{\Ball{\rho}{y}} u(z)\,dz$
yields
\[
|u(x) - \overline u| \le \fint_{\Ball{\rho}{y}} |u(x) - u(z)|\,dz
\le C_{\mathrm{H\ddot ol}}\,(2\rho)^{\alpha}.
\]
Multiplying by $\rho^{-\alpha}$ gives
$N_{\alpha}(u; y, \rho) \le C_{\mathrm{H\ddot ol}}\, 2^{\alpha}$.
\end{proof}

\begin{theorem}[H\"older implies Campanato seminorm bound,
\leanname{holder\_implies\_campanato}]
If $u$ is $\alpha$-H\"older on $\Ball{R}{x_0}$ with constant $C_{\mathrm{H\ddot ol}}$, then
\[
[u]_{\mathcal{L}^{\alpha}(\Ball{R}{x_0})} \le C_{\mathrm{H\ddot ol}}\, 2^{\alpha}.
\]
\end{theorem}
\begin{proof}
Combine the previous lemma with
\leanname{HasCampanatoBound.campanatoSeminorm\_le}.
\end{proof}


%% ========================================================================
\section{Ball Extension, Approximation, and Supporting Tools}\label{sec:extension}
%% ========================================================================

% Section 4: Ball Extension, Approximation, and Supporting Tools.
% Faithful translation of the De Giorgi formalization.

This section collects the analytic infrastructure underlying the De Giorgi
development: a Sobolev chain rule on open sets, smooth approximation of
$W^{1,p}$ functions on the unit ball, an explicit unit-ball extension operator
together with its $L^p$ and $W^{1,p}$ estimates, the positive-part machinery,
and a small toolkit for $L^p$ limits, finite covers, and pointwise component
analysis. Throughout we work in the Euclidean space $E := \R^d$ with $d \ge 1$.

\subsection{Sobolev chain rule}

The first goal is the chain rule for compositions $\Phi \circ u$ when $u$ is a
$W^{1,2}$ function and $\Phi$ is a smooth scalar function with bounded
derivative.

\begin{theorem}[\leanname{sobolev\_chain\_rule\_unitBall}]
\label{thm:chain-rule-unitBall}
Let $u \in W^{1,2}(\Bz{1})$ with weak gradient $\nabla u$ packaged as a
\leanname{MemW1pWitness}, let $i \in \{1,\dots,d\}$, and let $\Phi \in
C^{\infty}(\R)$ satisfy $\Phi(0)=0$ and $|\Phi'(t)| \le M$ for all $t \in \R$.
Then $\partial_i \Phi(u) = \Phi'(u)\,\partial_i u$ in the weak sense on
$\Bz{1}$.
\end{theorem}

\begin{proof}
Let $M$ be the prescribed bound on $|\Phi'|$, so $M \ge 0$. Since $\Bz{1}$ has
finite Lebesgue measure, by \leanname{exists\_smooth\_W12\_approx\_on\_unitBall}
there is a sequence $\psi_n \in C^{\infty}_c(\R^d)$ with
\begin{align*}
\eLpNormText{\psi_n - u}_{L^2(\Bz{1})} &\to 0,\\
\eLpNormText{\partial_i \psi_n - \partial_i u}_{L^2(\Bz{1})} &\to 0.
\end{align*}
Passing to a subsequence (still denoted $\psi_n$) we may also assume $\psi_n
\to u$ a.e.\ on $\Bz{1}$ via
\leanname{exists\_ae\_tendsto\_of\_eLpNorm\_tendsto}.

Since each $\psi_n$ is smooth and compactly supported, the classical chain
rule and integration by parts yield, for every test function
$\varphi \in C^{\infty}_c(\Bz{1})$,
\begin{multline*}
\int_{\Bz{1}} \Phi(\psi_n)\, \partial_i \varphi \,dx
= -\int_{\Bz{1}} \Phi'(\psi_n)\, \partial_i \psi_n\, \varphi \,dx.
\end{multline*}
We pass to the limit on each side.

\emph{Left-hand side.} The function $\Phi$ is $M$-Lipschitz because $\Phi'$
is bounded and $\Phi$ is differentiable. Hence
$|\Phi(\psi_n(x)) - \Phi(u(x))| \le M\, |\psi_n(x)-u(x)|$ pointwise, so
\[
\eLpNormText{\Phi(\psi_n)-\Phi(u)}_{L^2(\Bz{1})}
\le M\, \eLpNormText{\psi_n - u}_{L^2(\Bz{1})}
\to 0.
\]
The function $\partial_i \varphi$ lies in $L^2(\Bz{1})$ (it is bounded and
compactly supported), so by H\"older's inequality
$\int_{\Bz{1}} \Phi(\psi_n)\,\partial_i \varphi \,dx \to
\int_{\Bz{1}} \Phi(u)\,\partial_i\varphi \,dx$.

\emph{Right-hand side.} Decompose
\[
\Phi'(\psi_n)\,\partial_i \psi_n - \Phi'(u)\,\partial_i u
= \Phi'(\psi_n)\,(\partial_i \psi_n - \partial_i u)
+ (\Phi'(\psi_n) - \Phi'(u))\,\partial_i u.
\]
The first term is bounded in $L^2$-norm by
$M\,\eLpNormText{\partial_i \psi_n - \partial_i u}_{L^2(\Bz{1})}\to 0$.
For the second term, $\Phi'$ is continuous and $\psi_n \to u$ a.e., so
$\Phi'(\psi_n) \to \Phi'(u)$ a.e.\ with uniform bound $M$. The dominated
convergence theorem applied to $|\Phi'(\psi_n) - \Phi'(u)|^2 |\partial_i u|^2$
(dominated by $4M^2 |\partial_i u|^2 \in L^1$) yields convergence of the
second term in $L^2$. Multiplying by $\varphi \in L^\infty_c$ and integrating
gives convergence of the right-hand side to
$-\int_{\Bz{1}} \Phi'(u)\,\partial_i u\, \varphi\,dx$. The chain-rule
identity follows.
\end{proof}

\begin{theorem}[\leanname{sobolev\_chain\_rule}]
\label{thm:chain-rule-open}
Let $\Omega \subseteq \R^d$ be open, $u \in W^{1,2}(\Omega)$, and let $g \in
L^2(\Omega)$ be a representative of $\partial_i u$ in the weak sense. Let
$\Phi \in C^\infty(\R)$ with $\Phi(0)=0$ and $|\Phi'| \le M$. Then
$\Phi'(u)\,g$ is a weak $i$-th partial derivative of $\Phi(u)$ on $\Omega$.
\end{theorem}

\begin{proof}
Let $\nabla u$ be the canonical weak gradient extracted from the witness
$u \in W^{1,2}(\Omega)$, and write $g_i := (\nabla u)_i$. By uniqueness of
weak partial derivatives in $L^1_{\mathrm{loc}}$, we have $g = g_i$ a.e.\ on
$\Omega$, so it suffices to prove the statement with $g_i$ in place of $g$
(the difference of integrands vanishes a.e.\ in the test integral).

Local integrability of $\Phi(u)$ and of $\Phi'(u)\,g_i$ follows from the
$M$-Lipschitz bound $|\Phi(t)| \le M|t|$ and from $|\Phi'(u)\,g_i| \le M|g_i|$,
combined with the local integrability of $u$ and $g_i$.

It now suffices to verify the weak derivative identity locally, that is, on
each ball $\Ball{r}{x_0} \subseteq \Omega$
(\leanname{HasWeakPartialDeriv'\_of\_local}). On such a ball, restrict the
witness to obtain $u \in W^{1,2}(\Ball{r}{x_0})$ and apply
Theorem~\ref{thm:chain-rule-unitBall} (after rescaling to the unit ball via
\leanname{MemW1pWitness.rescale\_to\_unitBall}). The conclusion is the chain
rule for $\Phi(u)$ on $\Ball{r}{x_0}$, and gluing across all balls produces
the desired global identity on $\Omega$.
\end{proof}

\subsection{Smooth approximation on the unit ball}

We now produce smooth, compactly supported approximations of any $W^{1,p}$
function on the unit ball, with control on both function and gradient errors.
The construction proceeds in three layers: (i) a Lipschitz cutoff profile
$\eta_{r,R}$, (ii) inward dilation, and (iii) convolution with a mollifier.

\begin{definition}[\leanname{myCutoff}]
For $x_0 \in E$ and $0 < r < R$, define the radial cutoff profile
\[
\eta_{x_0,r,R}(x) := \mathrm{st}\!\left(\frac{R - \|x - x_0\|}{R - r}\right),
\]
where $\mathrm{st}$ is Mathlib's \emph{smooth transition function}, satisfying
$0 \le \mathrm{st} \le 1$, $\mathrm{st}(t) = 1$ for $t \ge 1$, and
$\mathrm{st}(t) = 0$ for $t \le 0$.
\end{definition}

\begin{lemma}[\leanname{smoothTransition\_nnnorm\_deriv\_bounded}, \leanname{Mst}]
There is a constant $\Mst < \infty$ with $\|\mathrm{st}'(t)\| \le \Mst$ for
all $t \in \R$. Consequently $\mathrm{st}$ is $\Mst$-Lipschitz.
\end{lemma}

\begin{proof}
The derivative $\mathrm{st}'$ is continuous and supported in $[0,1]$, hence
attains a maximum on this compact interval; outside $[0,1]$ it vanishes. Set
$\Mst$ to that maximum value.
\end{proof}

\begin{lemma}[\leanname{myCutoff\_eq\_one}, \leanname{myCutoff\_eq\_zero}, \leanname{myCutoff\_lipschitz}, \leanname{myCutoff\_fderiv\_bound}]
\label{lem:cutoff-properties}
Let $0 < r < R$. Then:
\begin{enumerate}
\item $\eta_{x_0,r,R}(x) = 1$ for $x \in \Ball{r}{x_0}$;
\item $\eta_{x_0,r,R}(x) = 0$ for $x \notin \overline{\Ball{R}{x_0}}$;
\item $\eta_{x_0,r,R}$ is $C^\infty$ on $E$;
\item $\eta_{x_0,r,R}$ is $\Mst (R-r)^{-1}$-Lipschitz;
\item $\|D\eta_{x_0,r,R}(x)\| \le \Mst (R-r)^{-1}$ for all $x \in E$.
\end{enumerate}
\end{lemma}

\begin{proof}
For (i): if $x \in \Ball{r}{x_0}$ then $\|x - x_0\| < r$ and so
$(R - \|x - x_0\|)/(R-r) > 1$, hence $\mathrm{st}$ takes the value $1$.
For (ii): if $\|x-x_0\| > R$ then the argument is negative, hence
$\mathrm{st}$ vanishes. Smoothness follows from
$\mathrm{st} \in C^\infty(\R)$ together with smoothness of the radial
argument away from $x_0$, and constancy near $x_0$. The radial map
$x \mapsto (R - \|x-x_0\|)/(R-r)$ is $1/(R-r)$-Lipschitz, so the composition
with the $\Mst$-Lipschitz $\mathrm{st}$ gives (iv); (v) follows by the
operator-norm bound for derivatives of Lipschitz maps.
\end{proof}

\begin{definition}[\leanname{Cutoff}]
A \emph{cutoff structure} on $\Ball{r}{x_0}$ inside $\Ball{R}{x_0}$ is a tuple
$\eta = (\eta, \mathrm{smooth}, \text{nonneg}, \text{le\_one}, \text{eq\_one},
\text{support}, \text{grad\_bound})$ where $\eta : E \to \R$ is
$C^\infty$, takes values in $[0,1]$, equals $1$ on $\Ball{r}{x_0}$, has
$\tsupport \eta \subseteq \overline{\Ball{R}{x_0}}$, and satisfies the
gradient bound $\|D\eta(x)\| \le \Mst (R-r)^{-1}$.
\end{definition}

\begin{theorem}[\leanname{Cutoff.exists}]
For every $0 < r < R$, the set of cutoff structures on $\Ball{r}{x_0}$ inside
$\Ball{R}{x_0}$ is nonempty.
\end{theorem}

\begin{proof}
The bundle constructed from $\eta_{x_0,r,R}$ together with the lemmas of
\ref{lem:cutoff-properties} is such a cutoff structure.
\end{proof}

\begin{theorem}[\leanname{Lp\_approx\_by\_continuous\_compactly\_supported}]
\label{thm:lp-approx-cc}
Let $1 \le p < \infty$ and $f \in \Lp{p}(\R^d)$. For every $\delta > 0$ there
exists a continuous compactly supported function $g$ with
$\eLpNormText{f-g}_{L^p} < \delta$.
\end{theorem}

\begin{proof}
This is the standard density of $C_c$ in $L^p$ for $p < \infty$. Choose
$0 < \varepsilon < \delta$ in $\R$, apply the Mathlib lemma
\leanname{exist\_eLpNorm\_sub\_le} to obtain a smooth compactly supported
$g$ with $\eLpNormText{f-g}_{L^p} \le \varepsilon < \delta$.
\end{proof}

\subsubsection*{Rescaling}

\begin{definition}[\leanname{unitBallDilate}]
For $\lambda > 0$ and $u : E \to \R$, define
$\unitBallDilate_\lambda u(x) := u(\lambda^{-1} x)$.
\end{definition}

\begin{lemma}[\leanname{smul\_inv\_mem\_unitBall}]
If $\lambda > 1$ and $x \in \Bz{1}$, then $\lambda^{-1} x \in \Bz{1}$.
\end{lemma}

\begin{proof}
$\|\lambda^{-1}x\| = \lambda^{-1}\|x\| \le \|x\| < 1$ since $\lambda^{-1} < 1$.
\end{proof}

\begin{lemma}[\leanname{MemW1pWitness.rescale\_to\_unitBall}]
\label{lem:rescale-witness}
Let $\hw \in W^{1,p}(\Ball{R}{x_0})$ with weak gradient $G$. Then the
function $z \mapsto u(x_0 + Rz)$ is in $W^{1,p}(\Bz{1})$ with weak
gradient $z \mapsto R\,G(x_0 + Rz)$.
\end{lemma}

\begin{proof}
Let $T(z) := x_0 + Rz$, a measurable embedding with
$T^{-1}(\Ball{R}{x_0}) = \Bz{1}$ and pushforward measure
$T_\sharp \mathrm{vol} = R^{-d}\, \mathrm{vol}$. The change of variables
yields $\eLpNormText{u\circ T}_{L^p(\Bz{1})}^p = R^{-d}
\eLpNormText{u}_{L^p(\Ball{R}{x_0})}^p$, hence $u\circ T \in L^p(\Bz{1})$.
For the weak derivative, let $\varphi \in C^\infty_c(\Bz{1})$ and define
$\psi(x) := \varphi(R^{-1}(x - x_0))$. Then $\psi \in C^\infty_c(\Ball{R}{x_0})$
and $\partial_i \varphi(z) = R\, \partial_i \psi(x_0 + Rz)$. Substituting
into the test definition for $G_i$ at $\psi$, performing the change of
variables, and dividing by the Jacobian yields the test definition with
$z \mapsto R\,G_i(x_0+Rz)$ at $\varphi$.
\end{proof}

\begin{lemma}[\leanname{eLpNorm\_rescale\_to\_unitBall}]
For $0 < R$, $p \in [1,\infty)$,
\[
\eLpNormText{u(x_0 + R\cdot)}_{L^p(\Bz{1})}
= R^{-d/p}\,\eLpNormText{u}_{L^p(\Ball{R}{x_0})}.
\]
\end{lemma}

\begin{proof}
Direct application of the change-of-variables formula for the affine map
$T$ above, using $|\det DT| = R^d$.
\end{proof}

\subsubsection*{Inward and outward dilation}

\begin{definition}[\leanname{MemW1pWitness.unitBallDilate}]
Let $\lambda > 1$ and $u \in W^{1,p}(\Bz{1})$. The inward dilate
$\unitBallDilate_\lambda u(x) = u(\lambda^{-1}x)$ defines a function in
$W^{1,p}(\Bz{1})$ obtained by restricting $u$ to $\Ball{\lambda^{-1}}{0}$
and rescaling by Lemma~\ref{lem:rescale-witness}.
\end{definition}

\begin{theorem}[\leanname{exists\_unitBallCutoff\_inside}]
Let $\lambda > 1$. There exists a cutoff structure $\eta$ on
$\Ball{1}{0}$ inside $\Ball{(1+\lambda)/2}{0}$ with
$\tsupport \eta \subseteq \Ball{\lambda}{0}$.
\end{theorem}

\begin{proof}
Apply \leanname{Cutoff.exists} with $r=1$, $R=(1+\lambda)/2$. The support
inclusion follows from $(1+\lambda)/2 < \lambda$.
\end{proof}

\begin{lemma}[\leanname{MemW1pWitness.unitBallDilate\_largeBall}]
\label{lem:dilate-large}
Let $\lambda > 1$ and $u \in W^{1,p}(\Bz{1})$. Then
$\unitBallDilate_\lambda u \in W^{1,p}(\Ball{\lambda}{0})$ with weak
gradient $x \mapsto \lambda^{-1}(\nabla u)(\lambda^{-1}x)$.
\end{lemma}

\begin{proof}
Use the measurable embedding $S(x) := \lambda^{-1}x$, push forward via
$\Ball{\lambda}{0}$, and apply the analogous test-function calculation as
in Lemma~\ref{lem:rescale-witness} with the radial scaling factor
$\lambda^{-1}$.
\end{proof}

\subsubsection*{One-shot smooth approximation}

\begin{theorem}[\leanname{exists\_smooth\_W1p\_oneShot\_on\_unitBall}]
\label{thm:oneshot}
Let $1 < p < \infty$, $u \in W^{1,p}(\Bz{1})$, and $\varepsilon > 0$. Then
there exists $\psi \in C^\infty_c(E)$ with
\begin{align*}
\eLpNormText{\psi - u}_{L^p(\Bz{1})} &< \varepsilon, \\
\eLpNormText{\partial_i \psi - \partial_i u}_{L^p(\Bz{1})} &< \varepsilon
\quad (i=1,\dots,d).
\end{align*}
\end{theorem}

\begin{proof}
Choose $\delta_u, \delta_G > 0$ small enough so that for any
$\lambda \in (1, 1+\delta/2)$ with $\delta := \min(1,\delta_u,\delta_G)$,
the inward dilate $\unitBallDilate_\lambda u$ approximates $u$ in $L^p(\Bz{1})$
to within $\varepsilon/2$, and similarly each rescaled gradient component is
within $\varepsilon/2$ of $(\nabla u)_i$. (Existence of such $\delta$ uses
strong-$L^p$-continuity of dilations.) Set $\lambda := 1 + \delta/2 \in (1,2)$.

Apply Lemma~\ref{lem:dilate-large} to get
$u_\lambda := \unitBallDilate_\lambda u \in W^{1,p}(\Ball{\lambda}{0})$.
Multiply $u_\lambda$ by a cutoff structure $\eta$ on $\Bz{1}$ inside
$\Ball{(1+\lambda)/2}{0}$ from \leanname{exists\_unitBallCutoff\_inside}; this
produces a $W^{1,p}$ function $\eta\, u_\lambda$ supported strictly inside
$\Ball{\lambda}{0}$, agreeing with $u_\lambda$ on $\Bz{1}$.

Convolve $\eta\, u_\lambda$ with a standard mollifier of radius $\rho$ small
enough that the support of the result lies inside $\Ball{\lambda}{0}$ and the
$L^p$ approximation error in both function and gradient is controlled by
$\varepsilon/4$. The result $\psi$ is smooth, compactly supported, and
satisfies, by the triangle inequality,
\begin{align*}
\eLpNormText{\psi - u}_{L^p(\Bz{1})}
&\le \eLpNormText{\psi - \eta u_\lambda}_{L^p(\Bz{1})}
+ \eLpNormText{\eta u_\lambda - u}_{L^p(\Bz{1})} \\
&< \varepsilon/4 + \varepsilon/2 < \varepsilon,
\end{align*}
and analogously for the gradient components, where on $\Bz{1}$ the cutoff
$\eta$ equals $1$, hence $\eta u_\lambda = u_\lambda$ and the gradient
error reduces to a sum of two terms each less than $\varepsilon/2$.
\end{proof}

\begin{theorem}[\leanname{exists\_smooth\_W1p\_approx\_on\_unitBall}]
\label{thm:wp-approx}
Let $1 < p < \infty$ and $u \in W^{1,p}(\Bz{1})$. There exists a sequence
$\psi_n \in C^\infty_c(E)$ such that $\psi_n \to u$ in $L^p(\Bz{1})$ and
$\partial_i \psi_n \to \partial_i u$ in $L^p(\Bz{1})$ for each $i$.
\end{theorem}

\begin{proof}
The one-shot Theorem~\ref{thm:oneshot} produces, for any $\varepsilon > 0$, a
single $\psi \in C^\infty_c(E)$ whose $L^p(\Bz{1})$ distance to $u$ and whose
componentwise $L^p(\Bz{1})$ distance from $\partial_i u$ are both strictly
less than $\varepsilon$. Apply this with the sequence
$\varepsilon_n := (n+1)^{-1}$ and use the axiom of choice
(\leanname{Classical.choose}) to extract a witness
$\psi_n \in C^\infty_c(E)$ for each $n$. By construction, the smoothness and
compact support of $\psi_n$ are inherited from the one-shot output, and we
have the strict bounds
\[
  \eLpNormText{\psi_n - u}_{L^p(\Bz{1})}
    < \mathrm{ENNReal.ofReal}(\varepsilon_n),
\]
together with the analogous bound for each gradient component
$(\fderiv{\psi_n}{x})(e_i) - (\nabla u)_i$.

To upgrade these strict numerical bounds to convergence in $L^p$, observe
that $1/(n+1) \to 0$ in $\R$
(\leanname{tendsto\_one\_div\_add\_atTop\_nhds\_zero\_nat}), so
$\mathrm{ENNReal.ofReal}(\varepsilon_n) \to 0$ in $[0,\infty]$. Each function
error is dominated by these vanishing constants, so the squeeze theorem
\leanname{tendsto\_of\_tendsto\_of\_tendsto\_of\_le\_of\_le}, applied with the
constant zero sequence as a lower bound, yields convergence to $0$. The same
argument applied componentwise to each
$\partial_i \psi_n - (\nabla u)_i$ gives the gradient convergence.
\end{proof}

\begin{theorem}[\leanname{exists\_smooth\_W12\_approx\_on\_unitBall}]
\label{thm:w12-approx}
The previous theorem holds in particular for $p = 2$.
\end{theorem}

\begin{proof}
Direct specialization of Theorem~\ref{thm:wp-approx} using $\ENNRealofReal(2)
= 2$ (after coercing the witness through the trivial isomorphism between the
two encodings of the exponent).
\end{proof}

\subsection{The unit-ball extension operator}

\subsubsection*{Definitions and elementary algebra}

\begin{definition}[\leanname{unitBallRetraction}]
Define $r : E \to E$ by
\[
r(x) := \begin{cases}
x & \text{if } \|x\| \le 1, \\
\|x\|^{-2}\, x & \text{if } 1 < \|x\| < 2, \\
0 & \text{if } \|x\| \ge 2.
\end{cases}
\]
On the annulus $1 < \|x\| < 2$, $r(x)$ is the spherical inversion centred at
the origin with radius $1$.
\end{definition}

\begin{definition}[\leanname{unitBallCutoff}]
Define $\chi : E \to \R$ by $\chi(x) := \min(1, \max(2 - \|x\|, 0))$.
\end{definition}

\begin{definition}[\leanname{unitBallExtension}]
For $u : E \to \R$, define
\[
\Eext u (x) := \chi(x)\, u(r(x)).
\]
\end{definition}

\begin{lemma}[\leanname{unitBallExtension\_zero}, \leanname{unitBallExtension\_add}, \leanname{unitBallExtension\_sub}]
The map $u \mapsto \Eext u$ is $\R$-linear: $\Eext 0 = 0$,
$\Eext{(u+v)} = \Eext u + \Eext v$, and $\Eext{(u-v)} = \Eext u - \Eext v$.
\end{lemma}

\begin{proof}
Pointwise from the definition: each is a multiplication of $\chi(x)$ by a
linear combination of values of $u$ and $v$ at $r(x)$.
\end{proof}

\begin{lemma}[\leanname{unitBallRetraction\_eq\_self\_of\_norm\_le\_one}, \leanname{unitBallRetraction\_eq\_self\_of\_mem\_ball}]
If $\|x\| \le 1$ then $r(x) = x$. In particular if $x \in \Bz{1}$ then
$r(x) = x$.
\end{lemma}

\begin{proof}
Direct from the first branch of the definition.
\end{proof}

\begin{lemma}[\leanname{unitBallExtension\_eq\_of\_mem\_ball}]
\label{lem:Eext-restrict}
If $x \in \Bz{1}$ then $\Eext u(x) = u(x)$.
\end{lemma}

\begin{proof}
Then $r(x) = x$ and $2 - \|x\| > 1$, so $\chi(x) = 1$, giving
$\Eext u (x) = 1 \cdot u(x)$.
\end{proof}

\begin{lemma}[\leanname{unitBallRetraction\_eq\_smul\_of\_annulus}]
If $1 < \|x\| < 2$ then $r(x) = \|x\|^{-2} x$.
\end{lemma}

\begin{proof}
Second branch of the definition.
\end{proof}

\begin{lemma}[\leanname{unitBallCutoff\_eq\_zero\_of\_two\_le\_norm}, \leanname{unitBallExtension\_eq\_zero\_of\_two\_le\_norm}]
If $\|x\| \ge 2$ then $\chi(x) = 0$ and $\Eext u (x) = 0$.
\end{lemma}

\begin{proof}
$2 - \|x\| \le 0$, so $\max(2-\|x\|,0) = 0$ and $\min(1,0) = 0$.
\end{proof}

\begin{lemma}[\leanname{unitBallExtension\_hasCompactSupport}]
$\Eext u$ has compact support contained in $\overline{\Ball{2}{0}}$.
\end{lemma}

\begin{proof}
$\Eext u$ vanishes outside $\overline{\Ball{2}{0}}$ by the previous lemma,
and $\overline{\Ball{2}{0}}$ is compact.
\end{proof}

\subsubsection*{Geometry of inversion}

\begin{definition}[\leanname{unitBallOuterShell}, \leanname{unitBallInnerShell}]
Set
\[
\Sigma^{\mathrm{out}} := \{x \in E : 1 < \|x\| < 2\},
\qquad
\Sigma^{\mathrm{in}} := \{x \in E : 1/2 < \|x\| < 1\}.
\]
Both are open and measurable.
\end{definition}

\begin{lemma}[\leanname{unitBallRetraction\_eq\_inversion\_of\_mem\_outerShell}]
For $x \in \Sigma^{\mathrm{out}}$, $r(x)$ equals the spherical inversion of
$x$ in the unit sphere centred at the origin, i.e. $r(x) = x/\|x\|^2$.
\end{lemma}

\begin{proof}
Combine \leanname{unitBallRetraction\_eq\_smul\_of\_annulus} with the formula
for the Euclidean inversion centred at $0$ with radius $1$, which is
$y \mapsto y/\|y\|^2$.
\end{proof}

\begin{lemma}[\leanname{inversion\_mem\_unitBallInnerShell\_of\_mem\_outerShell}, \leanname{inversion\_mem\_unitBallOuterShell\_of\_mem\_innerShell}]
The Euclidean inversion $I(x) := x/\|x\|^2$ maps
$\Sigma^{\mathrm{out}}$ bijectively onto $\Sigma^{\mathrm{in}}$ and vice
versa.
\end{lemma}

\begin{proof}
If $1 < \|x\| < 2$ then $\|I(x)\| = 1/\|x\| \in (1/2, 1)$, hence
$I(x) \in \Sigma^{\mathrm{in}}$. The opposite direction is symmetric, and
$I \circ I = \mathrm{id}$ on $E\setminus\{0\}$ gives the bijection.
\end{proof}

\begin{lemma}
(\leanname{inversion\_image\_unitBallOuterShell},
\leanname{inversion\_image\_unitBallInnerShell})\\
$I(\Sigma^{\mathrm{out}}) = \Sigma^{\mathrm{in}}$ and
$I(\Sigma^{\mathrm{in}}) = \Sigma^{\mathrm{out}}$.
\end{lemma}

\begin{proof}
Set equalities follow from the previous lemma together with the involution
$I^2 = \mathrm{id}$.
\end{proof}

\begin{lemma}
(\leanname{injOn\_inversion\_unitBallOuterShell},
\leanname{injOn\_inversion\_unitBallInnerShell})\\
The map $I$ is injective on $\Sigma^{\mathrm{out}}$ and on $\Sigma^{\mathrm{in}}$.
\end{lemma}

\begin{proof}
Both shells are subsets of $E\setminus\{0\}$, on which $I$ is a bijective
involution.
\end{proof}

\begin{lemma}[\leanname{hasFDerivWithinAt\_inversion\_unitBallInnerShell}]
For $x \in \Sigma^{\mathrm{in}}$, $I$ is differentiable at $x$ within
$\Sigma^{\mathrm{in}}$, and the Mathlib derivative formula for spherical
inversion applies.
\end{lemma}

\begin{proof}
$x \neq 0$ since $\|x\| > 1/2$, so the standard inversion derivative
(\leanname{EuclideanGeometry.hasFDerivAt\_inversion})
provides a Fr\'echet derivative at $x$, which restricts to the within-set
derivative.
\end{proof}

\begin{lemma}[\leanname{abs\_det\_fderiv\_inversion\_zero\_one}]
\label{lem:det-jacobian}
For $x \in E \setminus \{0\}$,
\[
\bigl|\det DI(x)\bigr| = \bigl(1/\|x\|\bigr)^{2d}.
\]
\end{lemma}

\begin{proof}
By the closed form of the inversion derivative,
\[
DI(x) = (1/\|x\|)^2 \, R_{(\R x)^\perp},
\]
where $R_{(\R x)^\perp}$ is the reflection across the orthogonal complement of
$\R x$. The determinant of this reflection equals $(-1)^{d-1}$ since the
fixed subspace has dimension $d-1$. Therefore
\[
\det DI(x) = (1/\|x\|)^{2d}\,\det R = (-1)^{d-1}(1/\|x\|)^{2d},
\]
whose absolute value is $(1/\|x\|)^{2d}$.
\end{proof}

\begin{lemma}[\leanname{abs\_det\_fderiv\_inversion\_zero\_one\_le\_of\_mem\_innerShell}]
For $x \in \Sigma^{\mathrm{in}}$, $|\det DI(x)| \le 2^{2d}$.
\end{lemma}

\begin{proof}
$\|x\| > 1/2$ gives $1/\|x\| < 2$, so $(1/\|x\|)^{2d} < 2^{2d}$.
\end{proof}

\begin{lemma}[\leanname{lintegral\_outerShell\_comp\_inversion\_le}]
\label{lem:cov-shell}
For any measurable $\Phi : E \to \R_{\ge 0}^\infty$,
\[
\int_{\Sigma^{\mathrm{out}}} \Phi(I(y))\,dy
\le 2^{2d}\,\int_{\Sigma^{\mathrm{in}}} \Phi(x)\,dx.
\]
\end{lemma}

\begin{proof}
By the change-of-variables formula for the differentiable injection $I$ on
$\Sigma^{\mathrm{in}}$ onto $\Sigma^{\mathrm{out}}$,
\[
\int_{\Sigma^{\mathrm{out}}} \Phi(I(y))\,dy
= \int_{\Sigma^{\mathrm{in}}} |\det DI(x)|\, \Phi(I(I(x)))\,dx
= \int_{\Sigma^{\mathrm{in}}} |\det DI(x)|\, \Phi(x)\,dx,
\]
using $I^2 = \mathrm{id}$. Now Lemma~\ref{lem:det-jacobian} together with the
bound from the inner shell gives $|\det DI(x)| \le 2^{2d}$ a.e., and the
result follows.
\end{proof}

\begin{lemma}[\leanname{unitBallCutoff\_eq\_two\_sub\_norm\_of\_mem\_outerShell}]
For $x \in \Sigma^{\mathrm{out}}$, $\chi(x) = 2 - \|x\|$.
\end{lemma}

\begin{proof}
$0 < 2 - \|x\| < 1$, so the $\min$ and $\max$ in the definition reduce to
$2 - \|x\|$.
\end{proof}

\begin{lemma}[\leanname{norm\_fderiv\_inversion\_zero\_one}, \leanname{norm\_fderiv\_inversion\_le\_one\_of\_mem\_outerShell}]
\label{lem:norm-DI}
For $x \neq 0$, $\|DI(x)\| \le (1/\|x\|)^2$. In particular, for
$x \in \Sigma^{\mathrm{out}}$ we have $\|DI(x)\| \le 1$.
\end{lemma}

\begin{proof}
$DI(x) = (1/\|x\|)^2 R$ with $R$ a reflection of operator norm $1$, so
$\|DI(x)\| \le (1/\|x\|)^2$. For $\|x\| \ge 1$ this is at most $1$.
\end{proof}

\begin{lemma}
(\leanname{hasFDerivAt\_unitBallCutoff\_of\_mem\_outerShell},
\leanname{norm\_fderiv\_unitBallCutoff\_of\_mem\_outerShell})\\
For $x \in \Sigma^{\mathrm{out}}$, $\chi$ is differentiable at $x$ with
$D\chi(x) = -D(\|\cdot\|)(x)$, and $\|D\chi(x)\| = 1$.
\end{lemma}

\begin{proof}
On $\Sigma^{\mathrm{out}}$, $\chi$ agrees with $y \mapsto 2 - \|y\|$, which
is smooth on $E\setminus\{0\}$. Differentiating gives
$D\chi = -D\|\cdot\|$, and the norm of the gradient of the Euclidean norm at
any nonzero point equals $1$.
\end{proof}

\begin{lemma}[\leanname{norm\_fderiv\_unitBallExtension\_le\_of\_mem\_outerShell}]
\label{lem:Eext-grad-bound}
Let $u \in C^1(E)$. For $x \in \Sigma^{\mathrm{out}}$,
\[
\|D\Eext u(x)\|
\le |u(I(x))| + \|Du(I(x))\|.
\]
\end{lemma}

\begin{proof}
On $\Sigma^{\mathrm{out}}$, $\Eext u(y) = \chi(y)\, u(I(y))$. By the product
rule and the chain rule,
\[
D\Eext u (x) = \chi(x)\, Du(I(x))\circ DI(x) + u(I(x))\, D\chi(x).
\]
Apply $\|\chi(x)\| \le 1$, $\|DI(x)\| \le 1$ (Lemma~\ref{lem:norm-DI}), and
$\|D\chi(x)\| = 1$ (previous lemma) and the operator-norm submultiplicativity
to obtain
$\|D\Eext u(x)\| \le \|Du(I(x))\| + |u(I(x))|$.
\end{proof}

\begin{lemma}[\leanname{ball\_zero\_two\_eq\_ball\_zero\_one\_union\_outerShell\_union\_sphere}]
The ball $\Ball{2}{0}$ decomposes as
$\Bz{1} \cup \Sigma^{\mathrm{out}} \cup \mathrm{S}_1$, where $\mathrm{S}_1$ is
the unit sphere; the union is disjoint up to the sphere, which has measure
zero.
\end{lemma}

\begin{proof}
A point $y$ with $\|y\| < 2$ is either in $\Bz{1}$ if $\|y\| < 1$, on the
sphere if $\|y\| = 1$, or in $\Sigma^{\mathrm{out}}$ if $1 < \|y\| < 2$.
\end{proof}

\begin{lemma}[\leanname{fderiv\_unitBallExtension\_eq\_of\_mem\_ball}]
For $u$ differentiable and $x \in \Bz{1}$, $D\Eext u (x) = Du(x)$.
\end{lemma}

\begin{proof}
On $\Bz{1}$, $\Eext u$ agrees with $u$ (Lemma~\ref{lem:Eext-restrict}).
Locally constant cutoff and identity retraction give the derivative
identity.
\end{proof}

\subsubsection*{Smooth-input extension estimates}

\begin{theorem}[\leanname{smooth\_unitBallExtension\_W1p\_estimate}]
\label{thm:smooth-extension-estimate}
Let $1 \le p < \infty$ and $u \in C^1(E)$. Then
\begin{align*}
\int_{\Ball{2}{0}} |\Eext u (x)|^p\,dx
&\le (1 + 2^{2d})\, \int_{\Bz{1}} |u(x)|^p\,dx, \\
\int_{\Ball{2}{0}} \|D\Eext u (x)\|^p\,dx
&\le \int_{\Bz{1}} \|Du(x)\|^p\,dx \\
&\quad + 2^{2d}\,2^{p-1}\,\biggl(\int_{\Bz{1}} |u(x)|^p\,dx
+ \int_{\Bz{1}} \|Du(x)\|^p\,dx\biggr).
\end{align*}
\end{theorem}

\begin{proof}
Decompose $\Ball{2}{0} = \Bz{1} \cup \Sigma^{\mathrm{out}} \cup \mathrm{S}_1$,
the sphere having measure zero. On $\Bz{1}$, $\Eext u = u$ and
$D\Eext u = Du$, so the integrals over $\Bz{1}$ contribute exactly the
$u$-side. On $\Sigma^{\mathrm{out}}$, use the shell formula
$\Eext u (x) = (2-\|x\|)\, u(I(x))$ with $|2-\|x\|| \le 1$. By
Lemma~\ref{lem:cov-shell},
\[
\int_{\Sigma^{\mathrm{out}}} |u(I(x))|^p\,dx
\le 2^{2d}\,\int_{\Sigma^{\mathrm{in}}} |u(x)|^p\,dx
\le 2^{2d}\,\int_{\Bz{1}} |u(x)|^p\,dx,
\]
which yields the function-side bound.

For the gradient side, use Lemma~\ref{lem:Eext-grad-bound} to dominate
\[
\|D\Eext u(x)\|^p \le 2^{p-1}\,(|u(I(x))|^p + \|Du(I(x))\|^p)
\quad (x \in \Sigma^{\mathrm{out}}),
\]
via the convexity inequality $(a+b)^p \le 2^{p-1}(a^p+b^p)$. Apply
Lemma~\ref{lem:cov-shell} to each of the two integrands and combine with
the contribution from $\Bz{1}$, where $D\Eext u = Du$.
\end{proof}

\subsubsection*{Smooth approximation of the extension on a rough input}

\begin{definition}[\leanname{unitBallShellFormula}]
For $u : E \to \R$, set
\[
\unitBallShellFormula u(x) := (2 - \|x\|)\, u(I(x)).
\]
For $x \in \Sigma^{\mathrm{out}}$, $\Eext u (x) = \unitBallShellFormula u(x)$.
\end{definition}

\begin{lemma}[\leanname{contDiffOn\_unitBallShellFormula}]
If $u \in C^1(E)$, then the function $\unitBallShellFormula u$ is of
class $C^1$ on $\Sigma^{\mathrm{out}}$.
\end{lemma}

\begin{proof}
$x \mapsto 2 - \|x\|$ and $x \mapsto I(x)$ are $C^1$ on $\Sigma^{\mathrm{out}}$
(both away from $0$), and $u$ is $C^1$ on $E$. Composition and product give
the result.
\end{proof}

\begin{definition}[\leanname{smoothUnitBallExtensionGrad}]
For $u : E \to \R$ smooth, define
$\smoothUnitBallExtensionGrad u (x) \in \R^d$ by its components
\[
  (\smoothUnitBallExtensionGrad u(x))_i := (D\Eext u (x))(e_i).
\]
\end{definition}

\begin{theorem}[\leanname{exists\_global\_smooth\_W1p\_approx\_of\_localizedWitness}]
\label{thm:loc-global}
Let $\Omega \subseteq E$ be open, $1 < p < \infty$, and let $u \in W^{1,p}(\Omega)$
have compact support contained in $\Omega$. Then there is a sequence $\psi_n
\in C^\infty_c(E)$ with
\[
\eLpNormText{\psi_n - u}_{L^p(E)} \to 0
\quad\text{and}\quad
\eLpNormText{\partial_i \psi_n - \mathbf{1}_\Omega (\partial_i u)}_{L^p(E)} \to 0
\]
for each $i$.
\end{theorem}

\begin{proof}
Since $\tsupport u$ is compact and contained in $\Omega$, the criterion
\leanname{memW01p\_of\_memW1p\_of\_tsupport\_subset} promotes the local
$W^{1,p}(\Omega)$ membership of $u$ to a $W^{1,p}_0(\Omega)$ membership: $u$
admits an approximating sequence of compactly supported smooth functions
inside $\Omega$. Pushing this through the zero-extension construction
\leanname{zeroExtend\_memW01p\_p} produces a global $W^{1,p}(E)$ witness for
the trivial extension of $u$ whose weak gradient is $\mathbf{1}_\Omega
\nabla u$, since classical partial derivatives of the zero-extended smooth
approximants vanish off $\Omega$ and converge in $L^p(E)$ to the indicator
of the original gradient.

The global density theorem
Theorem~\ref{thm:smooth_approx_univ}
(\leanname{exists\_smooth\_compactSupport\_W1p\_approx\_univ}) applied to the
extended witness on the open set $E$ then produces a sequence
$\psi_n \in C^\infty_c(E)$ with $\psi_n \to u$ in $L^p(E)$ and
$\partial_i \psi_n \to (\mathbf{1}_\Omega \nabla u)_i$ in $L^p(E)$ for each
component $i$. The compactness of $\tsupport u$ ensures that the gradient
limit is concentrated on $\Omega$, exactly as stated. The Lean source
additionally records a Fatou-type semicontinuity bound: a subsequence of the
$\psi_n$ converges a.e., and Fatou's lemma on $\|\cdot\|^p$ bounds the
$L^p$-energy of the limiting gradient by the liminf of the energies of the
approximants.
\end{proof}

\subsubsection*{Smooth blends and approximation control}

\begin{definition}[\leanname{sphereOneBlend}, \leanname{sphereTwoBlend}]
For $\varepsilon > 0$, define $\sigma^1_\varepsilon(x), \sigma^2_\varepsilon(x)
\in [0,1]$ via the smooth transition
\[
\sigma^1_\varepsilon(x) := \mathrm{st}\!\left(\frac{\|x\|-1+\varepsilon}{2\varepsilon}\right),
\qquad
\sigma^2_\varepsilon(x) := 1 - \mathrm{st}\!\left(\frac{\|x\|-2+\varepsilon}{2\varepsilon}\right),
\]
which transition smoothly across the spheres $\|x\|=1$ and $\|x\|=2$.
\end{definition}

\begin{definition}[\leanname{smoothUnitBallExtensionApprox}]
For $\varepsilon > 0$ and $\psi : E \to \R$ smooth, define
\begin{multline*}
\smoothUnitBallExtensionApprox_\varepsilon \psi (x)
:= \sigma^1_\varepsilon(x) \cdot \unitBallShellFormula \psi(x) \cdot \sigma^2_\varepsilon(x) \\
+ (1 - \sigma^1_\varepsilon(x))\, \psi(x).
\end{multline*}
On the inner core $\|x\| \le 1 - \varepsilon$ this equals $\psi(x)$, on the
outer core $1 + \varepsilon \le \|x\| \le 2 - \varepsilon$ it equals
$\unitBallShellFormula \psi(x)$, and outside $\|x\| \ge 2 + \varepsilon$ it
vanishes. The blend region is contained in two thin annuli around the
spheres of radius $1$ and $2$.
\end{definition}

\begin{lemma}
(\leanname{smoothUnitBallExtensionApprox\_contDiff},
\leanname{smoothUnitBallExtensionApprox\_hasCompactSupport})\\
For $\psi$ smooth and $0 < \varepsilon < 1$, $\smoothUnitBallExtensionApprox_\varepsilon
\psi$ is $C^\infty$ and has compact support contained in
$\overline{\Ball{2+\varepsilon}{0}}$.
\end{lemma}

\begin{proof}
Each of the two summands is smooth: the shell formula factor is smooth on
$E \setminus\{0\}$ and is multiplied by $\sigma^1_\varepsilon$, which
vanishes near the origin (since $\sigma^1_\varepsilon(x) = 0$ when
$\|x\| \le 1-\varepsilon$). The compact support follows from
$\sigma^2_\varepsilon$ vanishing for $\|x\| \ge 2+\varepsilon$.
\end{proof}

\begin{definition}[\leanname{unitBallApproxEps}]
$\varepsilon_n := (n+5)^{-1}$.
\end{definition}

\begin{lemma}[\leanname{unitBallApproxEps\_pos}, \leanname{unitBallApproxEps\_lt\_one}, \leanname{unitBallApproxEps\_le\_one\_fifth}, \leanname{tendsto\_unitBallApproxEps}, \leanname{unitBallApproxEps\_antitone}]
$\varepsilon_n > 0$, $\varepsilon_n < 1$, $\varepsilon_n \le 1/5$, $\varepsilon_n
\to 0$, and $(\varepsilon_n)$ is antitone.
\end{lemma}

\begin{proof}
All follow from $n \mapsto (n+5)^{-1}$.
\end{proof}

\begin{definition}[\leanname{unitBallBadAnnulusOne}, \leanname{unitBallBadAnnulusTwo}]
$\mathrm{Ann}^1_\varepsilon := \{x : 1-\varepsilon < \|x\| < 1 + \varepsilon\}$
and $\mathrm{Ann}^2_\varepsilon := \{x : 2-\varepsilon < \|x\| < 2+\varepsilon\}$.
Both are measurable, and antitone in $\varepsilon$ (with respect to inclusion
on $(0,1)$).
\end{definition}

\begin{lemma}
(\leanname{iInter\_unitBallBadAnnulusOne\_eq\_sphere},
\leanname{iInter\_unitBallBadAnnulusTwo\_eq\_sphere},
\leanname{tendsto\_measure\_unitBallBadAnnulusOne},
\leanname{tendsto\_measure\_unitBallBadAnnulusTwo})\\
$\bigcap_{n} \mathrm{Ann}^j_{\varepsilon_n} = \mathrm{S}_j$ for $j = 1, 2$, and
$|\mathrm{Ann}^j_{\varepsilon_n}| \to 0$ as $n \to \infty$.
\end{lemma}

\begin{proof}
A point lies in every $\mathrm{Ann}^j_{\varepsilon_n}$ iff $\|x\| = j$, since
the annulus shrinks to the sphere as $\varepsilon \to 0$. The measure of the
sphere is zero (Mathlib: \leanname{Measure.addHaar\_sphere}), so by continuity
of measure from above the annulus measures vanish.
\end{proof}

\begin{definition}[\leanname{sphereOneControl}, \leanname{sphereTwoControl}]
Compact buffer sets around the unit and double spheres,
$\mathrm{Buf}^1 := \{x : 4/5 \le \|x\| \le 6/5\}$ and
$\mathrm{Buf}^2 := \{x : 9/5 \le \|x\| \le 11/5\}$. Both are compact.
\end{definition}

\begin{lemma}[\leanname{badAnnulusOne\_subset\_sphereOneControl}, \leanname{badAnnulusTwo\_subset\_sphereTwoControl}, \leanname{isCompact\_sphereOneControl}, \leanname{isCompact\_sphereTwoControl}]
For all $n$, $\mathrm{Ann}^j_{\varepsilon_n} \subseteq \mathrm{Buf}^j$, and
both buffers are compact.
\end{lemma}

\begin{proof}
$\varepsilon_n \le 1/5$, so $j - \varepsilon_n \ge j - 1/5$ and similarly
above. Compactness is by closedness and boundedness.
\end{proof}

\begin{theorem}[\leanname{exists\_fderiv\_bound\_on\_compact}, \leanname{exists\_norm\_bound\_on\_compact}]
For any continuous $f$ (resp.\ $C^1$ $f$) on $E$, $f$ (resp.\ $\|Df\|$) is
bounded on every compact set.
\end{theorem}

\begin{proof}
Continuous functions attain bounded values on compact sets.
\end{proof}

\begin{theorem}[\leanname{exists\_shellFormula\_fderiv\_bound}]
For any smooth $\psi$, the derivative $D \unitBallShellFormula \psi$ is bounded
on the compact set $\mathrm{Buf}^2$, hence on each $\mathrm{Ann}^2_{\varepsilon_n}$.
\end{theorem}

\begin{proof}
Smoothness of $\unitBallShellFormula \psi$ on $\Sigma^{\mathrm{out}}$ extends
to smoothness on a neighborhood of $\mathrm{Buf}^2$, where $\|x\| \in
[9/5, 11/5] \subseteq (0, 2 + \delta)$, so the result follows from
\leanname{exists\_norm\_bound\_on\_compact}.
\end{proof}

\begin{lemma}[\leanname{tendsto\_smoothUnitBallExtensionApprox\_pointwise}, \leanname{tendsto\_fderiv\_smoothUnitBallExtensionApprox\_pointwise}]
\label{lem:smoothapprox-pointwise}
For $\psi \in C^\infty(E)$ and $x$ outside the unit and double spheres,
$\smoothUnitBallExtensionApprox_{\varepsilon_n} \psi (x) \to \Eext \psi (x)$
and $D\smoothUnitBallExtensionApprox_{\varepsilon_n} \psi (x) \to D\Eext\psi(x)$.
\end{lemma}

\begin{proof}
Eventually $x$ lies in either $\Bz{1}$ (so the inner branch picks up $\psi(x)
= \Eext \psi(x)$), or in $\Sigma^{\mathrm{out}}$ separated from both spheres
(so the outer branch agrees with $\unitBallShellFormula \psi(x) = \Eext\psi(x)$
and similarly for derivatives), or in $\|x\| > 2$ where the cutoff vanishes.
The blends become trivial in the limit and the derivative converges by
smoothness.
\end{proof}

\begin{definition}[\leanname{unitBallBadSet}]
$\mathrm{Bad}_\varepsilon := \mathrm{Ann}^1_\varepsilon \cup
\mathrm{Ann}^2_\varepsilon$.
\end{definition}

\begin{lemma}[\leanname{tendsto\_measure\_unitBallBadSet}]
$|\mathrm{Bad}_{\varepsilon_n}| \to 0$.
\end{lemma}

\begin{proof}
Sub-additivity together with the previous annulus convergence.
\end{proof}

\begin{lemma}[\leanname{smoothUnitBallExtensionApprox\_eq\_unitBallExtension\_of\_not\_mem\_badSet}]
For $\psi$ smooth and $x \notin \mathrm{Bad}_{\varepsilon_n}$,
$\smoothUnitBallExtensionApprox_{\varepsilon_n} \psi (x) = \Eext \psi(x)$.
\end{lemma}

\begin{proof}
Outside the bad set, $x$ falls cleanly into one of: inner core, outer core,
or the exterior $\|x\|>2+\varepsilon$, on each of which the blends become
constant ($0$ or $1$) and the formula collapses to $\Eext\psi$.
\end{proof}

\begin{theorem}[\leanname{exists\_shellFormula\_error\_bound}]
\label{thm:err-bound}
For any smooth $\psi$, there is a constant $K$ depending only on $\psi$
such that
\[
\bigl|\smoothUnitBallExtensionApprox_{\varepsilon_n} \psi(x) - \Eext\psi(x)\bigr|
\le K\,\varepsilon_n
\]
for all $x \in E$.
\end{theorem}

\begin{proof}
Outside $\mathrm{Bad}_{\varepsilon_n}$, both sides agree (so the error is
zero). On $\mathrm{Bad}_{\varepsilon_n}$, both sides differ from each other
only in the blend region, which has width $\sim \varepsilon_n$, and within a
ball of radius $\varepsilon_n$ around any point of the bad annulus the
shell formula is Lipschitz with constant uniformly controlled by the
$C^1$-bound of $\psi$ on the buffer compact $\mathrm{Buf}^2$. Taking the
mean-value bound on segments inside $\mathrm{Buf}^j$ yields the linear
control.
\end{proof}

\begin{theorem}[\leanname{exists\_fun\_error\_bound\_badSet}]
For each smooth $\psi$, there is $K$ such that
\[
\eLpNormText{\smoothUnitBallExtensionApprox_{\varepsilon_n}\psi - \Eext\psi}_{L^p(E)}
\le K\,\varepsilon_n\,|\mathrm{Bad}_{\varepsilon_n}|^{1/p}
\to 0.
\]
\end{theorem}

\begin{proof}
From Theorem~\ref{thm:err-bound} the pointwise difference
$|\smoothUnitBallExtensionApprox_{\varepsilon_n}\psi(x) - \Eext\psi(x)|$ is
bounded by $C\,\varepsilon_n \cdot \mathbf{1}_{\mathrm{Bad}_{\varepsilon_n}}(x)$
for some constant $C$ depending only on $\psi$. Pointwise both sides are
also dominated by the constant $C$ on the closed ball $\overline{\Ball{9/4}{0}}$:
outside the ball of radius $9/4$ both
$\smoothUnitBallExtensionApprox_{\varepsilon_n}\psi$ and $\Eext\psi$ vanish
(using \leanname{sphereOneControl\_subset\_closedBall},
\leanname{sphereTwoControl\_subset\_closedBall} for the bad annuli together
with \leanname{unitBallExtension\_hasCompactSupport}). Hence the family of
differences is uniformly $L^p$-dominated by the constant times indicator
$C\,\mathbf{1}_{\overline{\Ball{9/4}{0}}}$, which lies in $L^p(E)$ since the
underlying ball has finite Lebesgue measure
(\leanname{measure\_closedBall\_lt\_top}).

Since the bad annuli shrink to the spheres of radii $1$ and $2$ and the Haar
measure of those spheres is zero (\leanname{Measure.addHaar\_sphere}),
\leanname{tendsto\_measure\_unitBallBadSet} gives
$|\mathrm{Bad}_{\varepsilon_n}| \to 0$. Combining the pointwise bound with
the dominated convergence theorem in $L^p$, against the constant dominating
indicator, yields
\[
  \eLpNormText{\smoothUnitBallExtensionApprox_{\varepsilon_n}\psi
    - \Eext\psi}_{L^p(E)}
  \le K\,\varepsilon_n\,|\mathrm{Bad}_{\varepsilon_n}|^{1/p}
  \to 0,
\]
where $K$ is the constant from Theorem~\ref{thm:err-bound}, and the
$|\mathrm{Bad}_{\varepsilon_n}|^{1/p}$ factor arises from H\"older's inequality
applied to the indicator on the bad set.
\end{proof}

\subsubsection*{The smoothing theorem on a smooth input}

\begin{theorem}[\leanname{smooth\_input\_unitBallExtension\_smoothing}]
\label{thm:smooth-input-smoothing}
Let $\psi \in C^\infty(E)$ have compact support, $1 < p < \infty$. Then there
is a sequence
\[
  \Psi_n := \smoothUnitBallExtensionApprox_{\varepsilon_n}\psi \in C^\infty_c(E)
\]
such that:
\begin{enumerate}
\item $\Psi_n \to \Eext \psi$ in $L^p(E)$;
\item $\partial_i \Psi_n \to (\smoothUnitBallExtensionGrad\psi)_i$ in $L^p(E)$
for each $i$;
\item the $L^p$ norm of $\Psi_n$ is bounded uniformly by the smooth-input
extension estimate of Theorem~\ref{thm:smooth-extension-estimate}.
\end{enumerate}
\end{theorem}

\begin{proof}
The first item is Theorem~\ref{thm:err-bound} together with control of the
bad set's measure. For the second, the same blend-region argument applies to
$D\Psi_n - D\Eext\psi$: by Lemma~\ref{lem:smoothapprox-pointwise}, the
derivatives converge pointwise outside the spheres of radii $1$ and $2$ to
the smooth extension derivatives, while the bad-set contributions are
controlled by the uniform $C^1$ bounds on the buffer compacts. Item (3) is
an application of Theorem~\ref{thm:smooth-extension-estimate} to $\psi$.
\end{proof}

\begin{theorem}[\leanname{hasWeakPartials\_of\_global\_smoothApprox}]
Let $1 < p < \infty$, let $f \in L^p(E)$ have a candidate gradient $G \in
L^p(E; \R^d)$, and assume there exist $\Psi_n \in C^\infty_c(E)$ with
$\Psi_n \to f$ and $\partial_i \Psi_n \to G_i$ in $L^p(E)$ for each $i$. Then
$G_i$ is the weak $i$-th partial derivative of $f$ on $E$.
\end{theorem}

\begin{proof}
For $\varphi \in C^\infty_c(E)$, the integration-by-parts identity
$\int \Psi_n\, \partial_i \varphi = -\int (\partial_i \Psi_n)\, \varphi$
holds for each smooth $\Psi_n$. Pass to the limit using $L^p \times L^q$
duality with $1/p + 1/q = 1$ (note $\partial_i\varphi \in L^q$ and
$\varphi \in L^q$ since both are bounded compactly supported), giving
$\int f\, \partial_i \varphi = -\int G_i\, \varphi$, the weak derivative
identity.
\end{proof}

\subsubsection*{The main extension theorems}

\begin{theorem}[\leanname{exists\_smooth\_global\_approx\_of\_unitBallExtension}]
\label{thm:global-approx-extension}
Let $1 < p < \infty$ and $u \in W^{1,p}(\Bz{1})$. There exists
$G_{\mathrm{ext}} : E \to \R^d$ with the following properties.
\begin{enumerate}
\item $\Eext u \in L^p(E)$ and each component of $G_{\mathrm{ext}}$ lies in
$L^p(E)$.
\item Function-side estimate:
\[
\int_E |\Eext u(x)|^p\,dx
\le \CunitBallExtensionFun(d)\,\int_{\Bz{1}} |u(x)|^p\,dx,
\]
with $\CunitBallExtensionFun(d) := 1 + 2^{2d}$.
\item Gradient-side estimate:
\begin{multline*}
\int_E \|G_{\mathrm{ext}}(x)\|^p\,dx
\le \int_{\Bz{1}} \|\nabla u(x)\|^p\,dx \\
+ \CunitBallExtensionGrad(d, p)\,\biggl(\int_{\Bz{1}} |u(x)|^p\,dx
+ \int_{\Bz{1}} \|\nabla u(x)\|^p\,dx\biggr),
\end{multline*}
with $\CunitBallExtensionGrad(d,p) := 2^{2d}\cdot 2^{p-1}$.
\item There exists a sequence $\Phi_n \in C^\infty_c(E)$ with
$\Phi_n \to \Eext u$ and $\partial_i \Phi_n \to (G_{\mathrm{ext}})_i$ in
$L^p(E)$.
\end{enumerate}
\end{theorem}

\begin{proof}
Apply Theorem~\ref{thm:wp-approx} to obtain $\psi_n \in C^\infty_c(E)$ with
$\psi_n \to u$ and $\partial_i \psi_n \to (\nabla u)_i$ in $L^p(\Bz{1})$.
For each $\psi_n$, apply Theorem~\ref{thm:smooth-input-smoothing}, producing
a smooth extension sequence $\Phi_{n,m} \in C^\infty_c(E)$ approximating
$\Eext \psi_n$ in $W^{1,p}(E)$, and the smooth-input estimates of
Theorem~\ref{thm:smooth-extension-estimate}.

The smooth-input bound gives, by linearity of $\Eext$,
\begin{gather*}
\eLpNormText{\Eext \psi_n - \Eext \psi_m}_{L^p(E)}
\le (1 + 2^{2d})^{1/p}\,\eLpNormText{\psi_n - \psi_m}_{L^p(\Bz{1})},
\end{gather*}
\begin{multline*}
\eLpNormText{D\Eext \psi_n - D\Eext\psi_m}_{L^p(E)}
\le \eLpNormText{D\psi_n - D\psi_m}_{L^p(\Bz{1})} \\
+ C\,\bigl(\eLpNormText{\psi_n - \psi_m}_{L^p(\Bz{1})}
+ \eLpNormText{D\psi_n - D\psi_m}_{L^p(\Bz{1})}\bigr).
\end{multline*}
Since $(\psi_n)$ is Cauchy in $W^{1,p}(\Bz{1})$, the smooth-extension
sequence $\Eext\psi_n$ is Cauchy in $L^p(E)$, and the candidate gradient
sequence $\smoothUnitBallExtensionGrad \psi_n$ is Cauchy in
$L^p(E;\R^d)$. By the bare-function $L^p$ completeness theorems
(Theorems~\leanname{scalar\_cauchy\_to\_limit} and
\leanname{exists\_pi\_limit\_of\_cauchy\_eLpNorm}), there exist $L^p$ limits $f^*$
and $G_{\mathrm{ext}}$. By construction $f^* = \Eext u$ a.e.\ and the
estimates pass to the limit by lower semicontinuity of the $L^p$-norm under
weak/strong convergence and the smooth-input estimates applied to each
$\psi_n$. Finally, a diagonal extraction from $(\Phi_{n,m})$ produces the
sequence required in (4).
\end{proof}

\begin{lemma}[\leanname{aestronglyMeasurable\_unitBallExtension\_of\_memLp}]
If $u \in L^p(\Bz{1})$ then $\Eext u$ is a.e.\ strongly measurable on $E$.
\end{lemma}

\begin{proof}
Choose a strongly measurable representative $u'$ of $u$. Then $\Eext u'$ is
measurable: it is a composition of measurable maps. Now $\Eext u =
\Eext u'$ a.e.\ on $E$, by checking three cases.
\begin{itemize}
\item On $\Bz{1}$, $\Eext = \mathrm{id}$, so equality reduces to $u = u'$
a.e.
\item On $\|x\| \ge 2$, both vanish.
\item On $\Sigma^{\mathrm{out}}$, $\Eext u(x) = \chi(x)\, u(I(x))$. The
inversion $I$ is a smooth diffeomorphism between $\Sigma^{\mathrm{out}}$ and
$\Sigma^{\mathrm{in}}$. The set where $u \neq u'$ on $\Sigma^{\mathrm{in}}$
has measure zero, and its image under $I$ also has measure zero by
\leanname{addHaar\_image\_eq\_zero\_of\_differentiableOn\_of\_addHaar\_eq\_zero}.
\end{itemize}
The sphere $\mathrm{S}_1$ has measure zero, hence does not affect the a.e.
identity. Strong measurability of $\Eext u$ then follows from a.e. equality
to a measurable function.
\end{proof}

\begin{theorem}[\leanname{exists\_unitBall\_W1p\_extension}]
\label{thm:Eext-W1p}
Let $1 < p < \infty$ and $u \in W^{1,p}(\Bz{1})$. There is a $W^{1,p}$
witness $\hwExt$ for $\Eext u$ on $E$ with the following properties:
\begin{enumerate}
\item $\Eext u$ has compact support in $\overline{\Ball{2}{0}}$;
\item $\Eext u(x) = u(x)$ for $x \in \Bz{1}$;
\item the function-side bound from Theorem~\ref{thm:global-approx-extension} (2);
\item the gradient-side bound from Theorem~\ref{thm:global-approx-extension} (3).
\end{enumerate}
\end{theorem}

\begin{proof}
Combine Theorem~\ref{thm:global-approx-extension} (existence of
$G_{\mathrm{ext}}$ together with the approximating sequence) with the previous
two lemmas. The witness package follows from
\leanname{hasWeakPartials\_of\_global\_smoothApprox}, the support inclusion from
\leanname{unitBallExtension\_hasCompactSupport}, and the agreement on $\Bz{1}$
from \leanname{unitBallExtension\_eq\_of\_mem\_ball}.
\end{proof}

\begin{theorem}[\leanname{exists\_unitBall\_W1p\_extension'}]
A weakened formulation: for $u \in W^{1,p}(\Bz{1})$, there exists $u_{\mathrm{ext}}
\in W^{1,p}(E)$ with compact support and $u_{\mathrm{ext}} = u$ on $\Bz{1}$.
\end{theorem}

\begin{proof}
Project away the explicit constants from Theorem~\ref{thm:Eext-W1p}.
\end{proof}

\subsection{Positive-part machinery}

\begin{theorem}[\leanname{MemW1pWitness.weakGrad\_ae\_eq\_zero\_on\_zeroSet}]
Let $u \in W^{1,2}(\Omega)$ on an open set $\Omega \subseteq E$. Then for
each $i$, $(\nabla u)_i = 0$ a.e.\ on the zero set $\{u = 0\}\cap \Omega$.
\end{theorem}

\begin{proof}
Choose a measurable representative $u'$ of $u$. Apply Stampacchia's zero-set
theorem (\leanname{HasWeakPartialDeriv.toStampacchia}) to the witness with
the measurable representative, then transfer back via a.e.\ equality.
\end{proof}

\begin{lemma}[\leanname{MemW1pWitness.sub\_const}]
On a finite-measure open set $\Omega$, if $u \in W^{1,2}(\Omega)$ then
$u - c \in W^{1,2}(\Omega)$ with the same weak gradient.
\end{lemma}

\begin{proof}
Constants have weak gradient zero on $\Omega$. Subtract pointwise and use
that the test integral splits, the contribution from the constant being
zero.
\end{proof}

\begin{lemma}[\leanname{MemW1pWitness.mul\_smooth\_bounded}]
Let $u \in W^{1,2}(\Omega)$ and $\eta \in C^\infty(E)$ with
$|\eta| \le C_0$ and $\|D\eta\| \le C_1$. Then $\eta\, u \in W^{1,2}(\Omega)$
with weak gradient $\eta\, \nabla u + u\, \nabla \eta$.
\end{lemma}

\begin{proof}
The function-side $L^2$ membership follows from $|\eta\, u| \le C_0\,|u|$.
For each component, the candidate $(\eta\, (\nabla u)_i + (\partial_i \eta)\,u)$
is in $L^2$ by similar bounds. The weak-derivative identity follows from
the product rule for weak derivatives, \leanname{HasWeakPartialDeriv.mul\_smooth}.
\end{proof}

\begin{theorem}[\leanname{HasWeakPartialDeriv.of\_eLpNormApprox}]
Let $\Omega$ be open, $i$ an index, $f \in L^2(\Omega)$ with candidate
$g \in L^2(\Omega)$, and let $\psi_n$ be smooth functions whose pointwise
weak partial derivatives $g_n$ converge to $g$ in $L^2(\Omega)$ and whose
values converge to $f$ in $L^2(\Omega)$. Then $g$ is the weak $i$-th
partial derivative of $f$ on $\Omega$.
\end{theorem}

\begin{proof}
Fix a smooth compactly supported test $\varphi$ on $\Omega$, and write
$d\varphi(x) := (\fderiv{\varphi}{x})(e_i)$. The integration-by-parts identity for the
smooth approximants reads
\[
  \int_{\Omega} \psi_n \cdot d\varphi\,dx
    = -\int_{\Omega} g_n \cdot \varphi\,dx,
\]
which is built into the hypothesis $\mathrm{HasWeakPartialDeriv}\,i\,g_n\,\psi_n\,\Omega$.

Both $\varphi$ and $d\varphi$ are continuous with compact support, hence in
$L^2(\Omega)$. By the elementary $L^2$-pairing limit
\leanname{tendsto\_integral\_mul\_of\_eLpNorm\_tendsto\_zero}: if
$h_n - h \to 0$ in $L^2$ and $k$ is a fixed $L^2$ function, then
$\int k\,(h_n - h) \to 0$. Applying this with $k = d\varphi$ and
$h_n - h = \psi_n - f$ gives
\[
  \int_{\Omega} d\varphi\,(\psi_n - f)\,dx \to 0,
\]
and similarly with $k = \varphi$ and $h_n - h = g_n - g$ gives
$\int_{\Omega} \varphi\,(g_n - g)\,dx \to 0$. Adding the IBP identity and
$\int f\cdot d\varphi$ to the first limit gives
\[
  \int_{\Omega} \psi_n\cdot d\varphi\,dx \to \int_{\Omega} f \cdot d\varphi\,dx.
\]
Likewise the right-hand side of IBP converges to $-\int_{\Omega} g\cdot\varphi\,dx$.
Passing to the limit in the smooth IBP identity therefore yields
$\int f\,d\varphi = -\int g\,\varphi$, which is the defining identity for $g$
to be the weak $i$-th partial derivative of $f$ on $\Omega$.
\end{proof}

\begin{theorem}[\leanname{tendsto\_eLpNorm\_indicator\_diff\_mul\_of\_tendsto\_eLpNorm}]
Let $\psi_n \to u$ in $L^2(\Omega)$, with $u \in L^2(\Omega)$ and $G \in
L^2(\Omega)$ such that $G$ vanishes a.e.\ on $\{u = 0\}$. Then
\[
\eLpNormText{(\mathbf{1}_{\{\psi_n > 0\}} - \mathbf{1}_{\{u > 0\}})\,G}_{L^2(\Omega)}
\to 0.
\]
\end{theorem}

\begin{proof}
Write $F_n(x) := (\mathbf{1}_{\{\psi_n > 0\}}(x) - \mathbf{1}_{\{u > 0\}}(x))\,G(x)$.
The pointwise inequality $|F_n(x)| \le |G(x)|$ is immediate from the fact that
each indicator lies in $\{0,1\}$, so the family $(F_n)$ is uniformly dominated
by $|G| \in L^2(\Omega)$. Each $F_n$ is a.e.\ strongly measurable: the sets
$\{\psi_n > 0\}$ and $\{u > 0\}$ are null measurable since $\psi_n$ and $u$
are a.e.\ strongly measurable (the comparison with the constant $0$ is
handled by \leanname{AEStronglyMeasurable.nullMeasurableSet\_lt}), and $G$
is in $L^2$.

Vitali convergence in $L^2$ requires uniform integrability and uniform
tightness of the family $(F_n)$. Both are inherited from the constant family
$\{G\}$, which is uniformly integrable and uniformly tight by
\leanname{MeasureTheory.unifIntegrable\_const} and
\leanname{MeasureTheory.unifTight\_const} applied to the singleton
$\{G\} \subset L^2$. The pointwise domination $|F_n| \le |G|$ then transfers
both properties to $(F_n)$ via $\eLpNormText{\mathbf{1}_S F_n}_{L^2} \le
\eLpNormText{\mathbf{1}_S G}_{L^2}$ for any measurable set $S$.

It remains to show $F_n \to 0$ in measure. By \leanname{tendsto\_of\_subseq\_tendsto},
it suffices that every subsequence of $(F_n)$ admits an a.e.\ convergent
sub-subsequence with limit $0$. Given any subsequence, the $L^2$ hypothesis
$\psi_{n_k} \to u$ in $L^2$ supplies an a.e.\ convergent sub-sub-sequence
(\leanname{exists\_subseq\_tendsto\_ae\_of\_tendsto\_eLpNorm}). On the
sub-subsequence: at points where $u(x) > 0$, eventually $\psi_{n_{k_j}}(x) > 0$,
so both indicators equal $1$ and $F_{n_{k_j}}(x) \to 0$; at points where
$u(x) < 0$ both indicators eventually equal $0$; at points where $u(x) = 0$
the hypothesis hzero gives $G(x) = 0$ a.e., so $F_{n_{k_j}}(x) = 0$.
This produces an a.e.\ vanishing sub-subsequence. Combining with the Vitali
convergence theorem (uniform integrability + uniform tightness +
convergence in measure $\Rightarrow$ $L^2$ convergence) yields the desired
$\eLpNormText{F_n}_{L^2} \to 0$.
\end{proof}

\begin{definition}[\leanname{MemW1pWitness.posPart\_of\_aux}]
Given $u \in W^{1,2}(\Omega)$ with smooth approximation data $(\psi_n)$ and
the zero-set vanishing hypothesis from
\leanname{MemW1pWitness.weakGrad\_ae\_eq\_zero\_on\_zeroSet}, the positive part
$u^+ := \max(u,0) \in W^{1,2}(\Omega)$ has weak gradient
\[
(\nabla u^+)_i(x) := \begin{cases}
(\nabla u)_i(x) & \text{if } u(x) > 0, \\
0 & \text{otherwise}.
\end{cases}
\]
\end{definition}

\begin{theorem}[\leanname{MemW1pWitness.posPart}]
On a finite-measure open set $\Omega$, if $u \in W^{1,2}(\Omega)$ comes
equipped with a smooth approximation, then $u^+ \in W^{1,2}(\Omega)$ with
the gradient described above.
\end{theorem}

\begin{proof}
Apply \leanname{posPart\_of\_aux} with the zero-set vanishing input supplied by
\leanname{weakGrad\_ae\_eq\_zero\_on\_zeroSet}. The proof of the auxiliary
theorem proceeds in three steps.
\begin{enumerate}
\item Define $g(x) := \mathbf{1}_{\{u > 0\}}\, (\nabla u)_i(x) \in L^2(\Omega)$
and $g^{\psi_n}(x) := \mathbf{1}_{\{\psi_n > 0\}}\,(\partial_i \psi_n)(x)$.
\item Each $g^{\psi_n}$ is the classical weak partial derivative of
$\max(\psi_n, 0)$ by
\leanname{weakPartialDeriv\_positivePart\_of\_contDiff} (since $\psi_n$ is
$C^1$).
\item Show $\max(\psi_n, 0) \to u^+$ in $L^2(\Omega)$ by
\leanname{tendsto\_eLpNorm\_positivePart\_sub} and $g^{\psi_n} \to g$ in
$L^2(\Omega)$ by combining
\leanname{tendsto\_eLpNorm\_indicator\_diff\_mul\_of\_tendsto\_eLpNorm} (which uses
the zero-set hypothesis) with the $L^2$ convergence
$\partial_i \psi_n \to (\nabla u)_i$.
\item Apply Theorem~\leanname{HasWeakPartialDeriv.of\_eLpNormApprox} to deduce
that $g$ is a weak $i$-th partial derivative of $u^+$ on $\Omega$.
\end{enumerate}
\end{proof}

\subsection{$L^p$ function toolkit}

\begin{theorem}[\leanname{memLp\_of\_tendsto\_eLpNorm}]
Let $1 \le p \le \infty$, let $f_n \in L^p(\mu)$, and let $g$ be a.e.
strongly measurable. If $\eLpNormText{f_n - g}_{L^p} \to 0$, then
$g \in L^p(\mu)$.
\end{theorem}

\begin{proof}
For some $N$, $\eLpNormText{f_N - g}_{L^p} < 1$. The triangle inequality
gives $\eLpNormText{g}_{L^p} \le \eLpNormText{f_N}_{L^p} + 1 < \infty$.
\end{proof}

\begin{theorem}[\leanname{ae\_eq\_of\_tendsto\_eLpNorm\_sub}]
Lp limits are unique up to a.e.\ equality.
\end{theorem}

\begin{proof}
$\eLpNormText{g_1 - g_2}_{L^p} \le \eLpNormText{f_n - g_1}_{L^p} +
\eLpNormText{f_n - g_2}_{L^p} \to 0$, so $\eLpNormText{g_1 - g_2}_{L^p} = 0$,
and $g_1 - g_2 = 0$ a.e.
\end{proof}

\begin{theorem}[\leanname{scalar\_cauchy\_to\_limit}]
\label{thm:scalar-cauchy}
Let $1 \le p < \infty$, $E$ second-countable complete, and let $f_n \in
L^p(\mu; E)$ be a Cauchy sequence in $L^p$. There exists $g \in L^p(\mu; E)$
with $\eLpNormText{f_n - g}_{L^p}\to 0$.
\end{theorem}

\begin{proof}
Standard Riesz-Fischer construction. Extract a subsequence $f_{M(k)}$ with
$\eLpNormText{f_{M(k)} - f_{M(k+1)}}_{L^p} \le 2^{-(k+1)}$ via a controlled
$L^p$-Cauchy sub-extraction. Apply Mathlib's
\leanname{Lp.cauchy\_complete\_eLpNorm} to obtain a limit $g$ for the
subsequence. Upgrade to convergence of the full sequence by combining the
sub-sequence convergence with the original Cauchy property and a $\varepsilon/2$
triangle argument.
\end{proof}

\begin{theorem}
(\leanname{eLpNorm\_pi\_le\_sum\_component},
\leanname{eLpNorm\_pi\_component\_le},
\leanname{aestronglyMeasurable\_pi\_component},
\leanname{aestronglyMeasurable\_pi\_of\_components},
\leanname{memLp\_pi\_component})\\
For $F : \alpha \to \R^d$ and $1 \le p \le \infty$:
\begin{enumerate}
\item $\eLpNormText{F}_{L^p} \le \sum_i \eLpNormText{F_i}_{L^p}$;
\item $\eLpNormText{F_i}_{L^p} \le \eLpNormText{F}_{L^p}$;
\item $F$ is a.e.\ strongly measurable iff each component is;
\item $F \in L^p$ iff each component is.
\end{enumerate}
\end{theorem}

\begin{proof}
For the first item, the pointwise vector-norm bound on $\R^d$ reads
$\|F(x)\| \le \sum_i \|F(x)_i\|$ (\leanname{pi\_norm\_le\_sum\_norms}); applying
$\eLpNormText{\cdot}_{L^p}$ and using monotonicity of the eLpNorm under
pointwise inequalities of nonnegative functions
(\leanname{eLpNorm\_mono\_real}) bounds $\eLpNormText{F}_{L^p}$ above by
$\eLpNormText{x \mapsto \sum_i \|F(x)_i\|}_{L^p}$. Iterating the Minkowski
triangle inequality \leanname{eLpNorm\_add\_le} for $1 \le p \le \infty$ over
the finite index set bounds this in turn by
$\sum_i \eLpNormText{x \mapsto \|F(x)_i\|}_{L^p}$, and then
$\eLpNormText{x \mapsto \|F(x)_i\|}_{L^p} = \eLpNormText{F_i}_{L^p}$ since the
eLpNorm of $\|h\|$ equals that of $h$ (\leanname{eLpNorm\_norm}).

For the single-component bound $\eLpNormText{F_i}_{L^p} \le
\eLpNormText{F}_{L^p}$, the pointwise inequality $\|F(x)_i\| \le \|F(x)\|$
holds for any vector in $\R^d$ (the $i$-th component is dominated by the
sup norm, hence by any equivalent norm), and again $\eLpNormText{\cdot}_{L^p}$
is monotone. Strong measurability componentwise follows by composition with
the continuous projection $\R^d \to \R$ (\leanname{continuous\_apply}); the
converse direction assembles the components into a measurable function via
the product topology on $\R^d$. Membership in $L^p$ in either direction now
follows from the two norm comparisons together with strong measurability.
\end{proof}

\begin{theorem}[\leanname{tendsto\_eLpNorm\_pi\_component}]
Vector $L^p$ convergence implies componentwise $L^p$ convergence.
\end{theorem}

\begin{proof}
Component norms are bounded by vector norms.
\end{proof}

\begin{theorem}[\leanname{exists\_pi\_limit\_of\_cauchy\_eLpNorm}]
A Cauchy sequence $G_n : \alpha \to \R^d$ in $L^p$ admits a limit
$G_{\mathrm{ext}} \in L^p(\alpha; \R^d)$ with both vector and componentwise
$L^p$-convergence.
\end{theorem}

\begin{proof}
For each component, the projected sequence is scalar-valued $L^p$-Cauchy
(by the previous comparisons). Apply Theorem~\ref{thm:scalar-cauchy} to each
component, assemble the resulting limits into the vector function, and
verify the vector $L^p$ convergence by the componentwise sum bound.
\end{proof}

\begin{theorem}[\leanname{exists\_subseq\_tendsto\_ae\_of\_tendsto\_eLpNorm}]
\label{thm:ae-subseq}
Let $f_n, g \in L^2(\mu)$ with $\eLpNormText{f_n - g}_{L^2} \to 0$. There is
a subsequence $f_{n_k}$ with $f_{n_k} \to g$ a.e.
\end{theorem}

\begin{proof}
$L^2$ convergence implies convergence in measure
(\leanname{tendstoInMeasure\_of\_tendsto\_eLpNorm}), and convergence in measure
yields an a.e.\ convergent subsequence
(\leanname{TendstoInMeasure.exists\_seq\_tendsto\_ae}).
\end{proof}

\subsection{Finite covering}

\begin{theorem}[\leanname{exists\_finite\_inner\_ball\_cover}]
Let $0 < r$, $0 < \rho$, and $r + 2\rho < R$. Then there is a finite set
$T \subseteq \overline{\Ball{r}{x_0}}$ such that
\[
\overline{\Ball{r}{x_0}} \subseteq \bigcup_{c \in T} \Ball{\rho}{c},
\quad
\overline{\Ball{2\rho}{c}} \subseteq \Ball{R}{x_0} \text{ for all } c \in T.
\]
\end{theorem}

\begin{proof}
The closed ball $\overline{\Ball{r}{x_0}}$ is compact in the finite-dimensional
Euclidean space $E$ (\leanname{isCompact\_closedBall}), and the family
$\{\Ball{\rho}{c}: c \in \overline{\Ball{r}{x_0}}\}$ is an open cover of
$\overline{\Ball{r}{x_0}}$ since each centre $c$ is contained in its own ball
of radius $\rho > 0$ (\leanname{Metric.ball\_mem\_nhds}). By compactness
(\leanname{IsCompact.elim\_nhds\_subcover}) we extract a finite subset
$T \subseteq \overline{\Ball{r}{x_0}}$ for which the union of the
$\Ball{\rho}{c}$ over $c \in T$ still covers $\overline{\Ball{r}{x_0}}$.

It remains to verify the buffered containment
$\overline{\Ball{2\rho}{c}} \subseteq \Ball{R}{x_0}$ for each $c \in T$. Since
$c$ lies in $\overline{\Ball{r}{x_0}}$ we have $\dist(c,x_0) \le r$, and any
$x \in \overline{\Ball{2\rho}{c}}$ satisfies $\dist(x,c) \le 2\rho$. By the
triangle inequality
\[
  \dist(x, x_0) \le \dist(x, c) + \dist(c, x_0) \le 2\rho + r < R,
\]
where the strict inequality is exactly the buffer hypothesis $r + 2\rho < R$.
Hence $x \in \Ball{R}{x_0}$, completing the proof.
\end{proof}

\begin{theorem}[\leanname{exists\_finite\_inner\_ball\_cover\_with\_card}]
\label{thm:finite-cover-card}
Under the same hypotheses, the cover can be chosen with cardinality bounded
by $\lceil(4r/\rho + 1)^d\rceil$.
\end{theorem}

\begin{proof}
First obtain a fine cover by $\rho/4$-balls from the previous theorem. Apply
the maximality argument
\leanname{exists\_maximal\_separated\_subfinset} to extract a $(\rho/2)$-separated
subset $T \subseteq S$. The disjoint $\rho/4$-balls around the centers of $T$
all lie inside $\Ball{r + \rho/4}{x_0}$. By volume comparison,
\[
|T|\,(\rho/4)^d\, \omega_d \le (r + \rho/4)^d \omega_d,
\quad\text{i.e.}\quad
|T| \le (4r/\rho + 1)^d.
\]
Take ceilings.
\end{proof}

\begin{theorem}[\leanname{exists\_halfBall\_cover\_by\_eighth\_balls}]
There is a finite cover of $\overline{\Ball{1/2}{0}}$ by balls of radius
$1/8$, of cardinality at most $17^d$, such that the doubled balls of radius
$1/4$ each lie inside $\Bz{1}$.
\end{theorem}

\begin{proof}
Apply Theorem~\ref{thm:finite-cover-card} with $r = 1/2$, $\rho = 1/8$,
$R = 1$. The cardinality bound becomes $\lceil (4\cdot(1/2)/(1/8) + 1)^d
\rceil = \lceil 17^d\rceil = 17^d$.
\end{proof}

\begin{theorem}[\leanname{ae\_on\_set\_of\_ae\_on\_finite\_cover}]
If $S \subseteq \bigcup_{i \in T} U_i$ and the property $P(x)$ holds a.e.\
on each $U_i$, then $P$ holds a.e.\ on $S$.
\end{theorem}

\begin{proof}
The a.e.\ statement on a finite union is equivalent to the conjunction of
a.e.\ statements on each member, then restrict to $S$.
\end{proof}

\begin{theorem}[\leanname{lintegralOn\_le\_sum\_lintegralOn\_of\_finite\_cover}]
For a non-negative measurable $f$, $\int_S f \le \sum_{i \in T} \int_{U_i} f$
whenever $S \subseteq \bigcup_{i \in T} U_i$.
\end{theorem}

\begin{proof}
Bound from above by the integral over the union, then apply countable
sub-additivity (here finite) of the integral.
\end{proof}

\subsection{Foundations and shared prelude}

\begin{remark}[\leanname{Common}, \leanname{Foundations}]
The files \leanname{DeGiorgi.Common} and \leanname{DeGiorgi.Foundations}
serve as the shared prelude for the development. They open the namespaces
\texttt{MeasureTheory}, \texttt{Metric}, \texttt{Set}, and \texttt{Filter},
together with the scoped notations for $\R_{\ge 0}^\infty$, $\R_{\ge 0}$,
$\mathrm{Topology}$, and \texttt{InnerProductSpace}, so that subsequent
files can keep their imports minimal. No mathematical content beyond
namespace setup is introduced.
\end{remark}


%% ========================================================================
\section{Elliptic Coefficients and Weak Formulation}\label{sec:weakform}
%% ========================================================================

% Section 5: Elliptic Coefficients and Weak Formulation

Throughout this section we fix a dimension $d \in \N$ with $d \neq 0$ and work in the ambient Euclidean space $E := \R^d$ (denoted \leanname{AmbientSpace} in the formalization, where it is $\mathrm{EuclideanSpace}\,\R\,(\mathrm{Fin}\,d)$). For a $d \times d$ real matrix $M$ and $\xi \in E$, the matrix-vector action is denoted $M\xi$ (\leanname{matMulE}).

\subsection{Elliptic coefficient structure}

\begin{definition}[\leanname{EllipticCoeff}]
Let $\Omega \subseteq E$. A (nonsymmetric) \emph{elliptic coefficient field} on $\Omega$ consists of:
\begin{itemize}
  \item a matrix-valued field $a : E \to \R^{d \times d}$ with each entry $x \mapsto a(x)_{ij}$ measurable;
  \item a lower ellipticity constant $\lam > 0$ (\leanname{lam});
  \item an upper ellipticity constant $\Lambda$ with $\lam \leq \Lambda$;
  \item a coercivity statement: for almost every $x \in \Omega$,
  \[
    \lam \,\|\xi\|^{2} \leq \langle \xi, a(x)\xi\rangle \quad \text{for all } \xi \in E;
  \]
  \item an inverse coercivity statement: for almost every $x \in \Omega$,
  \[
    \Lambda^{-1}\,\|\xi\|^{2} \leq \langle \xi, a(x)^{-1}\xi\rangle \quad \text{for all } \xi \in E.
  \]
\end{itemize}
\end{definition}

\begin{definition}[\leanname{ellipticityRatio}, \leanname{NormalizedEllipticCoeff}]
The \emph{ellipticity ratio} of $A$ is $\mathrm{ER}(A) := \Lambda/\lam$. A coefficient field is \emph{normalized} if $\lam = 1$; in that case $\mathrm{ER}(A) = \Lambda$.
\end{definition}

\begin{lemma}[\leanname{EllipticCoeff.lam\_nonneg}, $\Lambda$\texttt{\_pos}, \leanname{ellipticityRatio\_pos}, \leanname{one\_le\_ellipticityRatio}]
For any $A$ as above one has $\lam \geq 0$, $\Lambda > 0$, $\mathrm{ER}(A) > 0$, and $1 \leq \mathrm{ER}(A)$.
\end{lemma}

\begin{proof}
Positivity of $\lam$ gives $\lam \geq 0$. From $0 < \lam \leq \Lambda$ we get $\Lambda > 0$ and $\Lambda \geq 0$. The ratio $\Lambda/\lam$ is then a quotient of positive reals, hence positive; from $\lam \leq \Lambda$ and $\lam > 0$ we get $1 \leq \Lambda/\lam$.
\end{proof}

\begin{lemma}[\leanname{ae\_coercive\_nonneg}, \leanname{ae\_coercive\_inv\_nonneg}]
For a.e.\ $x \in \Omega$ and every $\xi \in E$, both $\langle \xi, a(x)\xi\rangle$ and $\langle \xi, a(x)^{-1}\xi\rangle$ are nonnegative.
\end{lemma}

\begin{proof}
Apply coercivity (resp.\ inverse coercivity) and observe that $\lam\|\xi\|^2 \geq 0$ and $\Lambda^{-1}\|\xi\|^2 \geq 0$.
\end{proof}

\begin{lemma}[\leanname{det\_ne\_zero\_of\_coercive}, \leanname{inv\_matMulE\_matMulE}]
At any point $x$ where the lower coercivity bound holds, $\det a(x) \neq 0$ and $a(x)^{-1}\bigl(a(x)\xi\bigr) = \xi$ for all $\xi \in E$.
\end{lemma}

\begin{proof}
If $\det a(x) = 0$, choose a nonzero $v$ with $a(x) v = 0$. Then $\lam\|v\|^2 \leq \langle v, a(x)v\rangle = 0$, forcing $\|v\| = 0$ since $\lam > 0$, a contradiction. Once $a(x)$ is invertible, the identity $a(x)^{-1}a(x) = I$ gives the second claim.
\end{proof}

\begin{lemma}[\leanname{EllipticCoeff.mulVec\_sq\_le}: mixed quadratic bound]
For a.e.\ $x \in \Omega$ and all $\xi \in E$,
\[
  \|a(x)\xi\|^{2} \leq \Lambda \,\langle \xi, a(x)\xi\rangle.
\]
\end{lemma}

\begin{proof}
Set $\eta := a(x)\xi$. Inverse coercivity applied to $\eta$ gives $\Lambda^{-1}\|\eta\|^{2} \leq \langle \eta, a(x)^{-1}\eta\rangle$, hence $\|\eta\|^{2} \leq \Lambda \langle \eta, a(x)^{-1}\eta\rangle$. Since $a(x)^{-1}\eta = \xi$ by the previous lemma, the right-hand side equals $\Lambda\langle \eta, \xi\rangle = \Lambda\langle \xi, a(x)\xi\rangle$.
\end{proof}

\begin{lemma}[\leanname{EllipticCoeff.quadratic\_upper}]
For a.e.\ $x \in \Omega$ and all $\xi \in E$,
\[
  \langle \xi, a(x)\xi\rangle \leq \Lambda\,\|\xi\|^{2}.
\]
\end{lemma}

\begin{proof}
Let $q := \langle \xi, a(x)\xi\rangle$, $m := \|a(x)\xi\|$, $n := \|\xi\|$. Cauchy--Schwarz gives $q \leq n m$, and the previous lemma gives $m^{2} \leq \Lambda q$. If $q = 0$ the bound is trivial. Otherwise multiply $q \leq nm$ by $q$ to obtain $q^{2} \leq n^{2} m^{2} \leq \Lambda n^{2} q$, and divide by $q > 0$.
\end{proof}

\begin{lemma}[\leanname{EllipticCoeff.mixed\_bound}]
For a.e.\ $x \in \Omega$ and all $\xi, \zeta \in E$,
\[
  |\langle \zeta, a(x)\xi\rangle| \leq \Lambda\,\|\zeta\|\,\|\xi\|.
\]
\end{lemma}

\begin{proof}
Cauchy--Schwarz gives $|\langle \zeta, a(x)\xi\rangle| \leq \|\zeta\|\,\|a(x)\xi\|$, so it suffices to show $\|a(x)\xi\| \leq \Lambda\|\xi\|$. Squaring, $\|a(x)\xi\|^{2} \leq \Lambda \langle \xi, a(x)\xi\rangle \leq \Lambda^{2}\|\xi\|^{2}$ by the previous two lemmas, and taking square roots (using nonnegativity of both sides) gives the desired estimate.
\end{proof}

\subsection{Pointwise bilinear form on Sobolev witnesses}

Throughout, $u, v \in W^{1,2}(\Omega)$ are equipped with weak-gradient \emph{witnesses} $h_u, h_v$ in the sense of Section 4. We write $\nabla u$ for $h_u.\mathrm{weakGrad}$.

\begin{definition}[\leanname{bilinFormIntegrandOfCoeff}, \leanname{bilinFormOfCoeff}]
The pointwise bilinear-form integrand and the bilinear form are
\[
  b_A(u,v)(x) := \langle a(x)\,\nabla u(x), \nabla v(x)\rangle, \qquad
  \mathfrak{a}_A(u,v) := \int_{\Omega} b_A(u,v)(x)\,dx.
\]
\end{definition}

\begin{lemma}[\leanname{aestronglyMeasurable\_bilinFormIntegrandOfCoeff}, \leanname{integrable\_bilinFormIntegrandOfCoeff}]
The integrand $b_A(u,v)$ is a.e.\ strongly measurable and integrable on $\Omega$.
\end{lemma}

\begin{proof}
Measurability follows from the componentwise measurability of $a$ and the $\Lp{2}$-membership of the components of $\nabla u, \nabla v$. Integrability follows from the pointwise mixed bound $|b_A(u,v)| \leq \Lambda\,\|\nabla u\|\,\|\nabla v\|$ and the Cauchy--Schwarz product $\|\nabla u\|\,\|\nabla v\| \in \Lp{1}$.
\end{proof}

\begin{lemma}[\leanname{bilinForm\_bound}: continuity]
\[
  |\mathfrak{a}_A(u,v)| \leq \Lambda
  \Bigl(\int_{\Omega} \|\nabla u\|^{2}\,dx\Bigr)^{1/2}
  \Bigl(\int_{\Omega} \|\nabla v\|^{2}\,dx\Bigr)^{1/2}.
\]
\end{lemma}

\begin{proof}
By the triangle inequality for integrals and the pointwise bound $|b_A(u,v)| \leq \Lambda\,\|\nabla u\|\,\|\nabla v\|$,
\[
  |\mathfrak{a}_A(u,v)| \leq \Lambda \int_{\Omega} \|\nabla u\|\,\|\nabla v\|\,dx.
\]
H\"older's inequality with conjugate exponents $(2,2)$ bounds the right-hand side by $\Lambda \,\|\nabla u\|_{\Lp{2}}\,\|\nabla v\|_{\Lp{2}}$.
\end{proof}

\begin{lemma}[\leanname{bilinForm\_coercive}]
\[
  \lam \int_{\Omega} \|\nabla u\|^{2}\,dx \leq \mathfrak{a}_A(u,u).
\]
\end{lemma}

\begin{proof}
Coercivity gives $\lam\|\nabla u(x)\|^{2} \leq \langle \nabla u(x), a(x)\nabla u(x)\rangle = b_A(u,u)(x)$ for a.e.\ $x$. Integrate.
\end{proof}

\begin{lemma}[\leanname{bilinFormOfCoeff\_eq\_left}, \leanname{bilinFormOfCoeff\_eq\_right}]
On an open set $\Omega$, $\mathfrak{a}_A(u,v)$ is independent of the chosen $W^{1,2}$ witness for $u$ (resp.\ $v$).
\end{lemma}

\begin{proof}
Two witnesses for the same function on an open set have a.e.-equal weak gradients. The integrand depends only on these gradients, so the integrals agree.
\end{proof}

\begin{lemma}
(\leanname{bilinFormOfCoeff\_add\_left},
\leanname{bilinFormOfCoeff\_smul\_left},
\leanname{bilinFormOfCoeff\_add\_right},
\leanname{bilinFormOfCoeff\_smul\_right})\\
The bilinear form $\mathfrak{a}_A$ is bilinear: additive and homogeneous in both slots.
\end{lemma}

\begin{proof}
The weak gradient of $u_1 + u_2$ is the sum of the weak gradients (by witness construction), and matrix-vector multiplication and the inner product are linear, so the integrand splits. Linearity then transfers to the integral by integrability. The argument for the right slot and for scalar multiplication is identical.
\end{proof}

\subsection{Divergence-form right-hand sides}

\begin{definition}[\leanname{divergenceRHSIntegrandOfField}, \leanname{divergenceRHSOfField}]
For $F : E \to E$ in $\Lp{2}(\Omega; E)$ and a witness $h_v$ for $v \in W^{1,2}(\Omega)$, set
\[
  \ell_F(v)(x) := \langle F(x), \nabla v(x)\rangle, \qquad
  L_F(v) := -\int_{\Omega} \ell_F(v)(x)\,dx.
\]
\end{definition}

\begin{lemma}[\leanname{aestronglyMeasurable\_divergenceRHSIntegrandOfField}, \leanname{integrable\_divergenceRHSIntegrandOfField}]
$\ell_F(v)$ is a.e.\ strongly measurable and integrable on $\Omega$.
\end{lemma}

\begin{proof}
Pointwise, $|\ell_F(v)| \leq \|F\|\,\|\nabla v\|$, and the right-hand side belongs to $\Lp{1}$ by H\"older.
\end{proof}

\begin{lemma}[\leanname{divergenceRHSOfField\_bound}]
\[
  |L_F(v)| \leq \Bigl(\int_{\Omega}\|F\|^{2}\,dx\Bigr)^{1/2}
  \Bigl(\int_{\Omega}\|\nabla v\|^{2}\,dx\Bigr)^{1/2}.
\]
\end{lemma}

\begin{proof}
Bound $|L_F(v)|$ by the integral of $|\ell_F(v)|$, then by the integral of $\|F\|\,\|\nabla v\|$, and apply H\"older.
\end{proof}

\begin{lemma}[\leanname{divergenceRHSOfField\_add}, \leanname{divergenceRHSOfField\_smul}]
$L_F$ is additive and homogeneous in $v$.
\end{lemma}

\begin{proof}
The weak gradient is linear in $v$ (witness construction), and the inner product distributes; integrate.
\end{proof}

\begin{definition}[\leanname{weakProblemRHSOfField}, \leanname{weakProblemRHSOfFieldAndDatum}]
The \emph{raw-function} divergence RHS is $v \mapsto L_F(v)$ when $v \in H^{1}_{0}(\Omega)$ (using a chosen witness) and $0$ otherwise. The \emph{shifted} version with boundary datum $u_0$ is $v \mapsto L_F(v) - \mathfrak{a}_A(u_0, v)$.
\end{definition}

\begin{lemma}[\leanname{weakProblemRHSOfField\_eq\_of\_memH01}, \leanname{weakProblemRHSOfFieldAndDatum\_eq\_of\_memH01}]
For $v \in H^{1}_{0}(\Omega)$, the raw-function (resp.\ shifted) RHS agrees with $L_F(v)$ (resp.\ $L_F(v) - \mathfrak{a}_A(u_0,v)$) on every witness.
\end{lemma}

\begin{proof}
Two witnesses for $v$ on the open set $\Omega$ have a.e.-equal weak gradients, so $L_F$ and $\mathfrak{a}_A$ depend only on $v$ itself (Lemmas above).
\end{proof}

\begin{lemma}[\leanname{weakProblemRHSOfField\_bound}, \leanname{weakProblemRHSOfFieldAndDatum\_bound}]
For $v \in H^{1}_{0}(\Omega)$ with witness $h_v$,
\begin{align*}
  |L_F(v)| &\leq \|F\|_{\Lp{2}} \,\|\nabla v\|_{\Lp{2}}, \\
  |L_F(v) - \mathfrak{a}_A(u_0, v)| &\leq \bigl(\|F\|_{\Lp{2}} + \Lambda\,\|\nabla u_0\|_{\Lp{2}}\bigr)\,\|\nabla v\|_{\Lp{2}}.
\end{align*}
\end{lemma}

\begin{proof}
Combine the previous bounds using the triangle inequality and the bilinear-form continuity bound $|\mathfrak{a}_A(u_0,v)| \leq \Lambda\,\|\nabla u_0\|_{\Lp{2}}\|\nabla v\|_{\Lp{2}}$.
\end{proof}

\subsection{Smooth test functions and the L\texorpdfstring{$^2$}{2} coefficient operator}

\begin{definition}[\leanname{IsSmoothTestOn}]
A function $u : E \to \R$ is a \emph{smooth test on $\Omega$} if $u \in C^{\infty}(E)$, $u$ has compact support, and $\supp u \subseteq \Omega$.
\end{definition}

\begin{lemma}[\leanname{IsSmoothTestOn.zero}, \leanname{add}, \leanname{smul}, \leanname{sub}]
The smooth tests on $\Omega$ form a real vector space (closed under sums, scalar multiples, and differences).
\end{lemma}

\begin{proof}
Smoothness and compact support are stable under sums and scalar multiples; the support of a sum is contained in the union of supports, hence in $\Omega$.
\end{proof}

\begin{definition}[\leanname{smoothGradField}, \leanname{smoothTestWitness}]
For smooth $u$, the \emph{classical gradient field} is
\[
  (\nabla u)(x) := \bigl((\partial_{i} u)(x)\bigr)_{i=1}^{d} \in E.
\]
A smooth test on $\Omega$ carries a canonical $W^{1,2}(\Omega)$-witness with this classical gradient.
\end{definition}

\begin{lemma}[\leanname{smoothTest\_memH01}]
A smooth test on $\Omega$ belongs to $H^{1}_{0}(\Omega)$.
\end{lemma}

\begin{proof}
Approximate by its constant sequence; smoothness and compactly supported support in $\Omega$ provide the required convergence in $W^{1,2}$.
\end{proof}

\begin{definition}[\leanname{smoothFunToLp}, \leanname{smoothGradToLp}, \leanname{gradLpOfWitness}]
A smooth test $u$ defines:
\[
  [u] \in \Lp{2}(\Omega; \R), \qquad [\nabla u] \in \Lp{2}(\Omega; E),
\]
the latter being the $\Lp{2}$ class of any witness's weak gradient.
\end{definition}

\begin{lemma}[\leanname{smoothFunToLp\_zero}, \leanname{\_add}, \leanname{\_smul}, \leanname{\_sub}, and analogues for \leanname{smoothGradToLp}]
Both maps $u \mapsto [u]$ and $u \mapsto [\nabla u]$ are real-linear.
\end{lemma}

\begin{proof}
Linearity of the $\Lp{2}$ class follows from $\mathrm{toLp}$-linearity in Mathlib; for the gradient version one further uses linearity of the classical gradient.
\end{proof}

\begin{lemma}[\leanname{norm\_smoothFunToLp\_eq}, \leanname{norm\_smoothGradToLp\_eq}, \leanname{norm\_gradLpOfWitness\_eq}]
The norms in $\Lp{2}$ are
\[
  \bigl\|[u]\bigr\| = \Bigl(\int_{\Omega} |u|^{2}\,dx\Bigr)^{1/2}, \qquad
  \bigl\|[\nabla u]\bigr\| = \Bigl(\int_{\Omega} \|\nabla u\|^{2}\,dx\Bigr)^{1/2}.
\]
\end{lemma}

\begin{proof}
Direct from the definition of the $\Lp{2}$ norm via $\mathrm{toLp}$ in Mathlib.
\end{proof}

\subsection{The L\texorpdfstring{$^2$}{2} coefficient operator and the closed bilinear form}

\begin{definition}[\leanname{smoothGradSubmodule}]
Let $\Omega$ be open. Define
\[
  S(\Omega) := \bigl\{ [\nabla u] : u \text{ smooth test on } \Omega \bigr\} \subseteq \Lp{2}(\Omega; E),
\]
a real linear subspace, and let $H(\Omega)$ be its closure.
\end{definition}

\begin{lemma}[\leanname{coeffMulLpL}, \leanname{norm\_coeffMulLpL\_le}]
There is a bounded linear operator $T_A : \Lp{2}(\Omega;E) \to \Lp{2}(\Omega;E)$ given pointwise a.e.\ by $T_A g (x) = a(x)\,g(x)$, satisfying $\|T_A g\|_{\Lp{2}} \leq \Lambda\,\|g\|_{\Lp{2}}$.
\end{lemma}

\begin{proof}
For each component $i$, $|T_A g(x)_i| \leq \|T_A g(x)\| \leq \Lambda\|g(x)\|$ by the mixed bound, and componentwise measurability of $a$ and $g$ make $T_A g$ strongly measurable. Hence each component lies in $\Lp{2}$, so does $T_A g$, and the operator norm bound follows from $\|T_A g(x)\| \leq \Lambda\|g(x)\|$ a.e.\ via $\Lp{p}$ pointwise dominance.
\end{proof}

\begin{definition}[\leanname{coeffBilinSubmodule}]
For any closed real subspace $H \leq \Lp{2}(\Omega;E)$ define the continuous bilinear form
\[
  B_{A,H}(g,h) := \langle T_A g, h\rangle_{\Lp{2}}.
\]
\end{definition}

\begin{lemma}[\leanname{coeffBilinSubmodule\_coercive}]
$B_{A,H}$ is bounded by $\Lambda$ and coercive on $H$ with constant $\lam$:
\[
  |B_{A,H}(g,h)| \leq \Lambda\,\|g\|\,\|h\|, \qquad \lam\,\|g\|^{2} \leq B_{A,H}(g,g).
\]
\end{lemma}

\begin{proof}
The boundedness $|\langle T_A g, h\rangle| \leq \|T_A g\|\,\|h\| \leq \Lambda \|g\|\,\|h\|$ comes from Cauchy--Schwarz and the bound on $T_A$. For coercivity, write
\[
  B_{A,H}(g,g) = \int_{\Omega}\langle a(x) g(x), g(x)\rangle\,dx \geq \int_{\Omega} \lam\,\|g(x)\|^{2}\,dx = \lam\,\|g\|_{\Lp{2}}^{2}
\]
using the a.e.\ coercivity of $A$.
\end{proof}

\begin{lemma}[\leanname{tendsto\_eLpNorm\_vector\_of\_componentwise}]
Componentwise $\Lp{2}$-convergence in each coordinate of $\Fin d$ implies vector $\Lp{2}$-convergence: if for each $i$ the seminorm $\|g_n^{(i)} - g^{(i)}\|_{\Lp{2}} \to 0$, then $\|g_n - g\|_{\Lp{2}} \to 0$.
\end{lemma}

\begin{proof}
Pointwise, $\|g_n(x) - g(x)\|^{2} = \sum_i (g_n(x)_i - g(x)_i)^{2}$, hence $\|g_n(x) - g(x)\| \leq \sum_i |g_n(x)_i - g(x)_i|$. The $\Lp{2}$ norm of the right-hand side is bounded by $\sum_i \|g_n^{(i)} - g^{(i)}\|_{\Lp{2}}$, which tends to zero.
\end{proof}

\begin{lemma}[\leanname{coeffMulLpL\_inner\_gradLpOfWitness\_eq\_bilinFormOfCoeff}]
For witnesses $h_u, h_v$ for $u, v \in W^{1,2}(\Omega)$,
\[
  \langle T_A [\nabla u], [\nabla v]\rangle_{\Lp{2}} = \mathfrak{a}_A(u, v).
\]
\end{lemma}

\begin{proof}
Both sides equal $\int_\Omega \langle a(x)\,\nabla u(x), \nabla v(x)\rangle\,dx$; the left-hand side after using the a.e.\ identification $T_A[\nabla u] = a(\cdot)\nabla u(\cdot)$.
\end{proof}

\subsection{Weak problems and weak solutions}

\begin{definition}[\leanname{WeakProblem}]
A \emph{weak elliptic Dirichlet problem} consists of an open bounded set $\Omega \subseteq E$, an elliptic coefficient field $A$ on $\Omega$, and a functional $\mathrm{rhs} : (E \to \R) \to \R$.
\end{definition}

\begin{definition}[\leanname{IsWeakSolution}]
A function $u : E \to \R$ is a \emph{weak solution} of $P = (\Omega, A, \mathrm{rhs})$ if $u \in H^{1}_{0}(\Omega)$ and, for every witness $h_u$ of $u$ and every $v \in H^{1}_{0}(\Omega)$ with witness $h_v$,
\[
  \mathfrak{a}_A(u, v) = \mathrm{rhs}(v).
\]
\end{definition}

\begin{definition}[\leanname{IsSubsolution}, \leanname{IsSupersolution}, \leanname{IsSolution}]
For an elliptic field $A$ on $\Omega$, $u \in W^{1,2}(\Omega)$ is a (weak) \emph{subsolution} of $-\dv(a\nabla u) \leq 0$ if for all $\varphi \in H^{1}_{0}(\Omega)$ with $\varphi \geq 0$, $\mathfrak{a}_A(u, \varphi) \leq 0$. The notion of \emph{supersolution} reverses the inequality. A \emph{solution} is both a sub- and supersolution.
\end{definition}

\begin{definition}[\leanname{IsHomogeneousWeakSolution}]
$u \in W^{1,2}(\Omega)$ is a \emph{homogeneous weak solution} of $-\dv(a\nabla u) = 0$ on $\Omega$ if, for every $W^{1,2}$-witness $h_u$ of $u$ and every $\varphi \in H^{1}_{0}(\Omega)$ with $W^{1,2}$-witness $h_\varphi$, one has $\mathfrak{a}_A(u,\varphi) = 0$.
\end{definition}

\begin{lemma}[\leanname{isHomogeneousWeakSolution\_isSubsolution}, \leanname{isHomogeneousWeakSolution\_isSupersolution}, \leanname{isHomogeneousWeakSolution\_isSolution}]
A homogeneous weak solution is both a subsolution and a supersolution (hence a solution).
\end{lemma}

\begin{proof}
For any nonnegative test $\varphi$, $\mathfrak{a}_A(u,\varphi) = 0$, so both inequalities hold trivially.
\end{proof}

Note that \leanname{IsWeakSolution} is the existential predicate used by the Lax--Milgram existence theorem (it carries an explicit RHS functional and forces $u \in H^{1}_{0}$). \leanname{IsHomogeneousWeakSolution} is the equality-form local interface used by the Harnack and H\"older statements: it requires $u \in W^{1,2}(\Omega)$ but does \emph{not} require zero boundary values, and uses test functions in $H^{1}_{0}(\Omega)$.

\subsection{Existence via Lax--Milgram}

\begin{theorem}[\leanname{weakProblem\_exists}]
\label{thm:weakProblem-exists}
Let $d \geq 2$. Let $\Omega \subseteq E$ be open and bounded, $A$ an elliptic coefficient field on $\Omega$, and $\mathrm{rhs}$ a functional on $H^{1}_{0}(\Omega)$ that is real-linear and bounded:
\[
  |\mathrm{rhs}(v)| \leq C_{\mathrm{rhs}}\,\|\nabla v\|_{\Lp{2}}.
\]
Then there exists $u \in H^{1}_{0}(\Omega)$ that is a weak solution of $(\Omega, A, \mathrm{rhs})$.
\end{theorem}

\begin{proof}
Let $\mu := \mathrm{vol}\!\restriction\!\Omega$ and let $S = S(\Omega) \leq \Lp{2}(\Omega; E)$ be the smooth-gradient submodule with closure $H = H(\Omega)$. By the Poincar\'e inequality on bounded domains in dimension $d \geq 2$, there is a constant $C_P > 0$ with
\[
  \|u\|_{\Lp{2}} \leq C_P\,\|\nabla u\|_{\Lp{2}} \quad \text{for every smooth test } u \text{ on } \Omega.
\]
For each $g \in S$ choose, via Classical choice, a smooth test $u_g$ with $[\nabla u_g] = g$. The Poincar\'e inequality, lifted via $\mathrm{toReal}$, gives $\|[u_g]\|_{\Lp{2}} \leq C_P \|g\|$, and shows that two such choices with the same gradient class differ by a function with vanishing $\Lp{2}$-class.

Define linear maps on $S$ by
\[
  L_0(g) := \mathrm{rhs}(u_g), \qquad U_0(g) := [u_g] \in \Lp{2}(\Omega; \R).
\]
The hypotheses on $\mathrm{rhs}$ and the Poincar\'e inequality, together with the fact that two representatives with the same gradient class give the same value of $\mathrm{rhs}$ (because their difference has zero gradient class and the bound forces $|\mathrm{rhs}(u_1) - \mathrm{rhs}(u_2)| \leq C_{\mathrm{rhs}} \cdot 0 = 0$), make $L_0$ and $U_0$ well-defined linear maps with bounds
\[
  |L_0(g)| \leq C_{\mathrm{rhs}}\|g\|, \qquad \|U_0(g)\| \leq C_P\|g\|.
\]
Since the inclusion $S \hookrightarrow H$ has dense range (as $H$ is the closure of $S$), $L_0$ and $U_0$ extend uniquely to continuous maps $L : H \to \R$, $U : H \to \Lp{2}(\Omega;\R)$.

By the coercivity of the submodule bilinear form established above ($B_{A,H}$ is bounded and coercive on the Hilbert space $H$, which is complete; see \leanname{coeffBilinSubmodule\_coercive}), the Lax--Milgram theorem produces a unique $g_{\mathrm{sol}} \in H$ with
\[
  B_{A,H}(g_{\mathrm{sol}}, w) = L(w) \quad \text{for all } w \in H.
\]
Set $G_{\mathrm{sol}}$ to be the underlying $\Lp{2}(\Omega;E)$ element of $g_{\mathrm{sol}}$ and let $u_{\mathrm{Lp}} := U(g_{\mathrm{sol}}) \in \Lp{2}(\Omega;\R)$ with representative $u : E \to \R$.

Approximate $g_{\mathrm{sol}}$ by a sequence $G_n \in S$ (closure), with corresponding smooth tests $\psi_n$ such that $[\nabla \psi_n] = G_n$. By continuity of $U$, $[\psi_n] \to u_{\mathrm{Lp}}$ in $\Lp{2}$. By the componentwise convergence lemma, the components of $\nabla \psi_n$ converge to the corresponding components of $G_{\mathrm{sol}}$ in $\Lp{2}$, and $\psi_n \to u$ in $\Lp{2}$. By the standard transition principle for weak partial derivatives under $\Lp{p}$-approximation, $u$ has weak gradient given by $G_{\mathrm{sol}}$, providing a witness $h_u$ with $[\nabla u] = G_{\mathrm{sol}}$. Since $\psi_n$ are smooth tests on $\Omega$, $u \in H^{1}_{0}(\Omega)$.

For any smooth test $\varphi$ on $\Omega$ with corresponding $g_\varphi \in S$,
\[
  \mathfrak{a}_A(u, \varphi) = \langle T_A [\nabla u], [\nabla \varphi]\rangle_{\Lp{2}} = B_{A,H}(g_{\mathrm{sol}}, g_\varphi) = L(g_\varphi) = L_0(g_\varphi) = \mathrm{rhs}(\varphi).
\]
Finally, by approximation: any $v \in H^{1}_{0}(\Omega)$ admits an approximating sequence $\varphi_n$ of smooth tests in the $H^{1}_{0}$-sense; the bilinear form is continuous against gradient seminorms, and $\mathrm{rhs}$ is also bounded against the gradient seminorm. Passing to the limit in $\mathfrak{a}_A(u, \varphi_n) = \mathrm{rhs}(\varphi_n)$ (using independence of the bilinear form on the choice of witness in the right slot) gives $\mathfrak{a}_A(u, v) = \mathrm{rhs}(v)$ for any witness of $v$. This completes the proof.
\end{proof}

\begin{theorem}[\leanname{weakProblem\_exists\_of\_divergenceData}]
Under the hypotheses of Theorem~\ref{thm:weakProblem-exists} and $F \in \Lp{2}(\Omega; E)$, there exists a zero-Dirichlet weak solution of $-\dv(a\nabla u) = -\dv F$, i.e.\ a weak solution for the raw functional $v \mapsto L_F(v)$.
\end{theorem}

\begin{proof}
We invoke the abstract existence theorem
Theorem~\ref{thm:weakProblem-exists} (\leanname{weakProblem\_exists}) with
$\mathrm{rhs}$ chosen to be $\leanname{weakProblemRHSOfField}\,F$, the
linear functional $v \mapsto \int_{\Omega} \langle F, \nabla v\rangle\,dx$
acting on $H^{1}_{0}(\Omega)$. The abstract theorem requires three inputs.

\emph{Additivity.} For $u, v \in H^{1}_{0}(\Omega)$ pick canonical witnesses
through \leanname{MemW1p.someWitness}; the sum witness for $u + v$ is the
componentwise sum (\leanname{MemW1pWitness.add}). Both
$\mathrm{rhs}(u+v)$ and $\mathrm{rhs}(u) + \mathrm{rhs}(v)$ unfold via
\leanname{weakProblemRHSOfField\_eq\_of\_memH01} into integrals against the
weak gradient, so additivity reduces to
\leanname{divergenceRHSOfField\_add} applied to those witnesses.

\emph{Homogeneity.} Identical reasoning with the scalar witness
$\mathrm{MemW1pWitness.smul}$ and \leanname{divergenceRHSOfField\_smul}
gives $\mathrm{rhs}(c\,u) = c\,\mathrm{rhs}(u)$ for every $c \in \R$.

\emph{Boundedness against the gradient seminorm.} The integral
$\bigl|\int_{\Omega} \langle F, \nabla v\rangle\bigr|$ is controlled by
$\eLpNormText{F}_{L^2}\,\eLpNormText{\nabla v}_{L^2}$ via Cauchy--Schwarz; this is
recorded as \leanname{weakProblemRHSOfField\_bound}. We take
$C_F := \eLpNormText{F}_{L^2(\Omega)}$, which is finite by the hypothesis
$F \in L^2(\Omega; E)$.

With these three properties, Theorem~\ref{thm:weakProblem-exists} produces a
zero-Dirichlet weak solution $u \in H^{1}_{0}(\Omega)$ of
$-\dv(a\nabla u) = -\dv F$.
\end{proof}

\begin{theorem}[\leanname{dirichletProblem\_exists\_of\_divergenceData}, \leanname{dirichletProblem\_exists\_of\_divergenceData'}]
Under the same hypotheses, with prescribed boundary datum $u_0 \in W^{1,2}(\Omega)$ and $F \in \Lp{2}(\Omega;E)$, there exists $u \in W^{1,2}(\Omega)$ with $u - u_0 \in H^{1}_{0}(\Omega)$ and
\[
  \mathfrak{a}_A(u, v) = L_F(v) \quad \text{for all } v \in H^{1}_{0}(\Omega).
\]
\end{theorem}

\begin{proof}
Choose a witness $h_{u_0}$ for $u_0$. Apply the previous theorem to the shifted RHS $\mathrm{rhs}(v) = L_F(v) - \mathfrak{a}_A(u_0, v)$ to obtain a zero-Dirichlet weak solution $w$. Set $u := w + u_0$. Then $u - u_0 = w \in H^{1}_{0}(\Omega)$, and bilinearity of $\mathfrak{a}_A$ gives
\[
  \mathfrak{a}_A(u, v) = \mathfrak{a}_A(w, v) + \mathfrak{a}_A(u_0, v) = \mathrm{rhs}(v) + \mathfrak{a}_A(u_0, v) = L_F(v).
\]
\end{proof}

\begin{theorem}[\leanname{weakSolution\_stability}: energy estimate]
Let $A$ be elliptic on $\Omega$, $u \in H^{1}_{0}(\Omega)$ a weak solution against a functional $\mathrm{rhs}$ bounded by $C_F\|\nabla v\|_{\Lp{2}}$. Then
\[
  \|\nabla u\|_{\Lp{2}} \leq \frac{C_F}{\lam}.
\]
\end{theorem}

\begin{proof}
Test against $u$ itself. Coercivity gives $\lam \|\nabla u\|_{\Lp{2}}^{2} \leq \mathfrak{a}_A(u, u) = \mathrm{rhs}(u)$, while boundedness gives $|\mathrm{rhs}(u)| \leq C_F \|\nabla u\|_{\Lp{2}}$. The bilinear form value is nonnegative (by coercivity), hence $\mathrm{rhs}(u) \geq 0$ and equals $|\mathrm{rhs}(u)|$. Combining,
\[
  \lam \|\nabla u\|_{\Lp{2}}^{2} \leq C_F \|\nabla u\|_{\Lp{2}},
\]
and dividing (or using $\nabla u = 0$) yields $\|\nabla u\|_{\Lp{2}} \leq C_F/\lam$.
\end{proof}

\begin{theorem}[\leanname{ae\_eq\_zero\_of\_memH01\_of\_gradLpOfWitness\_eq\_zero}]
For $d \geq 2$, an open bounded $\Omega$, and $u \in H^{1}_{0}(\Omega)$ with vanishing $\Lp{2}$ gradient class, $u = 0$ a.e.\ on $\Omega$.
\end{theorem}

\begin{proof}
Take an approximating sequence of smooth tests $\varphi_n$ from the $H^{1}_{0}$ structure with $\varphi_n \to u$ in $\Lp{2}$ and $\nabla \varphi_n$ converging componentwise to a representative of $0$. By Poincar\'e on the bounded domain, $\|\varphi_n\|_{\Lp{2}} \leq C_P \|\nabla \varphi_n\|_{\Lp{2}} \to 0$, so $[\varphi_n] \to 0$ in $\Lp{2}$. Combining with $[\varphi_n] \to [u]$, we get $[u] = 0$, i.e.\ $u = 0$ a.e.
\end{proof}

\begin{theorem}[\leanname{weakProblem\_unique}]
For $d \geq 2$, weak solutions of a fixed weak problem $P$ are unique up to a.e.\ equality.
\end{theorem}

\begin{proof}
If $u, v$ are two weak solutions, $w := u - v$ lies in $H^{1}_{0}$ and tests against any $\varphi \in H^{1}_{0}$ to zero by linearity. By the stability estimate (with $C_F = 0$), $\|\nabla w\|_{\Lp{2}} = 0$, and the previous theorem gives $w = 0$ a.e.
\end{proof}

\begin{definition}[\leanname{IsDirichletWeakSolutionOfDivergenceData}]
$u$ is a Dirichlet weak solution for boundary datum $u_0$ and field $F$ if $u \in W^{1,2}(\Omega)$, $u - u_0 \in H^{1}_{0}(\Omega)$, and $\mathfrak{a}_A(u, v) = L_F(v)$ for every $v \in H^{1}_{0}(\Omega)$ on every witness.
\end{definition}

\begin{theorem}
(\leanname{isWeakSolution\_subDatum\_of\_dirichletProblem},
\leanname{dirichletProblem\_stability\_of\_divergenceData},
\leanname{dirichletProblem\_unique\_of\_divergenceData},
\leanname{dirichletProblem\_wellPosed\_of\_divergenceData})\\
The shifted unknown $u - u_0$ is a zero-boundary weak solution for the shifted RHS, satisfies the energy estimate
\[
  \|\nabla(u - u_0)\|_{\Lp{2}} \leq \frac{1}{\lam}\Bigl(\|F\|_{\Lp{2}} + \Lambda\,\|\nabla u_0\|_{\Lp{2}}\Bigr),
\]
and Dirichlet weak solutions for fixed $(u_0, F)$ are unique up to a.e.\ equality. Combined with existence, we obtain well-posedness.
\end{theorem}

\begin{proof}
The shifted-solution identity follows by linearity of $\mathfrak{a}_A$: $\mathfrak{a}_A(u - u_0, v) = \mathfrak{a}_A(u, v) - \mathfrak{a}_A(u_0, v) = L_F(v) - \mathfrak{a}_A(u_0, v)$. Apply \leanname{weakSolution\_stability} to the shifted unknown with $C_F = \|F\|_{\Lp{2}} + \Lambda\|\nabla u_0\|_{\Lp{2}}$ to obtain the energy estimate. Uniqueness reduces to uniqueness of zero-Dirichlet solutions for the shifted RHS, hence to \leanname{weakProblem\_unique}.
\end{proof}

\begin{theorem}[\leanname{aHarmonic\_replacement\_exists}, \leanname{aHarmonic\_replacement\_unique}]
For $d \geq 2$, an open bounded $\Omega$, an elliptic field $A$ on $\Omega$, and $u \in W^{1,2}(\Omega)$, there exists a homogeneous weak solution $h$ of $-\dv(a\nabla h) = 0$ with $h - u \in H^{1}_{0}(\Omega)$, and any two such replacements agree a.e.\ on $\Omega$.
\end{theorem}

\begin{proof}
Apply \leanname{dirichletProblem\_exists\_of\_divergenceData} with the
prescribed boundary datum $u_0 := u$ and the trivial divergence field
$F := 0 \in L^{2}(\Omega; E)$. The hypotheses are met because $u \in
W^{1,2}(\Omega)$ by assumption and the constant zero function is in $L^2$
(\leanname{MeasureTheory.MemLp.zero'}). The theorem produces a
$h \in W^{1,2}(\Omega)$ with $h - u \in H^{1}_{0}(\Omega)$ satisfying
\[
  \mathfrak{a}_A(h, v) = L_F(v) = \int_{\Omega}\langle 0, \nabla v\rangle\,dx = 0
  \quad \text{for every } v \in H^{1}_{0}(\Omega),
\]
i.e.\ $h$ is a homogeneous weak solution against any choice of witness for $v$.

To see that $h$ is moreover a \emph{homogeneous} weak solution in the sense of
\leanname{IsHomogeneousWeakSolution} (the local interface used by Harnack and
H\"older), unfold the definition: it requires that for every test
$\varphi \in H^{1}_{0}(\Omega)$ and every witness, the bilinear form vanishes.
This is exactly the displayed identity (with the right-hand side
$\leanname{divergenceRHSOfField}\,0$ unfolding to zero by
\leanname{divergenceRHSIntegrandOfField}).

Uniqueness follows directly from
\leanname{dirichletProblem\_unique\_of\_divergenceData} applied to the same
data $(u_0, F) = (u, 0)$: any two replacements have identical Dirichlet trace
along $u$ and identical zero divergence data, so they coincide a.e.\ on
$\Omega$.
\end{proof}

\subsection{Weighted bilinear-form estimates}

\begin{lemma}[\leanname{bilinFormIntegrand\_mul\_smooth\_eq}]
Let $u, w \in W^{1,2}(\Omega)$ with witnesses $h_u, h_w$, and $\eta \in C^{\infty}(E)$ satisfying $|\eta| \leq C_0$ and $\|\fderiv{\eta}{x}\| \leq C_1$. Then for every $x$,
\begin{multline*}
  b_A(u, \eta\!\cdot\!w)(x) = \eta(x)\,b_A(u,w)(x) + w(x)\,\bigl\langle a(x)\,\nabla u(x), \nabla \eta(x)\bigr\rangle,
\end{multline*}
where $\nabla(\eta\cdot w)$ is realized by the witness produced from \leanname{MemW1pWitness.mul\_smooth\_bounded\_p} as $\eta\,\nabla w + w\,\nabla \eta$.
\end{lemma}

\begin{proof}
The witness for $\eta w$ has weak gradient $\eta(x)\,\nabla w(x) + w(x)\,\nabla \eta(x)$ pointwise. Inserting into $\langle a(x)\nabla u(x), \cdot\rangle$ and using bilinearity of the inner product splits the integrand as claimed. The reformulation as a finite sum over $i$ is just unfolding the matrix-vector product and the $\Lp{2}$ inner product.
\end{proof}

\begin{theorem}[\leanname{bilinForm\_weighted\_cauchySchwarz}]
For $u, v \in W^{1,2}(\Omega)$ with witnesses $h_u, h_v$, a measurable weight $f : E \to \R$, and assuming $f^{2}\,b_A(u,u)$ is integrable, one has
\begin{multline*}
  \Bigl(\int_{\Omega} f(x)\,b_A(u,v)(x)\,dx\Bigr)^{2}
  \leq \Bigl(\int_{\Omega} f(x)^{2}\,b_A(u,u)(x)\,dx\Bigr)\\
  \cdot \Bigl(\Lambda \int_{\Omega} \|\nabla v(x)\|^{2}\,dx\Bigr).
\end{multline*}
\end{theorem}

\begin{proof}
If $f\,b_A(u,v)$ is non-integrable then the integral is defined to be zero and the inequality is trivial (the right-hand side is nonnegative). Otherwise, point\-wise Cauchy--Schwarz on the inner product gives
\[
  |f(x)\,b_A(u,v)(x)| \leq |f(x)|\,\|a(x)\nabla u(x)\|\,\|\nabla v(x)\|.
\]
By the mixed bound
$\|a(x)\nabla u(x)\|^{2} \leq \Lambda\,b_A(u,u)(x)$,
the function
\[
  g(x) := |f(x)|\,\|a(x)\nabla u(x)\|
\]
satisfies
\[
  g(x)^{2} = f(x)^{2}\,\|a(x)\nabla u(x)\|^{2} \leq \Lambda\,f(x)^{2}\,b_A(u,u)(x),
\]
so $g \in \Lp{2}$. H\"older's inequality with conjugate exponents $(2,2)$ applied to $g$ and $h(x) := \|\nabla v(x)\|$ gives
\[
  \int g\,h\,dx \leq \Bigl(\int g^{2}\,dx\Bigr)^{1/2}\Bigl(\int h^{2}\,dx\Bigr)^{1/2},
\]
and squaring (using $g^{2} \leq \Lambda f^{2}\,b_A(u,u)$) yields the claimed inequality.
\end{proof}

\begin{lemma}[\leanname{integrable\_bounded\_mul\_bilinFormIntegrand}]
If $|f(x)| \leq C$ everywhere and $u \in W^{1,2}(\Omega)$, then $f \cdot b_A(u,u) \in \Lp{1}(\Omega)$.
\end{lemma}

\begin{proof}
Pointwise, $|f(x)\,b_A(u,u)(x)| \leq C \,|b_A(u,u)(x)|$, and the right-hand side is integrable since $b_A(u,u)$ is.
\end{proof}

\subsection{Ball scaling}

We now record the affine pullback $z \mapsto x_0 + R z$ from the unit ball $\Bz{1}$ to $\Ball{R}{x_0}$ and its effect on weak (sub-/super-)solutions. Throughout, $R > 0$.

\begin{definition}[\leanname{rescaleToUnitBall}, \leanname{transportFromUnitBall}]
For $u : E \to \R$, write
\[
  (T^{*}_{x_0,R} u)(z) := u(x_0 + R z), \qquad
  (T_{x_0,R} u)(x) := u\bigl(R^{-1}(x - x_0)\bigr),
\]
the pullback to the unit ball and the pushforward from the unit ball, respectively.
\end{definition}

\begin{lemma}[\leanname{rescaleToUnitBall\_transportFromUnitBall}]
$T^{*}_{x_0,R} \circ T_{x_0,R} = \mathrm{id}$.
\end{lemma}

\begin{proof}
Direct unfolding: $u\bigl(R^{-1}((x_0 + Rz) - x_0)\bigr) = u(z)$.
\end{proof}

\begin{definition}[\leanname{rescaleCoeffToUnitBall}, \leanname{rescaleNormalizedCoeffToUnitBall}]
Pull back an elliptic field $A$ on $\Ball{R}{x_0}$ to the unit ball by $\widetilde A(z) := (A.a)(x_0 + R z)$, with the same constants $\lam, \Lambda$. The result is an elliptic coefficient field on $\Bz{1}$. The construction preserves normalization $\lam = 1$.
\end{definition}

\begin{lemma}[\leanname{MemW1pWitness.transportFromUnitBall}]
A $W^{1,p}$-witness for $u$ on $\Bz{1}$ transports to a $W^{1,p}$-witness for $T_{x_0,R} u$ on $\Ball{R}{x_0}$, with weak gradient $x \mapsto R^{-1}\,(\nabla u)\bigl(R^{-1}(x - x_0)\bigr)$.
\end{lemma}

\begin{proof}
The Lp norms transform by an explicit power of $R$ depending on the dimension; weak partial derivative compatibility is preserved by the chain rule for the affine map $R^{-1}(\cdot - x_0)$.
\end{proof}

\begin{theorem}[\leanname{rescaleToUnitBall\_isSubsolution}, \leanname{\_isSupersolution}, \leanname{\_isSolution}, \leanname{\_isHomogeneousWeakSolution}]
If $u$ is a weak subsolution (resp.\ supersolution, solution, homogeneous weak solution) of $A$ on $\Ball{R}{x_0}$, then $T^{*}_{x_0,R} u$ is the corresponding object for $\widetilde A$ on $\Bz{1}$.
\end{theorem}

\begin{proof}
The key computation is
\[
  \mathfrak{a}_{\widetilde A}\bigl(T^{*}_{x_0,R} u, \varphi\bigr) = \bigl(R^{-d} \cdot R^{2}\bigr)\,\mathfrak{a}_{A}\bigl(u, T_{x_0,R} \varphi\bigr),
\]
proven by changing variables in the integral defining $\mathfrak{a}$, using the Jacobian factor $R^{-d}$ and the chain-rule factor $R^{-1}$ in each gradient (giving an $R^{2}$ from the two slots). Since the multiplicative constant $R^{2-d} > 0$ has fixed sign, the inequality $\mathfrak{a}_{\widetilde A}(\cdot, \varphi) \leq 0$ on the unit ball is equivalent to $\mathfrak{a}_A(\cdot, T_{x_0,R}\varphi) \leq 0$ on $\Ball{R}{x_0}$, and similarly for the reverse inequality, equality, and the homogeneous case. Nonnegative test functions $\varphi$ on $\Bz{1}$ pull back to nonnegative test functions $T_{x_0,R}\varphi$ on $\Ball{R}{x_0}$ (and vice versa), so the sign condition is preserved.
\end{proof}

\subsection{Localization helpers}

\begin{definition}[\leanname{EllipticCoeff.restrict}, \leanname{NormalizedEllipticCoeff.restrict}]
Given $\Omega' \subseteq \Omega$ and $A$ elliptic on $\Omega$, the restriction $A\!\restriction\!\Omega'$ keeps the same matrix and constants and shrinks the a.e.\ coercivity domain to $\Omega'$. Normalization is preserved.
\end{definition}

\begin{lemma}[\leanname{essSup\_rescale\_halfBall}, \leanname{essInf\_rescale\_halfBall}]
For $R > 0$,
\[
  \esssup_{z \in \Bz{1/2}} u(x_0 + R z) = \esssup_{x \in \Ball{R/2}{x_0}} u(x),
\]
and similarly for $\essinf$.
\end{lemma}

\begin{proof}
The affine map $z \mapsto x_0 + R z$ is a measurable embedding sending the half unit ball onto $\Ball{R/2}{x_0}$ with the volume measure pushing forward to a strictly positive multiple of the restricted volume on $\Ball{R/2}{x_0}$. Strict positivity of the constant scaling implies that the a.e.\ filters coincide, so the essential suprema (and infima) match.
\end{proof}

\begin{lemma}[\leanname{memH01\_indicator\_ball\_of\_closedBall\_subset}, \leanname{ballIndicatorWitnessOn}]
If $\overline{\Ball{r}{c}} \subseteq \Omega$ and $\varphi \in H^{1}_{0}(\Ball{r}{c})$, then the zero extension $\mathbf{1}_{\Ball{r}{c}}\,\varphi$ lies in $H^{1}_{0}(\Omega)$ and carries an explicit $W^{1,2}(\Omega)$-witness whose weak gradient is the zero extension of the gradient of $\varphi$.
\end{lemma}

\begin{proof}
Zero-extension preserves $H^{1}_{0}$ on open sets (the approximating sequence of compactly supported smooth functions extends by zero). Compact support inside $\overline{\Ball{r}{c}} \subseteq \Omega$ together with the universal $W^{1,2}$-witness restricts to a $W^{1,2}(\Omega)$-witness with the same weak gradient.
\end{proof}

\begin{lemma}[\leanname{bilinForm\_ball\_restrict\_eq\_zeroExtend}]
Under the same hypotheses, for $u \in W^{1,2}(\Omega)$ and $\varphi \in H^{1}_{0}(\Ball{r}{c})$,
\[
  \mathfrak{a}_{A\restriction \Ball{r}{c}}\bigl(u\!\restriction, \varphi\bigr)
  = \mathfrak{a}_{A}\bigl(u, \mathbf{1}_{\Ball{r}{c}}\,\varphi\bigr).
\]
\end{lemma}

\begin{proof}
Pointwise, the two integrands agree on $\Ball{r}{c}$ and the right-hand integrand vanishes outside (the indicator-extended witness has zero weak gradient there). Hence both integrals over $\Omega$ equal the integral over $\Ball{r}{c}$ of the same integrand.
\end{proof}

\begin{lemma}[\leanname{MemW1pWitness.of\_ae\_eq}]
If $f =_{a.e.} g$ on $\Omega$ and $f$ has a $W^{1,2}(\Omega)$-witness, then $g$ has a witness with the same weak gradient.
\end{lemma}

\begin{proof}
Lp membership transfers under a.e.\ equality. The weak partial derivative identity $\int g\,\partial_i\varphi = -\int (\nabla f)_i\,\varphi$ follows from $\int g\,\partial_i\varphi = \int f\,\partial_i\varphi$ by a.e.\ equality.
\end{proof}

\begin{theorem}[\leanname{IsSubsolution.congr\_ae}, \leanname{IsSupersolution.congr\_ae}, \leanname{IsSolution.congr\_ae}]
The weak (sub-, super-) solution predicates are stable under a.e.\ equality of the unknown.
\end{theorem}

\begin{proof}
Use the previous lemma to transport witnesses, and observe that $\mathfrak{a}_A$ depends only on the witness's weak gradient, which is unchanged.
\end{proof}

\begin{theorem}[\leanname{IsSubsolution.restrict\_ball}, \leanname{IsSupersolution.restrict\_ball}, \leanname{IsSolution.restrict\_ball}]
Let $\Omega$ be open, $\overline{\Ball{r}{c}} \subseteq \Omega$, and $A$ elliptic on $\Omega$. If $u$ is a weak subsolution (resp.\ super-, solution) of $A$ on $\Omega$, then $u$ is a weak subsolution (resp.\ super-, solution) of $A\!\restriction\!\Ball{r}{c}$ on the smaller domain.
\end{theorem}

\begin{proof}
We treat the subsolution case; the supersolution and solution cases are
symmetric. Pick a witness $h_u^{\Omega}$ for $u$ on $\Omega$ via
\leanname{MemW1p.someWitness}. The restriction lemma
\leanname{MemW1pWitness.restrict} (specialised to the open subball
$\Ball{r}{c} \subseteq \Omega$) yields a witness $h_u^{\Ball{r}{c}}$ for $u$
on the smaller domain, with weak gradient given by the restriction of
$\nabla u$ to $\Ball{r}{c}$; this verifies the $W^{1,2}$ membership half of
\leanname{IsSubsolution}.

For the bilinear-form inequality, fix a nonnegative test
$\varphi \in H^{1}_{0}(\Ball{r}{c})$ together with any witness $h_\varphi$.
The zero extension $\mathbf{1}_{\Ball{r}{c}}\varphi$ lies in $H^{1}_{0}(\Omega)$
and admits an explicit witness $\hat h_\varphi$ produced by
\leanname{ballIndicatorWitnessOn}, whose weak gradient is the zero
extension of $\nabla \varphi$. The pointwise nonnegativity of $\varphi$ on
$\Ball{r}{c}$ extends to a pointwise nonnegativity of the indicator, so
$\mathbf{1}_{\Ball{r}{c}}\varphi$ is a valid nonnegative test on $\Omega$.

Apply the subsolution hypothesis on $\Omega$ to this extended test:
$\mathfrak{a}_{A}(h_u^{\Omega}, \hat h_\varphi) \le 0$. By
\leanname{bilinForm\_ball\_restrict\_eq\_zeroExtend}, the bilinear form for
the restricted coefficient field $A\restriction\Ball{r}{c}$ paired with
$h_u^{\Ball{r}{c}}$ and $h_\varphi$ equals the bilinear form for $A$ paired
with $h_u^{\Omega}$ and $\hat h_\varphi$. Therefore
$\mathfrak{a}_{A\restriction\Ball{r}{c}}(h_u^{\Ball{r}{c}}, h_\varphi) \le 0$,
which is the desired subsolution inequality on the smaller domain.
\end{proof}

\begin{theorem}[\leanname{IsSubsolution.sub\_const\_ball}, \leanname{IsSupersolution.sub\_const\_ball}, \leanname{IsSolution.sub\_const\_ball}]
For $u$ a weak (sub-, super-)solution on $\Ball{r}{c}$ and $k \in \R$, $u - k$ is a weak (sub-, super-)solution on $\Ball{r}{c}$ with the same coefficient field $A$.
\end{theorem}

\begin{proof}
The constant function $k$ has zero weak gradient on the ball (a finite-measure set), so the witness for $u - k$ has the same weak gradient as the witness for $u$. The bilinear form $\mathfrak{a}_A$ depends only on the gradient, so the inequality is unchanged.
\end{proof}

\begin{theorem}[\leanname{IsSubsolution.neg\_ball}, \leanname{IsSupersolution.neg\_ball}, \leanname{IsSolution.neg\_ball}]
On $\Ball{r}{c}$, negation swaps subsolutions and supersolutions: if $u$ is a subsolution, then $-u$ is a supersolution; similarly with the roles reversed; and a solution stays a solution under negation.
\end{theorem}

\begin{proof}
Suppose $u$ is a subsolution: pick a witness $h_u$ for $u$ on $\Ball{r}{c}$
via \leanname{MemW1p.someWitness}. The scaled witness
$h_u.\mathrm{smul}(-1)$ is a witness for $-u$ with weak gradient $-\nabla u$
(\leanname{MemW1pWitness.smul}), establishing the $W^{1,2}$ membership half
of \leanname{IsSupersolution}.

For the bilinear-form inequality, fix a nonnegative test $\varphi$ together
with a witness $h_\varphi$. The negation witness $h_{-(-u)}$ obtained by
applying $\mathrm{smul}(-1)$ a second time agrees a.e.\ with $h_u$
(after clearing the double negation by ring arithmetic), so by
\leanname{MemW1pWitness.of\_ae\_eq} it produces the same bilinear-form
value as $h_u$. The subsolution hypothesis applied to $u$ then gives
$\mathfrak{a}_A(h_u, h_\varphi) \le 0$.

By the linearity lemma \leanname{bilinFormOfCoeff\_smul\_left} applied with
scalar $-1$, we have
\[
  \mathfrak{a}_A(h_u, h_\varphi)
  = -\mathfrak{a}_A(h_{-u}, h_\varphi),
\]
so the inequality $\mathfrak{a}_A(h_u, h_\varphi) \le 0$ rewrites as
$\mathfrak{a}_A(h_{-u}, h_\varphi) \ge 0$, which is the supersolution
inequality for $-u$. The reverse direction (supersolution to subsolution) is
identical with the roles swapped; the solution case follows by combining
both.
\end{proof}


%% ========================================================================
\section{De Giorgi Iteration and Local Boundedness}\label{sec:degiorgi}
%% ========================================================================

% !TEX root = DeGiorgi_Lean_to_Tex.tex

This section is the formal core of the De Giorgi $L^2 \to L^\infty$
argument for non-negative subsolutions of a divergence-form linear
elliptic equation
$\dv(A \nabla u) \le 0$
on the unit ball $\Bz{1} \subset \R^d$, with $A$ measurable, symmetric,
and uniformly elliptic with constants $0 < \lambda \le \Lambda$ (here
$\lambda$ denotes always the global ellipticity constant, never a level
threshold).  The Lean development is organised into six layers, each
realised by a dedicated module:\\
\leanname{DeGiorgiIteration.CutoffAdmissibility},
\leanname{DeGiorgiIteration.Energy},\\
\leanname{DeGiorgiIteration.PreIteration},
\leanname{DeGiorgiIteration.Recurrence}, and
\leanname{DeGiorgiIteration.Linfty},
all reassembled into the public interface
\leanname{DeGiorgiIteration} and reexported by
\leanname{DeGiorgiTheory}.  We follow the same layering below.

Throughout the section $E = \R^d$ with $d \ge 1$, $\Omega \subset E$ is
an open set with finite Lebesgue measure, and $A$ is an elliptic
coefficient field with constants $\lambda$, $\Lambda$ encoded in
\leanname{EllipticCoeff}; we write $\mathrm{ellipticityRatio}(A) =
\Lambda/\lambda$.

\subsection{Positive-part truncation}

\begin{definition}[\leanname{positivePartSub}]
For any function $u : \alpha \to \R$ and any level $\ell \in \R$ we
define the positive-part truncation
\[
  (u-\ell)_+(x) := \max(u(x) - \ell, 0).
\]
\end{definition}

\begin{lemma}[\leanname{positivePartSub\_nonneg}]
$0 \le (u-\ell)_+(x)$ for every $x$.
\end{lemma}

\begin{proof}
Immediate from $\max(\,\cdot\,, 0) \ge 0$.
\end{proof}

\begin{lemma}[\leanname{positivePartSub\_sq\_ge\_gap\_sq\_of\_lt}]
If $\theta < \ell$ and $\ell < u(x)$, then
$(\ell - \theta)^2 \le |(u-\theta)_+(x)|^2$.
\end{lemma}

\begin{proof}
Under the hypotheses $u(x) - \theta > 0$, hence
$(u-\theta)_+(x) = u(x) - \theta \ge \ell - \theta > 0$.  Squaring this
nonnegative inequality yields the claim.
\end{proof}

\begin{lemma}[\leanname{positivePartSub\_antitone}]
If $\ell_1 \le \ell_2$, then
$(u-\ell_2)_+ \le (u-\ell_1)_+$ pointwise.
\end{lemma}

\begin{proof}
$\max(u-\ell_2, 0) \le \max(u-\ell_1, 0)$ since $u-\ell_2 \le u-\ell_1$.
\end{proof}

\begin{lemma}[\leanname{positivePartSub\_le\_posPartZero}]
If $0 \le \ell$, then $(u-\ell)_+ \le u_+ = \max(u,0)$.
\end{lemma}

\begin{lemma}[\leanname{le\_of\_positivePartSub\_eq\_zero}]
If $(u-\ell)_+(x) = 0$, then $u(x) \le \ell$.
\end{lemma}

\begin{proof}
If $u(x) - \ell \ge 0$ then $(u-\ell)_+(x) = u(x) - \ell$, so the
hypothesis forces $u(x) = \ell$.  Otherwise $u(x) < \ell$.
\end{proof}

\subsection{Cutoff admissibility}

The De Giorgi test functions are obtained by multiplying the
positive-part truncation by the square of a smooth scalar cutoff.

\begin{definition}[\leanname{deGiorgiCutoffTestGeneral}]
Given $\eta, u : E \to \R$ and a level $\ell$,
\[
  \varphi_{\eta, u, \ell}(x) := \eta(x) \cdot \bigl( \eta(x) \cdot
  (u-\ell)_+(x) \bigr) = \eta(x)^2 \, (u-\ell)_+(x).
\]
\end{definition}

\begin{lemma}[\leanname{deGiorgiCutoffTestGeneral\_nonneg}]
If $\eta \ge 0$ pointwise, then $\varphi_{\eta, u, \ell} \ge 0$.
\end{lemma}

\begin{lemma}[\leanname{deGiorgiCutoffTestGeneral\_hasCompactSupport}]
If $\eta$ has compact support, so does $\varphi_{\eta,u,\ell}$.
\end{lemma}

\begin{lemma}[\leanname{deGiorgiCutoffTestGeneral\_tsupport\_subset}]
If $\tsupp\,\eta \subseteq \Omega$, then $\tsupp\,\varphi_{\eta,u,\ell}
\subseteq \Omega$.
\end{lemma}

The technical core of this layer is a Sobolev membership statement: if
$u \in W^{1,2}(\Omega)$ and the positive-part truncation
$(u - \ell)_+$ is also in $W^{1,2}$, then for any smooth bounded
cutoff $\eta$ with $\tsupp\,\eta \subseteq \Omega$, the test function
$\varphi_{\eta,u,\ell}$ lies in $W^{1,2}_0(\Omega)$.  In Lean this is
phrased through the witness type \leanname{MemW1pWitness} and the
zero-trace structure \leanname{MemW01p}.

\begin{theorem}[\leanname{deGiorgiCutoffTest\_memW01p\_of\_truncWitness}]
Let $\Omega \subset E$ be open with finite measure.  Let $u$ be such
that $(u-\ell)_+$ admits a $W^{1,2}(\Omega)$ witness.  Let $\eta \in
C^\infty(E)$ satisfy $|\eta| \le C_0$, $\|\nabla \eta\| \le C_1$,
$\eta$ has compact support, and $\tsupp\,\eta \subseteq \Omega$.
Then $\varphi_{\eta,u,\ell} \in \Wop{2}(\Omega)$.
\end{theorem}

\begin{proof}
By two successive applications of the closure of $W^{1,2}$-witnesses
under multiplication by smooth bounded functions
(\leanname{MemW1pWitness.mul\_smooth\_bounded}), starting from $(u-\ell)_+$
and multiplying twice by $\eta$, one obtains a $W^{1,2}$ witness for
$\eta^2 (u-\ell)_+$.  Since $\varphi_{\eta,u,\ell}$ is compactly
supported with $\tsupp \subseteq \Omega$, the localisation lemma
\leanname{memW01p\_of\_memW1p\_of\_tsupport\_subset} promotes the
$W^{1,2}$-witness to a $W^{1,2}_0(\Omega)$ membership.
\end{proof}

\begin{theorem}[\leanname{deGiorgiCutoffTest\_memW01p\_of\_posPart}]
Same conclusion, but assuming a $W^{1,2}$-witness for $u$ together
with a generic positive-part closure rule
$\bigl(\, v \in W^{1,2}(\Omega) \implies (v)_+ \in W^{1,2}(\Omega)\,\bigr)$.
\end{theorem}

\begin{proof}
Apply the closure rule to $v = u-\ell$ to obtain a witness for
$(u-\ell)_+$, then invoke the previous theorem.
\end{proof}

A concrete version, \leanname{deGiorgiCutoffTest\_memW01p\_of\_posPartApprox},
is obtained by furnishing the $W^{1,2}$-witness for $(u-\ell)_+$ via the
Chapter~02 smooth approximation package for $u-\ell$.  The
ball-restricted variants
\leanname{deGiorgiCutoffTest\_memW01p\_on\_ball\_of\_ballPosPart} and
\leanname{deGiorgiCutoffTest\_memW01p\_on\_ball\_of\_ballPosPartApprox}
specialise these statements to $\Omega = \Ball{s}{x_0}$ and supply
finiteness of the restricted measure automatically.

\begin{definition}[\leanname{positivePartSub\_memW1pWitness\_on\_ball}]
On a ball $\Ball{s}{x_0}$, given a $W^{1,2}$-witness for $u$ and a
smooth approximation package for $u - \ell$, the lemma constructs a
$W^{1,2}$-witness for $(u - \ell)_+$ on the ball.
\end{definition}

\subsection{Weighted Caccioppoli inequality}

We now turn to the analytical heart of the chapter: the weighted
Caccioppoli inequality for truncations of subsolutions.  We use the
abbreviation
\[
  \nabla \eta(x) := (\partial_1 \eta(x), \dots, \partial_d \eta(x))
\]
identified with the inner-product representation of the Fréchet
derivative; see \leanname{deGiorgiFderivVec}.  The norm
$\|\nabla \eta(x)\|$ then equals the operator norm
$\|D\eta(x)\|$, which is recorded as
\leanname{deGiorgi\_norm\_fderivVec\_eq\_norm\_fderiv}.

\begin{lemma}[\leanname{weighted\_caccioppoli\_pointwise} (Young's inequality)]
Let $\Lambda > 0$, $\eta, \zeta, E, M, v \in \R$, and assume the
pointwise coercivity condition $|M|^2 \le \Lambda E$. Then
\begin{align*}
  2 \eta |v| \zeta |M|
   \le \tfrac{1}{2}\, \eta^2 E + 2\Lambda\, \zeta^2 |v|^2.
\end{align*}
\end{lemma}

\begin{proof}
Apply Young's inequality $2ab \le a^2 + b^2$ with $a =
\eta |M| / \sqrt{2\Lambda}$ and $b = \sqrt{2\Lambda}\, \zeta |v|$ to
get
\[
  2 \eta |v| \zeta |M|
   \le \frac{\eta^2 |M|^2}{2\Lambda} + 2\Lambda\, \zeta^2 |v|^2.
\]
The hypothesis $|M|^2 \le \Lambda E$ then yields $\eta^2 |M|^2 /
(2\Lambda) \le \tfrac{1}{2} \eta^2 E$, completing the proof.
\end{proof}

\begin{theorem}[Abstract absorption,
\leanname{weighted\_caccioppoli\_absorb}]
Let $(\alpha, \mu)$ be a measure space and let $E, M, \eta, \zeta, v
: \alpha \to \R$ be measurable functions with $\Lambda > 0$. Assume
the integrability conditions and:
\begin{enumerate}
\item $\int \eta^2 E \, d\mu \le \int 2 \eta |v| \zeta |M| \, d\mu$,
\item $|M|^2 \le \Lambda E$ almost everywhere.
\end{enumerate}
Then $\int \eta^2 E \, d\mu \le 4 \Lambda \int \zeta^2 |v|^2 \, d\mu$.
\end{theorem}

\begin{proof}
Apply the pointwise Young inequality almost everywhere, integrate to
get
\[
  \int 2\eta|v|\zeta|M|\, d\mu
   \le \tfrac{1}{2}\int \eta^2 E\, d\mu
        + 2\Lambda \int \zeta^2 |v|^2 \, d\mu,
\]
combine with the assumed core estimate to obtain
$\int \eta^2 E\, d\mu \le \tfrac12 \int \eta^2 E\, d\mu + 2\Lambda
\int \zeta^2 |v|^2 \, d\mu$, and absorb to conclude.
\end{proof}

The PDE-facing weighted Caccioppoli inequality is the following
statement, which is the analytical core of the De Giorgi argument.

\begin{theorem}[Weighted Caccioppoli for subsolutions,
\leanname{caccioppoli\_weighted\_of\_subsolution}]
Let $\Omega \subset E$ be open with finite measure, $A$ an elliptic
coefficient field on $\Omega$, $u$ a subsolution
($\leanname{IsSubsolution}\, A\, u$), and $\eta$ a smooth bounded
nonneg function with $|\eta| \le C_0$, $\|\nabla \eta\| \le C_1$,
$\tsupp\,\eta \subseteq \Omega$.  Let $\ell \in \R$ and assume
$(u - \ell)_+ \in W^{1,2}(\Omega)$ and that $\varphi_{\eta,u,\ell}$
admits a $W^{1,2}_0(\Omega)$ structure.  Then
\begin{multline*}
  \int_\Omega \eta(x)^2 \, \bigl\| \nabla(u-\ell)_+(x) \bigr\|^2 \, dx
   \le \\
   4\, \frac{\Lambda}{\lambda}\,
     \int_\Omega \|\nabla \eta(x)\|^2 \,
       \bigl|(u-\ell)_+(x)\bigr|^2 \, dx.
\end{multline*}
\end{theorem}

\begin{proof}
Test the subsolution inequality against
$\varphi := \varphi_{\eta,u,\ell} = \eta^2 (u-\ell)_+$, which is
admissible by the cutoff admissibility theorem and is nonnegative.
The subsolution condition gives
$\int_\Omega \langle A\nabla u, \nabla \varphi\rangle \, dx \le 0$.
Computing the gradient of $\varphi$,
\begin{align*}
  \nabla \varphi(x)
   = \eta(x)^2 \, \nabla (u-\ell)_+(x)
     + 2\, \eta(x)\, (u-\ell)_+(x) \, \nabla \eta(x),
\end{align*}
which is the content of
\leanname{deGiorgiCutoffTestWitnessWeighted\_grad}. Combining with the
identification $\nabla u = \nabla(u - \ell)_+$ on $\{u > \ell\}$ and
$\nabla(u-\ell)_+ = 0$ on $\{u \le \ell\}$ (the lemmas
\leanname{truncGrad\_eq\_on\_superlevel} and
\leanname{truncGrad\_zero\_on\_sublevel}), we may rewrite
\begin{multline*}
  0 \ge
   \int_\Omega \eta(x)^2 \,
     \langle A(x) \nabla(u-\ell)_+(x), \nabla(u-\ell)_+(x)\rangle \, dx \\
   + \int_\Omega 2 \eta(x) (u-\ell)_+(x)\,
     \langle A(x) \nabla(u-\ell)_+(x), \nabla \eta(x)\rangle \, dx.
\end{multline*}
By the coercivity bound for $A$ (lemma
$\Lambda$-coercive: $\langle Aw,w\rangle \ge \lambda \|w\|^2$) and the
upper bound (lemma $\langle Aw,w\rangle \le \Lambda\|w\|^2$ together
with $\|Aw\|^2 \le \Lambda \langle Aw,w\rangle$), the first term is
bounded below by $\lambda \int \eta^2 \|\nabla(u-\ell)_+\|^2$, while
the second is dominated in absolute value by
\[
  \int_\Omega 2\, \eta(x) (u-\ell)_+(x) \,
    \|A(x)\nabla(u-\ell)_+(x)\|\,\|\nabla \eta(x)\| \, dx.
\]
Setting $E := \langle A\nabla(u-\ell)_+, \nabla(u-\ell)_+\rangle$ and
$M := \|A\nabla(u-\ell)_+\|$, we have $|M|^2 \le \Lambda E$, and
applying the abstract absorption lemma above with these choices
yields
\[
  \int_\Omega \eta(x)^2\, E(x)\, dx
   \le 4\,\Lambda \int_\Omega \|\nabla \eta(x)\|^2\,
     |(u-\ell)_+(x)|^2 \, dx.
\]
Coercivity now bounds the left-hand side from below by
$\lambda \int_\Omega \eta^2 \|\nabla(u-\ell)_+\|^2$, and dividing by
$\lambda$ produces the announced inequality with constant
$4 \Lambda/\lambda$.
\end{proof}

\begin{theorem}[Ball Caccioppoli,
\leanname{caccioppoli\_weighted\_on\_ball\_of\_ballPosPart}]
The same conclusion holds with $\Omega = \Ball{s}{x_0}$ and a witness
for $(u-\ell)_+$ on the ball, provided $\eta$ is smooth, nonnegative,
$|\eta| \le 1$, $\|\nabla \eta\| \le C_\eta$, and $\tsupp\,\eta
\subseteq \Ball{s}{x_0}$.
\end{theorem}

The approximation-based variant is
\leanname{caccioppoli\_weighted\_on\_ball\_of\_posPartApprox}.

\begin{theorem}[Localized energy estimate]
(\leanname{deGiorgi\_energy\_estimate\_on\_concentricBalls\_of\_ballPosPart})\\
Let $0 < r < s$, $A$ elliptic on $\Ball{s}{x_0}$, $u$ a subsolution.
Let $\eta$ be smooth, $0 \le \eta \le 1$ on $E$, $\eta = 1$ on
$\Ball{r}{x_0}$, $\tsupp\,\eta \subset \Ball{s}{x_0}$, and
$\|\nabla \eta\| \le C_\eta$. Then
\begin{multline*}
  \int_{\Ball{r}{x_0}} \bigl\|\nabla(u-\ell)_+(x)\bigr\|^2 \, dx \\
  \le 4\,\frac{\Lambda}{\lambda} \, C_\eta^2 \,
    \int_{\Ball{s}{x_0}} \bigl|(u-\ell)_+(x)\bigr|^2 \, dx.
\end{multline*}
\end{theorem}

\begin{proof}
Since $\eta = 1$ on $\Ball{r}{x_0}$, the integral on the left equals
$\int_{\Ball{r}{x_0}} \eta^2 \|\nabla(u-\ell)_+\|^2$. Monotonicity in
the integration set and the weighted Caccioppoli inequality bound
this by $4(\Lambda/\lambda)\int_{\Ball{s}{x_0}} \|\nabla \eta\|^2
|(u-\ell)_+|^2$. The pointwise bound $\|\nabla\eta\|^2 \le C_\eta^2$
on $\Ball{s}{x_0}$ then gives the claim.
\end{proof}

\subsection{Chebyshev bound on superlevel sets}

\begin{theorem}[\leanname{superlevel\_measureReal\_le\_integral\_posPart}]
Let $(\alpha, \mu)$ be a finite measure space, $u : \alpha \to \R$,
and let $\theta < \ell$. If $|(u-\theta)_+|^2$ is integrable, then
\[
  \mu\bigl(\{x : \ell < u(x)\}\bigr)
   \le \frac{1}{(\ell-\theta)^2}\,
       \int |(u-\theta)_+(x)|^2 \, d\mu(x).
\]
\end{theorem}

\begin{proof}
Set $f(x) := |(u-\theta)_+(x)|^2$ and $\varepsilon := (\ell-\theta)^2$. The
hypothesis $\theta < \ell$ ensures $\varepsilon > 0$. By
\leanname{positivePartSub\_sq\_ge\_gap\_sq\_of\_lt}, on the superlevel set
$\{x : \ell < u(x)\}$ the value of $f$ is at least
$(\ell - \theta)^2 = \varepsilon$, so we have the set inclusion
\[
  \{x : \ell < u(x)\}
    \subseteq \{x : \varepsilon \le f(x)\}.
\]
By monotonicity of measure (\leanname{measureReal\_mono}) this gives
$\mu\{\ell < u\} \le \mu\{\varepsilon \le f\}$.

Markov's inequality for nonnegative integrable functions
(\leanname{MeasureTheory.mul\_meas\_ge\_le\_integral\_of\_nonneg}) applied to
$f$ at threshold $\varepsilon$ yields
\[
  \varepsilon \cdot \mu\{\varepsilon \le f\}
   \le \int_{\alpha} f(x)\,d\mu(x).
\]
Combining the set inclusion with monotone multiplication by $\varepsilon \ge 0$,
\[
  \varepsilon \cdot \mu\{\ell < u\}
   \le \varepsilon \cdot \mu\{\varepsilon \le f\}
   \le \int_{\alpha} f\,d\mu.
\]
Dividing by the positive constant $\varepsilon$ produces the desired bound
$\mu\{\ell < u\} \le ((\ell - \theta)^2)^{-1}\int |(u-\theta)_+|^2\,d\mu$.
\end{proof}

\subsection{Pre-iteration estimate}

We now combine the energy estimate with a Sobolev/Holder
interpolation step to obtain a single-step nonlinear inequality
between truncation energies at two levels and on two concentric balls.

\begin{theorem}[Abstract pre-iteration combination,
\leanname{deGiorgi\_preiter\_abstract}]
Let $d > 0$ and let $C_{\mathrm{sob}}, C_{\mathrm{en}}, I_\theta, m,
\theta < \ell$ all be nonneg numbers. Suppose
\begin{align*}
  I_\ell &\le C_{\mathrm{sob}} \cdot G \cdot m^{2/d}, \\
  m &\le (\ell - \theta)^{-2}\, I_\theta, \\
  G &\le C_{\mathrm{en}} \cdot I_\theta.
\end{align*}
Then
\[
  I_\ell \le C_{\mathrm{sob}}\, C_{\mathrm{en}}\,
    \bigl((\ell-\theta)^{-2} I_\theta\bigr)^{2/d}\, I_\theta.
\]
\end{theorem}

\begin{proof}
Raise the second inequality to the power $2/d \ge 0$, multiply by
$G$, then by $C_{\mathrm{sob}}$, using $C_{\mathrm{en}}$ to bound
$G$.
\end{proof}

\begin{theorem}[Packaged pre-iteration,
\leanname{deGiorgi\_preiter\_of\_energy}]
With $\mu$ a finite measure and the same notation, if
\begin{itemize}
\item $I_\ell \le C_{\mathrm{sob}}\, G\, \mu(\{\ell < u\})^{2/d}$,
\item $G \le C_{\mathrm{en}} \cdot I_\theta$,
\item $\int |(u-\theta)_+|^2 \, d\mu \le I_\theta$,
\end{itemize}
then
\(
  I_\ell \le C_{\mathrm{sob}}\, C_{\mathrm{en}}\,
    ((\ell-\theta)^{-2} I_\theta)^{2/d}\, I_\theta.
\)
\end{theorem}

\begin{proof}
By the Chebyshev bound, $\mu(\{\ell < u\}) \le (\ell-\theta)^{-2}\,
\int|(u-\theta)_+|^2 \, d\mu \le (\ell-\theta)^{-2}\, I_\theta$. The
abstract combination lemma now applies.
\end{proof}

The cutoff Sobolev step required to discharge the Sobolev hypothesis
$I_\ell \le (C_{\mathrm{gns}})^2 \int \|\nabla \varphi\|^2\,
\mu(\{\ell < u\})^{2/d}$ is

\begin{theorem}[\leanname{deGiorgi\_cutoffSobolev\_on\_concentricBalls\_of\_ballPosPart}]
Under the standing hypotheses ($d > 2$, $0 < r < s$, $u$ admitting a
$W^{1,2}$ witness, $(u-\theta)_+$ admitting a $W^{1,2}$ witness on
$\Ball{s}{x_0}$, $\theta < \ell$, $\eta$ smooth with $\eta = 1$ on
$\Ball{r}{x_0}$, $|\eta| \le 1$, $\tsupp \eta \subseteq
\Ball{s}{x_0}$), there exists a $W^{1,2}$-witness $w_{\eta,\theta}$
for $\eta^2 (u-\theta)_+$ such that
\begin{multline*}
  \int_{\Ball{r}{x_0}} \bigl|(u-\ell)_+(x)\bigr|^2 \, dx
  \le (C_{\mathrm{gns}}(d,2))^2 \cdot \\
   \Bigl(\int_{\Ball{s}{x_0}} \|\nabla w_{\eta,\theta}(x)\|^2 \, dx
   \Bigr) \cdot
    \mu_s\bigl(\{\ell < u\}\bigr)^{2/d},
\end{multline*}
where $\mu_s$ denotes Lebesgue measure restricted to $\Ball{s}{x_0}$.
\end{theorem}

\begin{proof}
The cutoff $\varphi := \eta^2 (u-\theta)_+$ has compact support inside
$\Ball{s}{x_0}$ and lies in $\Wop{2}(\Ball{s}{x_0})$ by
the cutoff admissibility theorem, hence its zero extension to $\R^d$
lies in $W^{1,2}(\R^d)$. The Gagliardo--Nirenberg--Sobolev embedding
gives the global $L^{2^*}$ estimate
$\|\varphi\|_{L^{2^*}} \le C_{\mathrm{gns}}(d,2) \|\nabla
\varphi\|_{L^2}$, where $2^* = 2d/(d-2)$. On $\Ball{r}{x_0}$ one has
$\eta = 1$ and $u > \ell > \theta$ implies $\varphi(x) = (u(x) -
\theta) > \ell - \theta$, so $|(u-\ell)_+|^2 \le |\varphi|^2$ on
$\{\ell < u\}$. Restricting the $L^{2^*}$ estimate to $\Ball{r}{x_0}$
and applying Holder's inequality with exponents $2^*/2$ and its dual
gives the announced bound, where the volume factor reduces to
$\mu_s(\{\ell < u\})^{2/d}$.
\end{proof}

\begin{theorem}[PDE pre-iteration on concentric balls,
\leanname{deGiorgi\_preiter\_on\_concentricBalls\_of\_ballPosPart}]
Under the same hypotheses ($d > 2$, $0 < r < s$, $\theta < \ell$,
$\eta$ as above), one has
\begin{multline*}
  \int_{\Ball{r}{x_0}} |(u-\ell)_+(x)|^2 \, dx
  \le (C_{\mathrm{gns}}(d,2))^2
    \cdot \bigl(8(\Lambda/\lambda + 1)\, C_\eta^2\bigr) \cdot \\
    \bigl((\ell - \theta)^{-2}
       \int_{\Ball{s}{x_0}} |(u-\theta)_+|^2 \, dx\bigr)^{2/d}
    \cdot \int_{\Ball{s}{x_0}} |(u-\theta)_+|^2 \, dx.
\end{multline*}
\end{theorem}

\begin{proof}
Apply the cutoff Sobolev step with witness $w_{\eta,\theta}$, the
energy estimate $\int_{\Ball{s}{x_0}} \|\nabla w_{\eta,\theta}\|^2 \le
8(\Lambda/\lambda+1) C_\eta^2 \int_{\Ball{s}{x_0}} |(u-\theta)_+|^2$
(coming from the gradient identity for $\varphi$ together with the
weighted Caccioppoli inequality and the elementary bound
$(a+b)^2 \le 2a^2 + 2b^2$), and the packaged pre-iteration. The
identification of the gradient bound coefficient is recorded as
\leanname{deGiorgi\_cutoff\_gradient\_bound\_of\_weighted}.
\end{proof}

\subsection{De Giorgi iteration sequences}

We now fix the canonical scales of the iteration.

\begin{definition}[\leanname{deGiorgiTail}, \leanname{deGiorgiRadius},
\leanname{deGiorgiLevel}, \leanname{deGiorgiRecurrenceBase},
\leanname{deGiorgiRecurrenceCoeff}]
For $n \in \N$, $\ell^* > 0$, and $K > 0$, set
\begin{align*}
  \tau_n &:= (1/2)^{n+1}, \\
  r_n &:= (1 + \tau_n)/2, \\
  \ell_n(\ell^*) &:= \ell^* (1 - \tau_n), \\
  B_d &:= 2^{2 + 4/d}, \\
  C_d(K, \ell^*) &:= K \cdot 2^{6 + 8/d} \cdot (\ell^*)^{-4/d}.
\end{align*}
\end{definition}

\begin{lemma}
The radii $r_n$ satisfy $1/2 < r_n < 1$ and decrease to $1/2$, the
levels increase from $\ell_0 < \ell^*$ to $\ell^*$, and one has
the gap identities (\leanname{deGiorgiRadius\_gap},
\leanname{deGiorgiLevel\_gap})
\[
  r_n - r_{n+1} = (1/2)^{n+3}, \qquad
  \ell_{n+1}(\ell^*) - \ell_n(\ell^*) = \ell^* (1/2)^{n+2}.
\]
\end{lemma}

\begin{definition}[\leanname{deGiorgiEnergySeq}]
For a function $u$, base point $x_0$, and threshold $\ell^* > 0$, set
\[
  A_n(u, x_0, \ell^*) := \int_{\Ball{r_n}{x_0}}
     \bigl|(u - \ell_n(\ell^*))_+(x)\bigr|^2 \, dx.
\]
\end{definition}

\subsection{Recurrence and abstract closeout}

\begin{theorem}[Geometric closeout for nonlinear recurrences,
\leanname{deGiorgi\_recurrence\_closeout}]
Let $Y_n \ge 0$, $C > 0$, $B > 1$, $\alpha > 0$, and assume
\[
  Y_{n+1} \le C\, B^n\, Y_n^{1+\alpha}, \qquad
  Y_0 \le C^{-1/\alpha}\, B^{-1/\alpha^2}.
\]
Then $Y_n \to 0$ as $n \to \infty$.
\end{theorem}

\begin{proof}
Set $X := C B^{1/\alpha}$, $D := X^{1/\alpha}$, $q :=
B^{-1/\alpha} < 1$, and define $W_n := B^{n/\alpha} Y_n$ and $Z_n :=
D W_n$. The hypothesis on $Y_0$ is equivalent to $Y_0 \le D^{-1}$, so
$Z_0 \le 1$. From the recurrence
\[
  W_{n+1} \le X \, W_n^{1+\alpha},
\]
multiply by $D$ and use $D X = D^{1+\alpha}$ to get $Z_{n+1} \le
Z_n^{1+\alpha}$, hence by induction $Z_n \le 1$ for all $n$. This
yields $W_n \le D^{-1}$ and consequently
$Y_n \le D^{-1} q^n$, which tends to $0$ as $n \to \infty$.
\end{proof}

\begin{theorem}[Recurrence rewriting,
\leanname{deGiorgi\_preiter\_to\_recurrence}]
If a sequence $A_n$ satisfies, for all $n$,
\begin{multline*}
  A_{n+1} \le
   \frac{K}{(r_n - r_{n+1})^2 (\ell_{n+1} - \ell_n)^{4/d}}
   \cdot A_n^{1 + 2/d},
\end{multline*}
then with $C_d(K,\ell^*)$ and $B_d$ as above one has
\[
  A_{n+1} \le C_d(K, \ell^*)\, B_d^n\, A_n^{1 + 2/d}.
\]
\end{theorem}

\begin{proof}
By the gap identities,
$(r_n - r_{n+1})^2 = 2^{-(2n+6)}$ and
$(\ell_{n+1}-\ell_n)^{4/d} = (\ell^*)^{4/d}\, 2^{-(n+2)\cdot 4/d}$,
so the denominator equals
$(\ell^*)^{4/d}\, 2^{-(6+8/d + n(2+4/d))}$. Inverting and rearranging
yields the claimed form.
\end{proof}

\begin{theorem}[Vanishing of the canonical recurrence,
\leanname{deGiorgi\_preiter\_vanishing}]
Under $d > 2$, $K > 0$, $\ell^* > 0$, the canonical recurrence above,
and the smallness condition
\[
  A_0 \le C_d(K,\ell^*)^{-d/2}\, B_d^{-(d/2)^2 / 2},
\]
the sequence $A_n$ tends to zero.
\end{theorem}

\begin{proof}
The pre-iteration recurrence is rewritten using
\leanname{deGiorgi\_preiter\_to\_recurrence} into the standard form with
$B := B_d$, $C := C_d(K,\ell^*)$, $\alpha := 2/d$. The smallness
hypothesis is exactly the threshold required by
\leanname{deGiorgi\_recurrence\_closeout}, so $A_n \to 0$.
\end{proof}

\subsection{The De Giorgi $L^\infty$ bound}

The next layer assembles the $L^\infty$ bound on the half-ball from
the recurrence machinery.

\begin{lemma}[\leanname{ae\_eq\_zero\_of\_forall\_setIntegral\_sq\_le\_of\_tendsto\_zero}]
Let $f : E \to \R$ with $|f|^2 \in L^1(s)$ and let $Y_n \to 0$ be a
sequence of upper bounds for $\int_s |f|^2$. Then $f = 0$ a.e. on
$s$.
\end{lemma}

\begin{lemma}[\leanname{integral\_sq\_halfBall\_le\_deGiorgiEnergySeq}]
For $\ell^* > 0$, the half-ball energy is monotone in the iteration:
\[
  \int_{\Ball{1/2}{x_0}} |(u-\ell^*)_+|^2 \, dx
   \le A_n(u, x_0, \ell^*),
\qquad \forall n.
\]
\end{lemma}

\begin{lemma}[\leanname{deGiorgiEnergySeq\_zero\_le\_integral\_sq\_posPartZero\_on\_unitBall}]
The first term satisfies
$A_0 \le \int_{\Bz{1}} (u_+)^2 \, dx$.
\end{lemma}

\begin{definition}[\leanname{C\_DeGiorgiSmallness}]
For $K > 0$ and $d \ge 1$, set
\[
  C_{\mathrm{small}}(d, K) :=
   \bigl(K \cdot 2^{6 + 8/d}\bigr)^{d/4}
   \cdot B_d^{d^2/8}.
\]
\end{definition}

This is the height constant required to convert a unit-ball $L^2$
smallness assumption into the abstract initial recurrence threshold.
We collect three monotonicity/scaling identities used later:
\leanname{C\_DeGiorgiSmallness\_pos},
\leanname{C\_DeGiorgiSmallness\_mono} (monotonicity in $K$), and
\leanname{C\_DeGiorgiSmallness\_mul\_right}:
\[
  C_{\mathrm{small}}(d, K t) =
    C_{\mathrm{small}}(d, K) \, t^{d/4} \quad (K, t \ge 0).
\]

\begin{definition}[Iteration constants for subsolutions,
\leanname{K\_DeGiorgi\_subsolution} and \leanname{C\_DeGiorgi\_subsolution}]
For an elliptic coefficient field $A$ on $\Bz{1}$ set
\[
  K_{\mathrm{DG}}(d, A) := (C_{\mathrm{gns}}(d, 2))^2 \cdot
    8\bigl(\Lambda/\lambda + 1\bigr) \cdot (2 M_{\mathrm{st}})^2 + 1,
\]
\[
  C_{\mathrm{DG}}(d, A) := C_{\mathrm{small}}(d, K_{\mathrm{DG}}(d, A)).
\]
Both are positive (\leanname{K\_DeGiorgi\_subsolution\_pos},
\leanname{C\_DeGiorgi\_subsolution\_pos}). The constant
$M_{\mathrm{st}}$ is the standard cutoff gradient bound from the
Stampacchia layer.
\end{definition}

\begin{theorem}[Iteration package,
\leanname{deGiorgi\_iteration\_package\_on\_unitBall\_of\_subsolution}]
Let $d > 2$, $A$ elliptic on $\Bz{1}$, and $u$ a subsolution. Then
$K_{\mathrm{DG}}(d,A) > 0$ and for every $\ell^* > 0$:
\begin{enumerate}
\item $|(u-\ell^*)_+|^2$ is integrable on $\Ball{1/2}{0}$;
\item for every $n \in \N$, $|(u-\ell_n(\ell^*))_+|^2$ is integrable
on $\Ball{r_n}{0}$;
\item the canonical pre-iteration recurrence holds:
\begin{multline*}
  A_{n+1}(u,0,\ell^*) \le \\
   \frac{K_{\mathrm{DG}}(d,A)}{(r_n - r_{n+1})^2 \,
     (\ell_{n+1}(\ell^*)-\ell_n(\ell^*))^{4/d}} \cdot
   A_n(u,0,\ell^*)^{1 + 2/d}.
\end{multline*}
\end{enumerate}
\end{theorem}

\begin{proof}
Take a $W^{1,2}(\Bz{1})$ witness for $u$ from
$\leanname{IsSubsolution.memW1p}$. The smooth approximation package
on the unit ball produces, for any shift $\theta$, an approximating
sequence $\psi_n \in C^1_c$ with $\psi_n \to u - \theta$ in $L^2$ and
$\nabla \psi_n \to \nabla u$ in $L^2$. From this one constructs
witnesses for $(u-\ell^*)_+$ and $(u-\ell_n(\ell^*))_+$ on $\Bz{1}$
via \leanname{positivePartSub\_memW1pWitness\_on\_ball}, and the
required integrability follows by restriction to $\Ball{r_n}{0}$. For
the recurrence, fix $n$ and apply the PDE pre-iteration theorem
\leanname{deGiorgi\_preiter\_on\_concentricBalls\_of\_ballPosPart} to $r =
r_{n+1}$, $s = r_n$, $\theta = \ell_n(\ell^*)$, $\ell =
\ell_{n+1}(\ell^*)$, with a smooth Stampacchia cutoff $\eta_n$
satisfying $\eta_n = 1$ on $\Ball{r_{n+1}}{0}$, $\tsupp\,\eta_n
\subset \Ball{r_n}{0}$, and $\|\nabla \eta_n\| \le 2 M_{\mathrm{st}}
/ (r_n - r_{n+1})$, so that the constant becomes
$(C_{\mathrm{gns}}(d,2))^2 \cdot 8(\Lambda/\lambda + 1)
(2 M_{\mathrm{st}})^2 \le K_{\mathrm{DG}}(d,A)$.
\end{proof}

\begin{theorem}[Vanishing version,
\leanname{linfty\_subsolution\_DeGiorgi\_ae\_zero\_of\_smallness}]
With $d > 2$, $K, \ell^* > 0$, the integrability assumptions on the
truncations, the canonical recurrence, and the abstract smallness
condition
\[
  A_0(u,x_0,\ell^*) \le
    C_d(K,\ell^*)^{-d/2}\, B_d^{-(d/2)^2 / 2},
\]
one has $(u-\ell^*)_+ = 0$ almost everywhere on $\Ball{1/2}{x_0}$.
\end{theorem}

\begin{proof}
By \leanname{deGiorgi\_preiter\_vanishing}, $A_n \to 0$. Combined with
the half-ball monotonicity lemma, this gives
$\int_{\Ball{1/2}{x_0}} |(u-\ell^*)_+|^2 = 0$, whence
$(u-\ell^*)_+ = 0$ a.e.\ on the half-ball by the abstract zero
lemma.
\end{proof}

\begin{corollary}[\leanname{linfty\_subsolution\_DeGiorgi\_ae\_of\_smallness}]
Under the same hypotheses, $\max(u, 0) \le \ell^*$ a.e.\ on
$\Ball{1/2}{x_0}$.
\end{corollary}

\begin{proof}
$(u-\ell^*)_+(x) = 0$ implies $u(x) \le \ell^*$ by
\leanname{le\_of\_positivePartSub\_eq\_zero}, and $0 \le \ell^*$.
\end{proof}

\begin{theorem}[Unit-ball smallness packaging,
\leanname{linfty\_subsolution\_DeGiorgi\_ae\_of\_unitBall\_initial\_energy\_smallness}]
Replacing the abstract smallness on $A_0$ by the same threshold for
$\int_{\Bz{1}} (u_+)^2$, the same conclusion holds.
\end{theorem}

\begin{proof}
The bound $A_0 \le \int_{\Bz{1}} (u_+)^2$ converts the assumption
into the previous form.
\end{proof}

\begin{theorem}[Height-threshold form,
\leanname{linfty\_subsolution\_DeGiorgi\_ae\_of\_unitBall\_height\_bound}]
Under the same setup, if
\[
  C_{\mathrm{small}}(d,K)\,
    \Bigl(\int_{\Bz{1}} (u_+)^2\Bigr)^{1/2} \le \ell^*,
\]
then $\max(u,0) \le \ell^*$ a.e.\ on $\Ball{1/2}{x_0}$.
\end{theorem}

\begin{proof}
Squaring and using the formula for $C_{\mathrm{small}}$ shows that
the height-bound implies the abstract smallness condition above.
This step is the lemma
\leanname{deGiorgi\_initial\_energy\_small\_of\_height\_bound}.
\end{proof}

\begin{theorem}[De Giorgi $L^\infty$ bound,
\leanname{linfty\_subsolution\_DeGiorgi}]
Let $d > 2$, $A$ elliptic on $\Bz{1}$, and $u$ a subsolution with
$|\max(u,0)|^2 \in L^1(\Bz{1})$. Then
\[
  \max(u(x), 0) \le C_{\mathrm{DG}}(d, A) \cdot
    \Bigl(\int_{\Bz{1}} |\max(u,0)|^2\, dx\Bigr)^{1/2}
\]
for almost every $x \in \Ball{1/2}{0}$.
\end{theorem}

\begin{proof}
Set $I_0 := \int_{\Bz{1}} |u_+|^2 \, dx$. If $I_0 = 0$ then by
monotonicity $|u_+|^2 = 0$ a.e.\ on $\Ball{1/2}{0}$, so the
inequality holds trivially. Otherwise let $\ell^* := C_{\mathrm{DG}}
(d,A) \sqrt{I_0} > 0$. The iteration package supplies the integrability
of all canonical truncations and the canonical recurrence. The
height-threshold theorem then yields $\max(u,0) \le \ell^* =
C_{\mathrm{DG}}(d,A) \sqrt{I_0}$ a.e.\ on $\Ball{1/2}{0}$.
\end{proof}

\begin{corollary}[Existence of a level threshold,
\leanname{linfty\_subsolution\_DeGiorgi\_exists\_threshold}]
Under the same hypotheses there exists $\ell^* > 0$ with
$C_{\mathrm{DG}}(d,A) \sqrt{I_0} \le \ell^*$ such that
$\max(u,0) \le \ell^*$ a.e.\ on $\Ball{1/2}{0}$.
\end{corollary}

\subsection{Normalised constants}

\begin{definition}[\leanname{K\_DeGiorgi\_subsolution\_normalized},
\leanname{C\_DeGiorgi\_subsolution\_normalized}]
The dimension-only counterparts are
\[
  K^{\mathrm{norm}}_{\mathrm{DG}}(d) :=
    2 \cdot (C_{\mathrm{gns}}(d,2))^2 \cdot 8\, (2 M_{\mathrm{st}})^2 + 1,
\]
\[
  C^{\mathrm{norm}}_{\mathrm{DG}}(d) :=
    C_{\mathrm{small}}(d, K^{\mathrm{norm}}_{\mathrm{DG}}(d)).
\]
\end{definition}

\begin{lemma}[\leanname{C\_DeGiorgi\_subsolution\_le\_normalized}]
For a normalised elliptic coefficient field $A$ on $\Bz{1}$ (i.e.\
$\lambda = 1$, so $\Lambda \ge 1$ and the ratio
$\Lambda/\lambda = \Lambda$),
\[
  C_{\mathrm{DG}}(d, A) \le
    C^{\mathrm{norm}}_{\mathrm{DG}}(d) \cdot \Lambda^{d/4}.
\]
\end{lemma}

\begin{proof}
Comparing the formulas, $K_{\mathrm{DG}}(d,A) = c (\Lambda+1) + 1$ and
$K^{\mathrm{norm}}_{\mathrm{DG}}(d) = 2c + 1$ where $c =
(C_{\mathrm{gns}}(d,2))^2 \cdot 8 (2 M_{\mathrm{st}})^2$. Since
$\Lambda \ge 1$ one has $c(\Lambda + 1) + 1 \le (2c+1)\Lambda$,
i.e.\ $K_{\mathrm{DG}}(d,A) \le K^{\mathrm{norm}}_{\mathrm{DG}}(d) \cdot
\Lambda$. Monotonicity of $C_{\mathrm{small}}$ in $K$ followed by the
multiplicative scaling identity
$C_{\mathrm{small}}(d, K \cdot \Lambda) = C_{\mathrm{small}}(d, K)\,
\Lambda^{d/4}$ gives the claim.
\end{proof}

\begin{theorem}[Normalised De Giorgi $L^\infty$ bound,
\leanname{linfty\_subsolution\_DeGiorgi\_normalized}]
Let $d > 2$, $A$ a normalised elliptic coefficient field on
$\Bz{1}$, and $u$ a subsolution with
$|\max(u,0)|^2 \in L^1(\Bz{1})$. Then
\[
  \max(u(x), 0) \le
    C^{\mathrm{norm}}_{\mathrm{DG}}(d) \cdot \Lambda^{d/4} \cdot
    \Bigl(\int_{\Bz{1}} |\max(u,0)|^2 \, dx\Bigr)^{1/2}
\]
for almost every $x \in \Ball{1/2}{0}$.
\end{theorem}

\begin{proof}
Combine \leanname{linfty\_subsolution\_DeGiorgi} with the previous
lemma comparing $C_{\mathrm{DG}}(d,A)$ to
$C^{\mathrm{norm}}_{\mathrm{DG}}(d) \, \Lambda^{d/4}$, multiplying
by the nonnegative factor $\sqrt{I_0}$.
\end{proof}

This concludes the formal De Giorgi $L^2 \to L^\infty$ chapter. The
two top-level interfaces \leanname{DeGiorgiIteration} and
\leanname{DeGiorgiTheory} simply re-export the constructions and
theorems above so that downstream chapters (weak Harnack, Harnack,
Holder) can rely on the De Giorgi bound directly.


%% ========================================================================
\section{Moser Iteration and Supersolution Estimates}\label{sec:moser}
%% ========================================================================

%% --- Section 7 ---

\subsection{Moser Iteration}

Throughout this section we fix a dimension $d \in \N$ with $d \geq 1$ and
$d > 2$, and write $E = \R^d$ for the ambient space. All balls are taken
around the origin in $E$, and we abbreviate $\Bz{r} = \Ball{r}{0}$. Functions
$u : E \to \R$ are real valued; $A$ denotes a normalized elliptic coefficient
field on $\Bz{1}$, with ellipticity ratio $\Lambda = \Lambda_A \geq 1$
\leanname{NormalizedEllipticCoeff}. The notion of subsolution is
\leanname{IsSubsolution} and that of supersolution is
\leanname{IsSupersolution}. The starting point is the normalized De Giorgi
$L^2 \to L^\infty$ estimate proved in Section 5.

\subsubsection{Moser constants}

\begin{definition}[Moser anchor constant, \leanname{C\_MoserAnchor}]
The \emph{Moser anchor} in dimension $d$ is
\[
  C_{\mathrm{MA}}(d) := \max\!\left\{1,\;
    C_{\mathrm{DG}}^{\mathrm{sub}}(d)^{2},\;
    8\,\Mst^{2}\,C_{\mathrm{gns}}(d,2)^{2}\right\},
\]
where $C_{\mathrm{DG}}^{\mathrm{sub}}(d)$ is the normalized De Giorgi
subsolution constant of Section 5 and $C_{\mathrm{gns}}(d,2)$ is the
Gagliardo--Nirenberg--Sobolev constant for $p=2$.
\end{definition}

\begin{lemma}[\leanname{one\_le\_C\_MoserAnchor}]
$1 \leq C_{\mathrm{MA}}(d)$.
\end{lemma}

\begin{proof}
Immediate from $C_{\mathrm{MA}}(d) \geq \max\{1,\dots\} \geq 1$.
\end{proof}

\begin{definition}[Sobolev gain factor and decay ratio]
Set
\[
  \chi := \frac{d}{d-2} \quad (\text{\leanname{moserChi}}),
  \qquad
  q := \frac{d-2}{d} \quad (\text{\leanname{moserDecayRatio}}).
\]
\end{definition}

\begin{lemma}[Basic algebraic facts on $\chi,q$]
Assume $d > 2$. Then
\begin{enumerate}
  \item $0 \leq q$ \leanname{moserDecayRatio\_nonneg};
  \item $q < 1$ \leanname{moserDecayRatio\_lt\_one};
  \item $0 < \chi$ \leanname{moserChi\_pos};
  \item $1 < \chi$ \leanname{one\_lt\_moserChi};
  \item $q = \chi^{-1}$ \leanname{moserDecayRatio\_eq\_inv\_moserChi}.
\end{enumerate}
\end{lemma}

\begin{proof}
For (1)--(2) note $0 \leq d-2 < d$, so $0 \leq (d-2)/d < 1$.
For (3)--(4), $d/(d-2) > 0$ and $d > d-2$ gives $d/(d-2) > 1$.
For (5), $\big((d-2)/d\big)^{-1} = d/(d-2)$.
\end{proof}

\begin{definition}[Dimensional Moser constant, \leanname{C\_Moser}]
Set $C_{\mathrm{MA}} := C_{\mathrm{MA}}(d)$. If $d > 2$, define
\[
  C_{\mathrm{Moser}}(d) := \max\!\left\{C_{\mathrm{MA}},\;
    (32\,C_{\mathrm{MA}})^{d/2}\,
    4^{\sum_{n=0}^{\infty} n\,q^{n}}\right\};
\]
otherwise $C_{\mathrm{Moser}}(d) := C_{\mathrm{MA}}$.
\end{definition}

\begin{lemma}[\leanname{C\_MoserAnchor\_le\_C\_Moser}, \leanname{one\_le\_C\_Moser}, \leanname{C\_DeGiorgi\_subsolution\_normalized\_sq\_le\_C\_Moser}]
$C_{\mathrm{MA}}(d) \leq C_{\mathrm{Moser}}(d)$, $1 \leq C_{\mathrm{Moser}}(d)$, and
$C_{\mathrm{DG}}^{\mathrm{sub}}(d)^2 \leq C_{\mathrm{Moser}}(d)$.
\end{lemma}

\begin{proof}
Both branches of the definition return a value $\geq C_{\mathrm{MA}}(d) \geq 1$.
The third statement follows by transitivity from the corresponding bound for
$C_{\mathrm{MA}}$.
\end{proof}

\subsubsection{Geometric sequences}

\begin{definition}[Nested Moser radii, \leanname{moserRadius}]
\[
  r_n := \frac{1 + 2^{-n}}{2}, \qquad n \geq 0.
\]
Thus $r_0 = 1$ and $r_n \searrow 1/2$.
\end{definition}

\begin{lemma}[Properties of the radii]
For all $n \in \N$:
\begin{enumerate}
  \item $r_0 = 1$ \leanname{moserRadius\_zero};
  \item $r_n \leq 1$ \leanname{moserRadius\_le\_one};
  \item $0 < r_n$ \leanname{moserRadius\_pos};
  \item $1/2 < r_n$ \leanname{moserRadius\_gt\_half};
  \item $r_n - r_{n+1} = 2^{-(n+2)}$ \leanname{moserRadius\_gap};
  \item $r_{n+1} < r_n$ \leanname{moserRadius\_succ\_lt}.
\end{enumerate}
\end{lemma}

\begin{proof}
Each follows by direct algebra: e.g.\ for (5),
$r_n - r_{n+1} = (2^{-n} - 2^{-(n+1)})/2 = 2^{-(n+2)}$.
\end{proof}

\begin{definition}[Exponent sequence, \leanname{moserExponentSeq}]
Given $p_0 \in \R$ and $n \in \N$, set
\[
  p_n := p_0 \chi^{n}.
\]
\end{definition}

\begin{lemma}
With $\chi = d/(d-2)$ and $q = \chi^{-1}$:
\begin{enumerate}
  \item $p_0 = p_0$ and $p_{n+1} = \chi p_n$ \leanname{moserExponentSeq\_zero}, \leanname{moserExponentSeq\_succ};
  \item if $d>2$ and $p_0 > 0$, then $p_n > 0$ \leanname{moserExponentSeq\_pos};
  \item if $d>2$ and $p_0 \geq 0$, then $p_0 \leq p_n$ \leanname{moserExponentSeq\_ge\_initial};
  \item if $d>2$ and $p_0 > 1$, then $p_n > 1$ \leanname{one\_lt\_moserExponentSeq};
  \item if $d>2$ and $p_0 > 0$, then $1/p_n = q^n / p_0$ \leanname{inv\_moserExponentSeq}.
\end{enumerate}
\end{lemma}

\begin{proof}
All four are immediate consequences of $\chi > 1$ and $p_n = p_0 \chi^n$;
identity (5) uses $\chi^{-n} = q^n$.
\end{proof}

\subsubsection{Regularized power cutoffs}

For the cutoff--power test functions of the iteration we set
\begin{equation*}
  \mathrm{PC}_d(\eta,u,p)(x) := \eta(x)\,|\!\max(u(x),0)|^{p/2}
  \quad (\text{\leanname{moserPowerCutoff}}),
\end{equation*}
the smooth cutoff $\eta$ being supported in a slightly larger ball than the
inner ball where the bound is sought.

\begin{lemma}[Support property, \leanname{moserPowerCutoff\_tsupport\_subset}]
If $\supp \eta \subseteq \Omega$, then
$\supp \mathrm{PC}_d(\eta,u,p) \subseteq \Omega$.
\end{lemma}

\begin{proof}
Pointwise $\mathrm{PC}_d(\eta,u,p)(x) = 0$ whenever $\eta(x)=0$.
\end{proof}

\subsubsection{Sobolev preparation and the energy bound}

The hard nonlinear core of the Moser iteration is the construction of a
$W^{1,2}$ witness for the cutoff power $\mathrm{PC}_d(\eta,u,p)$ together with
the matching energy estimate.

\begin{theorem}[Cutoff--power energy bound, \leanname{moser\_powerCutoff\_memW1p\_energy\_of\_subsolution\_core}]
\label{thm:moser-cutoff-energy}
Let $d > 2$, $A$ a normalized elliptic coefficient on $\Bz{1}$, $u$ a
subsolution of $A$, and $\eta \in C^{\infty}_c(\Bz{s})$ a smooth cutoff with
$0 \leq \eta \leq 1$ and $\|\nabla \eta\|_{L^\infty} \leq C_\eta$. Suppose
$0 < s \leq 1$, $1 < p$, and $\int_{\Bz{s}} |u_+|^p < \infty$. Then there
exists a $W^{1,2}$ witness for $v := \mathrm{PC}_d(\eta,u,p)$ on $\Bz{s}$
satisfying
\begin{equation*}
  \int_{\Bz{s}} \|\nabla v\|^{2}\,dx
  \;\leq\;
  2\,C_\eta^{2}\,\Big(\Lambda\,\Big(\frac{p}{p-1}\Big)^{2} + 1\Big)
  \int_{\Bz{s}} |u_+|^{p}\,dx.
\end{equation*}
\end{theorem}

\begin{proof}
The witness for $v := \mathrm{PC}_d(\eta, u, p) = \eta\,(u_+)^{p/2}$ is
constructed in \leanname{moserPowerCutoff\_memW1pWitness}, which we now
unpack. First, a $W^{1,2}$ witness for $(u_+)^{p/2}$ on $\Bz{s}$ is obtained
from the truncation $u_+$ (whose witness comes from
\leanname{positivePartSub\_memW1pWitness\_on\_ball}) composed with the
power $t \mapsto t^{p/2}$ via the $W^{1,2}$ chain rule
\leanname{sobolev\_chain\_rule\_unitBall}; the local integrability
$\int (u_+)^{p} < \infty$ together with $p > 1$ ensures the chain-rule
hypotheses (locally Lipschitz factor on $[0,\infty)$). The product
$\eta\,(u_+)^{p/2}$ is then an honest $W^{1,2}$ function via
\leanname{MemW1pWitness.mul\_smooth\_bounded\_p}, with weak gradient
$\nabla v = \eta\,(p/2)(u_+)^{p/2-1}\nabla u_+ + (u_+)^{p/2}\nabla \eta$.

For the energy bound, test the subsolution inequality against
$\varphi := \eta^{2}(u_+)^{p-1}$ (this is in $H^{1}_{0}(\Bz{s})$ thanks to
the support inclusion of $\eta$ and the integrability hypothesis on $u_+$).
The pointwise expansion of $\langle A\nabla u, \nabla \varphi\rangle$ via
\leanname{bilinFormIntegrand\_mul\_smooth\_eq} yields
\begin{multline*}
  \int_{\Bz{s}} (p-1)\,\eta^{2}\,(u_+)^{p-2}\,\langle A\nabla u_+, \nabla u_+\rangle\,dx
  \\
  \le -2\int_{\Bz{s}} \eta\,(u_+)^{p-1}\,\langle A\nabla u_+, \nabla \eta\rangle\,dx,
\end{multline*}
where on the left we discarded the supersolution-favourable term and on the
right we used the subsolution inequality. The right-hand side is bounded by
Cauchy--Schwarz against $\langle A\cdot, \cdot\rangle$ followed by AM--GM
with weight $\varepsilon = (p-1)/2$:
\[
  2|\eta\,(u_+)^{p-1}\langle A\nabla u_+,\nabla \eta\rangle|
  \le \frac{p-1}{2}\eta^{2}(u_+)^{p-2}\langle A\nabla u_+,\nabla u_+\rangle
  + \frac{2}{p-1}(u_+)^{p}\langle A\nabla \eta,\nabla \eta\rangle.
\]
Absorbing the first term into the left-hand side leaves
$(p-1)/2 \cdot \int \eta^{2}(u_+)^{p-2}\langle A\nabla u_+,\nabla u_+\rangle$,
and the residual term is bounded by $\Lambda C_\eta^{2}\int (u_+)^{p}$ via
the upper ellipticity bound for $A$ and $\|\nabla \eta\| \le C_\eta$.

Finally, expanding $\nabla v = \eta\,(p/2)(u_+)^{p/2-1}\nabla u_+ +
(u_+)^{p/2}\nabla \eta$, the pointwise inequality
$\|\nabla v\|^{2} \le 2\,\eta^{2}(p/2)^{2}(u_+)^{p-2}\|\nabla u_+\|^{2}
+ 2(u_+)^{p}\|\nabla \eta\|^{2}$ together with the absorbed Caccioppoli
estimate gives
\[
  \int_{\Bz{s}} \|\nabla v\|^{2}\,dx
   \le 2\,C_\eta^{2}\Bigl(\Lambda\bigl(\tfrac{p}{p-1}\bigr)^{2} + 1\Bigr)
   \int_{\Bz{s}} (u_+)^{p}\,dx,
\]
the displayed energy inequality.
\end{proof}

\begin{theorem}[Zero trace upgrade, \leanname{moser\_powerCutoff\_memW01p\_energy\_of\_subsolution}]
Under the hypotheses of Theorem~\ref{thm:moser-cutoff-energy}, the witness
moreover satisfies $v \in W^{1,2}_{0}(\Bz{s})$.
\end{theorem}

\begin{proof}
Since $\supp \eta \subseteq \Bz{s}$ is compact, $\supp v \subseteq \Bz{s}$ is
likewise compactly contained. The criterion
\leanname{memW01p\_of\_memW1p\_of\_tsupport\_subset} then upgrades $v$ from
$W^{1,p}$ to $W^{1,p}_0$.
\end{proof}

\subsubsection{One-step Moser pre-estimate}

\begin{theorem}[Concentric ball pre-estimate, \leanname{moser\_preMoser}]
\label{thm:moser-preMoser}
Let $d > 2$, $A$ as above, $u$ a subsolution, and $0 < r < s \leq 1$.
Suppose $1 < p$ and $\int_{\Bz{s}} |u_+|^p < \infty$. Then $|u_+|^{\chi p}$
is integrable on $\Bz{r}$ and
\begin{multline*}
  \Big(\int_{\Bz{r}} |u_+|^{\chi p}\,dx\Big)^{1/(\chi p)}
  \\ \leq\;
  \Big(\frac{C_{\mathrm{MA}}(d)}{(s-r)^{2}}\,
    \Big(\Lambda\Big(\frac{p}{p-1}\Big)^{2} + 1\Big)\Big)^{1/p}
  \,\Big(\int_{\Bz{s}} |u_+|^{p}\,dx\Big)^{1/p}.
\end{multline*}
\end{theorem}

\begin{proof}
Choose the midpoint radius $R := (r+s)/2 \in (r,s)$ and a smooth radial cutoff
$\eta$ with $\eta = 1$ on $\Bz{r}$, $\supp \eta \subseteq \overline{\Bz{R}}
\subseteq \Bz{s}$, $0 \leq \eta \leq 1$, and (by
\leanname{Cutoff.exists}) $\|\nabla \eta\|_{L^\infty} \leq C_\eta := 2\Mst/(s-r)$.

Set $v := \mathrm{PC}_d(\eta,u,p) = \eta\,(u_+)^{p/2}$. By
Theorem~\ref{thm:moser-cutoff-energy} we get $v \in W^{1,2}_0(\Bz{s})$ with
\begin{equation*}
  \int_{\Bz{s}} \|\nabla v\|^{2}\,dx
  \;\leq\;
  2\,C_\eta^{2}\,\Big(\Lambda\Big(\frac{p}{p-1}\Big)^{2}+1\Big)
  \int_{\Bz{s}} (u_+)^{p}\,dx.
\end{equation*}
The Sobolev inequality \leanname{sobolev\_prepare\_on\_ball} on $\Bz{s}$ in
exponent $2 \to 2^{*} = 2d/(d-2)$ gives, with $C_{\mathrm{gns}} = C_{\mathrm{gns}}(d,2)$,
\begin{equation*}
  \Big(\int_{\Bz{s}} v^{2^{*}}\,dx\Big)^{1/2^{*}}
  \;\leq\;
  C_{\mathrm{gns}}\,
  \Big(\int_{\Bz{s}} \|\nabla v\|^{2}\,dx\Big)^{1/2}.
\end{equation*}
On $\Bz{r}$ we have $\eta = 1$, so $v^{2^{*}} = (u_+)^{p\,d/(d-2)} = (u_+)^{\chi p}$.
Combining with the energy bound and using $C_\eta^2 = 4\Mst^2/(s-r)^2$ together
with $8\Mst^{2}\,C_{\mathrm{gns}}^{2} \leq C_{\mathrm{MA}}(d)$
(\leanname{cutoff\_sobolev\_anchor\_le\_C\_MoserAnchor}) yields
\begin{equation*}
  \Big(\int_{\Bz{r}} (u_+)^{\chi p}\,dx\Big)^{2/2^{*}}
  \;\leq\;
  \frac{C_{\mathrm{MA}}(d)}{(s-r)^{2}}
  \Big(\Lambda\Big(\frac{p}{p-1}\Big)^{2}+1\Big)
  \int_{\Bz{s}} (u_+)^{p}\,dx.
\end{equation*}
Raising to the power $1/p$ and using $2/(2^{*}\,p)= 1/(\chi p)$ gives the
claim.
\end{proof}

\subsubsection{Iteration}

We now iterate the pre-estimate along $(r_n,p_n)$.

\begin{definition}[Iteration quantities]
Set
\begin{align*}
  I_n(u) &:= \int_{\Bz{r_n}} (u_+)^{p_n}\,dx
    && (\text{\leanname{moserIterIntegral}}), \\
  N_n(u) &:= I_n(u)^{1/p_n}
    && (\text{\leanname{moserIterNorm}}), \\
  \mathrm{S}(A,p_0) &:= 32\,C_{\mathrm{MA}}(d)\,\Lambda\,\Big(\frac{p_0}{p_0-1}\Big)^{2}
    && (\text{\leanname{moserStepConst}}).
\end{align*}
The product majorant produced by iteration is, for $n \geq 0$,
\begin{equation*}
  M_n(A,u,p_0) :=
  \Big(
    \mathrm{S}(A,p_0)^{\sum_{i=0}^{n-1} q^{i}}\,
    4^{\sum_{i=0}^{n-1} i\,q^{i}}\,
    \int_{\Bz{1}} (u_+)^{p_0}\,dx
  \Big)^{1/p_0}.
\end{equation*}
We further define the target majorant
\begin{align*}
  B_\infty(A,u,p_0) &:= C_{\mathrm{Moser}}(d)\,\Lambda^{d/2}\,
    \Big(\frac{p_0}{p_0-1}\Big)^{d}\,\int_{\Bz{1}} (u_+)^{p_0}\,dx
    && (\text{\leanname{moserLinftyBoundPow}}), \\
  M_\infty(A,u,p_0) &:= B_\infty(A,u,p_0)^{1/p_0}
    && (\text{\leanname{moserLinftyMajorant}}).
\end{align*}
\end{definition}

\begin{theorem}[Step--majorant inequality, \leanname{moser\_step\_majorant\_le}]
\label{thm:moser-step-maj}
For $d > 2$, $1 < p_0$ and any $n \geq 0$,
\begin{multline*}
  \Big(\frac{C_{\mathrm{MA}}(d)}{(r_n-r_{n+1})^{2}}\,
    \Big(\Lambda\Big(\frac{p_n}{p_n-1}\Big)^{2}+1\Big)\Big)^{1/p_n}
  M_n(A,u,p_0) \\
  \leq\;
  M_{n+1}(A,u,p_0).
\end{multline*}
\end{theorem}

\begin{proof}
First, $r_n - r_{n+1} = 2^{-(n+2)}$ so
$C_{\mathrm{MA}}/(r_n-r_{n+1})^2 = C_{\mathrm{MA}}\,4^{n+2} = 16\,C_{\mathrm{MA}}\,4^{n}$.
Second, $p_n/(p_n-1) \leq p_0/(p_0-1)$ since $t \mapsto t/(t-1)$ is decreasing
on $(1,\infty)$, hence
\[
  \Lambda\Big(\frac{p_n}{p_n-1}\Big)^{2}+1
  \;\leq\; 2\,\Lambda\Big(\frac{p_0}{p_0-1}\Big)^{2}.
\]
Combining gives
\begin{equation*}
  \frac{C_{\mathrm{MA}}}{(r_n-r_{n+1})^{2}}\,
    \Big(\Lambda\Big(\frac{p_n}{p_n-1}\Big)^{2}+1\Big)
  \;\leq\;
  \mathrm{S}(A,p_0)\,4^{n}.
\end{equation*}
Raising to $1/p_n = q^n/p_0$ and multiplying by $M_n$ produces the inner factor
$\mathrm{S}^{q^n/p_0} 4^{n q^n/p_0}$, which precisely promotes the partial sums
in the exponents from $\sum_{i<n}$ to $\sum_{i\leq n}$. The resulting expression
equals $M_{n+1}(A,u,p_0)$ by the definition of $M_{n+1}$.
\end{proof}

\begin{theorem}[Iteration bound, \leanname{moser\_iteration\_bound}]
\label{thm:moser-iter-bound}
For $d > 2$, $1 < p_0$, $u$ a subsolution of $A$, and
$\int_{\Bz{1}} |u_+|^{p_0} < \infty$, for every $n \in \N$ we have
$|u_+|^{p_n}$ integrable on $\Bz{r_n}$ and
\[
  N_n(u) \;\leq\; M_n(A,u,p_0).
\]
\end{theorem}

\begin{proof}
By induction on $n$. The case $n=0$ uses $r_0=1$, $p_0 = p_0$, and
$M_0(A,u,p_0)= (\int_{\Bz{1}} (u_+)^{p_0})^{1/p_0}$ by definition. For the
inductive step, the pre-estimate Theorem~\ref{thm:moser-preMoser} applied with
$(p,r,s) = (p_n,r_{n+1},r_n)$ gives integrability of $|u_+|^{\chi p_n}$ on
$\Bz{r_{n+1}}$ and
\begin{multline*}
  N_{n+1}(u)
  \;\leq\;
  \Big(\frac{C_{\mathrm{MA}}}{(r_n-r_{n+1})^{2}}
    \Big(\Lambda\Big(\frac{p_n}{p_n-1}\Big)^{2}+1\Big)\Big)^{1/p_n}
  N_n(u).
\end{multline*}
The inductive hypothesis $N_n(u) \leq M_n(A,u,p_0)$ and
Theorem~\ref{thm:moser-step-maj} then yield $N_{n+1}(u) \leq M_{n+1}(A,u,p_0)$.
\end{proof}

\begin{lemma}[Geometric majorant absorption, \leanname{moser\_geometric\_majorant}]
\label{lem:moser-geom-maj}
For $d > 2$, $1 < p_0$, and all $n$,
\[
  M_n(A,u,p_0) \;\leq\; M_\infty(A,u,p_0).
\]
\end{lemma}

\begin{proof}
Set $S := \sum_{i=0}^{n-1} q^{i}$, $T := \sum_{i=0}^{n-1} i q^{i}$,
$B := 32\,C_{\mathrm{MA}}(d)$, $r := p_0/(p_0-1)$. Then:
$S \leq d/2$ \leanname{moserDecayRatio\_sum\_le\_half\_dim}, since
$\sum_{i\geq 0} q^{i} = 1/(1-q) = d/2$; also $T \leq \sum_{i\geq 0} i\,q^{i}$
\leanname{moserDecayRatio\_weighted\_sum\_le\_tsum} by summability.
Now $\mathrm{S}(A,p_0) = B \Lambda r^{2}$, so
\[
  \mathrm{S}(A,p_0)^{S}
  = B^{S} \Lambda^{S} (r^{2})^{S}.
\]
Using $1 \leq B$, $1 \leq \Lambda$, and $1 \leq r$, monotonicity of $t \mapsto
t^{e}$ in the exponent gives $B^{S} \leq B^{d/2}$, $\Lambda^{S} \leq
\Lambda^{d/2}$, and $(r^{2})^{S} \leq r^{d}$ (because $2S \leq d$). Also $4^{T}
\leq 4^{\sum_{i\geq 0} i q^{i}}$. Multiplying these inequalities and using the
explicit definition of $C_{\mathrm{Moser}}(d)$ in the case $d>2$:
\[
  B^{d/2}\,4^{\sum_{i\geq 0} i q^{i}} \;\leq\; C_{\mathrm{Moser}}(d),
\]
we obtain
\[
  \mathrm{S}(A,p_0)^{S}\,4^{T}\,I_0(u)
  \;\leq\;
  C_{\mathrm{Moser}}(d)\,\Lambda^{d/2}\,r^{d}\,I_0(u)
  = B_\infty(A,u,p_0).
\]
Raising to the power $1/p_0$ yields $M_n \leq M_\infty$.
\end{proof}

\subsubsection{Closeout to $L^\infty$}

\begin{lemma}[Closeout from iterated bounds, \leanname{moser\_ae\_closeout}]
\label{lem:moser-closeout}
Fix $d > 2$, $1 < p_0$, and assume that for every $n \in \N$,
$|u_+|^{p_n}$ is integrable on $\Bz{r_n}$ and $N_n(u) \leq M_\infty(A,u,p_0)$.
Then for almost every $x \in \Bz{1/2}$,
\[
  |u_+(x)|^{p_0} \;\leq\; B_\infty(A,u,p_0).
\]
\end{lemma}

\begin{proof}
Set $\mu := \mathrm{vol}\,|\,\Bz{1/2}$, $f(x) := |u_+(x)|$, $g(x) := f(x)^{p_0/2}$,
and $K := B_\infty(A,u,p_0)^{1/2}$. Note $K \geq 0$.
For each $n$, write $q_n := 2 \chi^n$. We claim that for every $n$
and every $c > K$,
\[
  \mu_{\mathrm{r}}\{x : c \leq g(x)\} \;\leq\; (K/c)^{q_n}.
\]
Indeed, on $\{c \leq g\}$ we have $c^{q_n} \leq g^{q_n}$, and Markov's
inequality gives
\[
  c^{q_n}\,\mu_{\mathrm{r}}\{c \leq g\}
  \;\leq\;
  \int g^{q_n}\,d\mu
  = \int_{\Bz{1/2}} f^{p_n}\,dx
  \;\leq\;
  \int_{\Bz{r_n}} f^{p_n}\,dx = I_n(u),
\]
where we used $g^{q_n} = f^{p_0 q_n / 2} = f^{p_n}$ and
$\Bz{1/2} \subseteq \Bz{r_n}$. By hypothesis $I_n(u)^{1/p_n} = N_n(u) \leq
M_\infty$, so raising to $p_n$ and using $1/q_n = (1/p_n)(p_0/2)$ we get
$I_n(u) \leq M_\infty^{q_n p_0/2 \cdot 2/p_0} = M_\infty^{q_n p_0/p_0}$\,;
more cleanly, $I_n^{1/q_n} \leq M_\infty^{p_0/2} = K$. Hence
$\mu_{\mathrm{r}}\{c\leq g\} \leq (K/c)^{q_n}$.

As $\chi > 1$, $q_n \to \infty$, so $(K/c)^{q_n} \to 0$ for every $c > K$.
Therefore $\mu_{\mathrm{r}}\{c \leq g\} = 0$ for all $c > K$, and a countable
union over $c = K + 1/(m+1)$ shows $g \leq K$ a.e.\ on $\Bz{1/2}$. Squaring,
$f^{p_0} = g^{2} \leq K^{2} = B_\infty$.
\end{proof}

\begin{theorem}[Moser $L^p \to L^\infty$ for subsolutions, \leanname{linfty\_subsolution\_Moser}]
\label{thm:linfty-Moser}
Let $d > 2$, $A$ a normalized elliptic coefficient on $\Bz{1}$, $u$ a
subsolution, and $1 < p_0$ with $\int_{\Bz{1}} |u_+|^{p_0} < \infty$. Then for
almost every $x \in \Bz{1/2}$,
\[
  |u_+(x)|^{p_0}
  \;\leq\;
  C_{\mathrm{Moser}}(d)\,\Lambda^{d/2}\,
  \Big(\frac{p_0}{p_0-1}\Big)^{d}
  \int_{\Bz{1}} |u_+|^{p_0}\,dx.
\]
\end{theorem}

\begin{proof}
The headline Moser bound is the composition of three pieces, all assembled
along the discrete sequence of radii $r_n$ and exponents $p_n$.

\emph{Step 1: iterated step bound.} Theorem~\ref{thm:moser-iter-bound}
(\leanname{moser\_iteration\_bound}) is applied as a single induction on
$n$: at each step, the concentric-ball pre-estimate
Theorem~\ref{thm:moser-preMoser} on $(r_{n+1}, r_n)$ with exponent $p_n$
upgrades $L^{p_n}$ control of $u_+$ on $\Bz{r_n}$ to $L^{p_{n+1}} = L^{\chi p_n}$
control on $\Bz{r_{n+1}}$, with multiplicative constant the step constant
$\mathrm{S}_n$ depending on $\Lambda$, $C_{\mathrm{MA}}(d)$, the dimension,
and the gap $r_n - r_{n+1} = 2^{-(n+2)}$. Telescoping over $i < n$ produces
the bound
\[
  N_n(u)
   \le \Bigl(\prod_{i=0}^{n-1} \mathrm{S}_i\Bigr) \cdot N_0(u)
   = M_n(A, u, p_0).
\]
At every stage the integrability of $|u_+|^{p_n}$ on $\Bz{r_n}$ is propagated
forward by the same induction, providing the second half of the conclusion.

\emph{Step 2: geometric absorption of the product.}
Lemma~\ref{lem:moser-geom-maj} (\leanname{moser\_geometric\_majorant})
bounds the infinite product $\prod_{i \ge 0} \mathrm{S}_i$ by the closed-form
constant $M_\infty(A,u,p_0)$. The key two summation identities are
$\sum_{i \ge 0} q^{i} \le d/2$
(\leanname{moserDecayRatio\_sum\_le\_half\_dim}) and the weighted variant
$\sum_{i \ge 0} i\,q^{i}$
(\leanname{moserDecayRatio\_weighted\_sum\_le\_tsum}); these are precisely
what is needed to convert the product of $\mathrm{S}_i$, after taking
logarithms and using $1/p_i = q^{i}/p_0$, into a finite expression involving
$C_{\mathrm{Moser}}(d)$, $\Lambda^{d/2}$ and the explicit
$\bigl(p_0/(p_0-1)\bigr)^{d}$ factor visible in the conclusion. Composing with
Step 1 gives $N_n(u) \le M_\infty(A, u, p_0)$ for every $n$.

\emph{Step 3: a.e.\ closeout via Markov.}
Lemma~\ref{lem:moser-closeout} (\leanname{moser\_ae\_closeout}) converts the
uniform $L^{p_n}$ bound at every stage into an a.e.\ pointwise bound on the
intersection $\Bz{1/2} \subseteq \bigcap_n \Bz{r_n}$. The argument is a
Markov-type estimate: for every $c > K := M_\infty(A,u,p_0)^{1/2}$, the
super-level set $\{c \le f\}$ on $\Bz{1/2}$ has measure at most
$(K/c)^{q_n}$, and since $q_n \to \infty$ this measure is zero. Taking a
countable union over $c = K + 1/(m+1)$ shows $f \le K$ a.e.\ on $\Bz{1/2}$,
and squaring gives $|u_+|^{p_0} \le K^{2}$ a.e.\ on $\Bz{1/2}$, which is the
displayed Moser bound.
\end{proof}

\begin{remark}[$p_0 = 2$ anchor, \leanname{linfty\_subsolution\_Moser\_two}]
The case $p_0 = 2$ already follows from the normalized De Giorgi $L^2 \to
L^\infty$ estimate of Section 5: with $c := C_{\mathrm{DG}}^{\mathrm{sub}}(d)
\Lambda^{d/4}$, the De Giorgi bound gives
$\max(u(x),0) \leq c\,(\int_{\Bz{1}} |u_+|^{2}\,dx)^{1/2}$ a.e.\ on
$\Bz{1/2}$. Squaring and using
$C_{\mathrm{DG}}^{\mathrm{sub}}(d)^{2} \leq C_{\mathrm{Moser}}(d)$ together with
$2/(2-1) = 2$ recovers the $p_0=2$ case of Theorem~\ref{thm:linfty-Moser}.
\end{remark}

%% --- Section 8 ---

\subsection{Supersolution Estimates}

We now turn to positive supersolutions of the same operator. The main outputs
of this section are the two stages of the weak Harnack inequality:
\emph{stage one (inverse)} produces an a.e.\ lower bound on $u$ from inverse
$L^{p}$ control of $u$, and \emph{stage two (forward)} produces a low-power
$L^{p}$-on-$L^{p}$ Caccioppoli-type bound. The architecture mirrors Section 7,
with the role of $u_+$ replaced by $u^{-1}$ (resp.\ $u$ at low powers).

Throughout this section, $u : E \to \R$ is assumed to be strictly positive on
$\Bz{1}$ and a supersolution of $A$. We write $\mu_{1/2}$ for
$\mathrm{vol}\,|\,\Bz{1/2}$.

\subsubsection{The weighted Caccioppoli inequality}

The abstract energy tool driving every iteration in this section is the
following weighted Caccioppoli inequality.

\begin{theorem}[Weighted Caccioppoli, \leanname{weighted\_caccioppoli\_of\_supersolution}]
\label{thm:weighted-cacc}
Let $\Omega \subseteq E$ be open, $A$ a normalized elliptic coefficient on
$\Omega$, $u$ a supersolution with $u \in W^{1,2}(\Omega)$, and let
$\Psi : \R \to \R$ be smooth with $\Psi(0)=0$, $|\Psi'| \leq M$, $\Psi(u) \geq 0$
on $\Omega$, and $\Psi'(u) \leq 0$ on $\Omega$ (with $\Psi(u)=0$ wherever
$\Psi'(u)=0$). Let $\varphi \in C^{\infty}_c(\Omega)$ satisfy $|\varphi| \leq 1$
and $\|\nabla \varphi\|_{L^\infty} \leq C_\varphi$, with the integrability
\[
  \int_{\Omega} \|\nabla \varphi\|^{2}\,
  \frac{\Psi(u)^{2}}{-\Psi'(u)}\,dx < \infty
\]
(interpreted as $0$ where $\Psi'(u)=0$). Then
\begin{multline*}
  \int_{\Omega} \varphi^{2}\,(-\Psi'(u))\,\langle A\nabla u,\nabla u\rangle\,dx
  \\ \leq\;
  4\,\Lambda \int_{\Omega} \|\nabla \varphi\|^{2}\,
  \frac{\Psi(u)^{2}}{-\Psi'(u)}\,dx.
\end{multline*}
\end{theorem}

\begin{proof}
By the supersolution property tested against $\varphi^{2}\Psi(u) \geq 0$,
\[
  0 \leq \int_{\Omega} \langle A\nabla u, \nabla(\varphi^{2}\Psi(u))\rangle\,dx.
\]
Using the chain rule for $W^{1,2}$ functions, $\nabla(\Psi(u)) = \Psi'(u)\nabla u$,
and the product rule $\nabla(\varphi^{2}\Psi(u)) = \varphi^{2}\Psi'(u)\nabla u
+ 2\varphi\Psi(u)\nabla \varphi$. The bilinear form expansion
(\leanname{bilinFormIntegrand\_mul\_smooth\_eq}) yields the pointwise identity
\[
  \langle A\nabla u, \nabla(\varphi^{2}\Psi(u))\rangle
  = \varphi^{2}\Psi'(u)\,\langle A\nabla u,\nabla u\rangle
  + 2\varphi\Psi(u)\,\langle A\nabla u,\nabla \varphi\rangle.
\]
Since $\Psi'(u) \leq 0$, the principal term is non-positive; rearranging
\[
  \int \varphi^{2}\,(-\Psi'(u))\,\langle A\nabla u,\nabla u\rangle\,dx
  \;\leq\;
  2 \int \varphi\,\Psi(u)\,\langle A\nabla u,\nabla \varphi\rangle\,dx.
\]
By Cauchy--Schwarz with respect to the symmetric form $\langle A\cdot,\cdot\rangle$
and AM--GM with weight $\varepsilon = 1/2$,
\begin{multline*}
  2|\varphi\Psi(u)\,\langle A\nabla u,\nabla \varphi\rangle|
  \\ \leq\;
  \tfrac12\,\varphi^{2}(-\Psi'(u))\langle A\nabla u,\nabla u\rangle
  + 2\,\frac{\Psi(u)^{2}}{-\Psi'(u)}\,\langle A\nabla \varphi,\nabla \varphi\rangle.
\end{multline*}
Absorbing the first term into the left-hand side and bounding
$\langle A\nabla \varphi,\nabla \varphi\rangle \leq \Lambda \|\nabla \varphi\|^{2}$
(ellipticity \leanname{NormalizedEllipticCoeff.upper}) yields the stated inequality
with constant $4\Lambda$.
\end{proof}

\subsubsection{Reverse Holder inequalities (one--step)}

\paragraph{Inverse case.}

\begin{definition}[Inverse-power cutoff, \leanname{superPowerCutoff}]
For a smooth $\eta$, positive $u$, and $p > 0$,
\[
  \mathrm{PCI}_d(\eta,u,p)(x) := \eta(x)\,|u(x)^{-1}|^{p/2}.
\]
\end{definition}

\begin{theorem}[Inverse-power one-step gain, \leanname{supersolution\_preMoser\_inverse}]
\label{thm:preMoser-inv}
Let $d>2$, $A$ as above, $u > 0$ a supersolution, $p > 0$, and
$0 < r < s \leq 1$ with $\int_{\Bz{s}} |u^{-1}|^{p}\,dx < \infty$. Then
$|u^{-1}|^{\chi p}$ is integrable on $\Bz{r}$ and
\begin{multline*}
  \Big(\int_{\Bz{r}} |u^{-1}|^{\chi p}\,dx\Big)^{1/(\chi p)}
  \\ \leq\;
  \Big(\frac{C_{\mathrm{MA}}(d)}{(s-r)^{2}}\,
    \Big(\Lambda\Big(\frac{p}{1+p}\Big)^{2}+1\Big)\Big)^{1/p}
  \Big(\int_{\Bz{s}} |u^{-1}|^{p}\,dx\Big)^{1/p}.
\end{multline*}
\end{theorem}

\begin{proof}
Apply Theorem~\ref{thm:weighted-cacc} with $\Psi(t) := t^{-p}$ on $(0,\infty)$
(suitably truncated), so that $\Psi'(t) = -p\,t^{-p-1} \leq 0$ and
$\Psi(t)^{2}/(-\Psi'(t)) = t^{-p-1}/p$. Choose a smooth cutoff $\eta$ with
$\eta = 1$ on $\Bz{r}$, $\supp \eta \subseteq \Bz{R}$ where $R := (r+s)/2$,
$0 \leq \eta \leq 1$, $\|\nabla \eta\|_{L^\infty} \leq C_\eta := 2\Mst/(s-r)$.
The weighted Caccioppoli inequality applied to $\varphi := \eta$ then yields
\begin{equation*}
  \int_{\Bz{s}} \eta^{2}\,p\,u^{-p-1}\,\langle A\nabla u,\nabla u\rangle\,dx
  \;\leq\;
  \frac{4\Lambda}{p}
  \int_{\Bz{s}} \|\nabla \eta\|^{2}\,u^{-p-1}\,dx.
\end{equation*}
Combined with the chain-rule identity
$\nabla(\eta\,u^{-p/2}) = (\nabla \eta) u^{-p/2} - (p/2) \eta\,u^{-p/2-1}\nabla u$
and the elementary $|a-b|^{2} \leq 2(a^{2}+b^{2})$ this gives, for the witness
$v := \mathrm{PCI}_d(\eta,u,p)$, the energy bound
\[
  \int_{\Bz{s}} \|\nabla v\|^{2}\,dx
  \;\leq\;
  2\,C_\eta^{2}\,\Big(\Lambda\Big(\frac{p}{1+p}\Big)^{2}+1\Big)
  \int_{\Bz{s}} |u^{-1}|^{p}\,dx;
\]
this is the content of \leanname{superPowerCutoff\_energy\_bound}. Now apply
the Sobolev inequality \leanname{sobolev\_prepare\_on\_ball} to $v$ on $\Bz{s}$
and use $\eta = 1$ on $\Bz{r}$, exactly as in the proof of Theorem~\ref{thm:moser-preMoser}.
The constant $8\Mst^{2}\,C_{\mathrm{gns}}^{2} \leq C_{\mathrm{MA}}$ absorbs the
gradient bound, yielding the displayed reverse Holder inequality.
\end{proof}

\paragraph{Forward (low-power) case.}

\begin{definition}[Forward low-power cutoff, \leanname{superPowerCutoffFwd}]
For smooth $\eta$, $u > 0$, and $p \in (0,1)$,
\[
  \mathrm{PCF}_d(\eta,u,p)(x) := \eta(x)\,|u(x)|^{p/2}.
\]
\end{definition}

\begin{theorem}[Forward low-power one-step gain, \leanname{supersolution\_preMoser\_forward}]
\label{thm:preMoser-fwd}
Let $d>2$, $A$ as above, $u>0$ a supersolution, $0 < p < 1$, and
$0 < r < s \leq 1$ with $\int_{\Bz{s}} |u|^{p}\,dx < \infty$. Then
$|u|^{\chi p}$ is integrable on $\Bz{r}$ and
\begin{multline*}
  \Big(\int_{\Bz{r}} |u|^{\chi p}\,dx\Big)^{1/(\chi p)}
  \\ \leq\;
  \Big(\frac{C_{\mathrm{MA}}(d)}{(s-r)^{2}}\,
    \Big(\Lambda\Big(\frac{p}{1-p}\Big)^{2}+1\Big)\Big)^{1/p}
  \Big(\int_{\Bz{s}} |u|^{p}\,dx\Big)^{1/p}.
\end{multline*}
\end{theorem}

\begin{proof}
This time use $\Psi(t) := -t^{p}$ on $(0,\infty)$, so $\Psi'(t) = -p\,t^{p-1}
\leq 0$ and $\Psi(t)^{2}/(-\Psi'(t)) = t^{p+1}/p$. (Note that for $p \in (0,1)$
the function $t \mapsto t^{p}$ is concave with $\Psi'(0)=+\infty$, but
truncation issues are absorbed in \leanname{superPowerCutoffFwd\_energy\_bound}.)
The weighted Caccioppoli inequality with $\varphi := \eta$ supplies the
analogous energy estimate
\[
  \int_{\Bz{s}} \|\nabla v\|^{2}\,dx
  \;\leq\;
  2 C_\eta^{2}\,\Big(\Lambda\Big(\frac{p}{1-p}\Big)^{2}+1\Big)
  \int_{\Bz{s}} |u|^{p}\,dx
\]
for $v := \mathrm{PCF}_d(\eta,u,p)$. Sobolev embedding then concludes as in the
inverse case.
\end{proof}

\subsubsection{Stage one (inverse iteration)}

We collect the inverse-power constants used in the iteration.

\begin{definition}[Inverse iteration quantities, \leanname{superIterIntegralInv}, \leanname{superIterNormInv}, \leanname{superStepConstInv}]
For $u>0$, $p_0 > 0$, $n \in \N$, set
\begin{align*}
  I_n^{\mathrm{inv}}(u) &:= \int_{\Bz{r_n}} |u^{-1}|^{p_n}\,dx, \\
  N_n^{\mathrm{inv}}(u) &:= \big(I_n^{\mathrm{inv}}(u)\big)^{1/p_n}, \\
  \mathrm{S}_n^{\mathrm{inv}}(A,p_0) &:=
   \Big(\frac{C_{\mathrm{MA}}(d)}{(r_n-r_{n+1})^{2}}\,
     \Big(\Lambda\Big(\frac{p_n}{1+p_n}\Big)^{2}+1\Big)\Big)^{1/p_n}.
\end{align*}
The constant
\begin{equation*}
  C_{\mathrm{wH}}^{0}(d) :=
  \begin{cases}
    C_{\mathrm{Moser}}(d)\,(\chi^{2})^{\sum_{n\geq 0} n q^{n}}, & d > 2,\\
    C_{\mathrm{Moser}}(d), & \text{otherwise},
  \end{cases}
\end{equation*}
(\leanname{C\_weakHarnack0}) absorbs the extra geometric factor coming from the
$(\chi^{2})^{i}$ term in the inverse step bound.
\end{definition}

\begin{lemma}[Step bound, \leanname{superStepConstInv\_le}]
\label{lem:stepInv-le}
For $d>2$ and $p_0 > 0$,
\[
  \mathrm{S}_n^{\mathrm{inv}}(A,p_0)
  \;\leq\;
  \Big(16\,C_{\mathrm{MA}}(d)\,(\Lambda p_0^{2}+1)\,(4\chi^{2})^{n}\Big)^{1/p_n}.
\]
\end{lemma}

\begin{proof}
The factor $C_{\mathrm{MA}}/(r_n - r_{n+1})^{2} = 16\,C_{\mathrm{MA}}\,4^{n}$ as
before. Using $p_n/(1+p_n) \leq p_n$ and $p_n^{2} = p_0^{2}\chi^{2n}$,
\[
  \Lambda\Big(\frac{p_n}{1+p_n}\Big)^{2}+1
  \;\leq\; \Lambda p_n^{2}+1
  \;\leq\; (\Lambda p_0^{2}+1)\,(\chi^{2})^{n}.
\]
Multiplying gives the inner expression $16 C_{\mathrm{MA}}(\Lambda p_0^{2}+1)
(4\chi^{2})^{n}$; raising to $1/p_n$ yields the claim.
\end{proof}

\begin{theorem}[Inverse iteration, \leanname{supersolution\_iteration\_inverse}]
\label{thm:iter-inv}
For $d>2$, $u>0$ a supersolution with $\int_{\Bz{1}} |u^{-1}|^{p_0}\,dx < \infty$,
and every $n \in \N$, $|u^{-1}|^{p_n}$ is integrable on $\Bz{r_n}$ and
\[
  N_n^{\mathrm{inv}}(u)
  \;\leq\;
  \prod_{i=0}^{n-1}\mathrm{S}_i^{\mathrm{inv}}(A,p_0)\;\cdot\; N_0^{\mathrm{inv}}(u).
\]
\end{theorem}

\begin{proof}
Induction on $n$. The base case $n = 0$ uses $r_0 = 1$
(\leanname{moserRadius\_zero}) and $p_0 = p_0$
(\leanname{moserExponentSeq\_zero}): the integrability hypothesis for
$|u^{-1}|^{p_0}$ on $\Bz{1}$ is exactly what is assumed, and the empty product
is $1$ so the iteration bound reads $N_0^{\mathrm{inv}}(u) \le N_0^{\mathrm{inv}}(u)$.

For the inductive step from $n$ to $n+1$, set $p_n := \mathrm{moserExponentSeq}\,d\,p_0\,n$
and observe $p_n > 0$ by \leanname{moserExponentSeq\_pos} (using $d > 2$ and
$p_0 > 0$). Apply the inverse-power one-step pre-estimate
Theorem~\ref{thm:preMoser-inv}
(\leanname{supersolution\_preMoser\_inverse}) with parameters $(p, r, s) =
(p_n, r_{n+1}, r_n)$. The hypotheses required are: positivity of $p_n$ (from
the previous remark), positivity of $r_{n+1}$ (\leanname{moserRadius\_pos}),
the strict inequality $r_{n+1} < r_n$ (\leanname{moserRadius\_succ\_lt}), the
upper bound $r_n \le 1$ (\leanname{moserRadius\_le\_one}), the strict
positivity of $u$ on the unit ball, the supersolution property of $u$, and
the integrability of $|u^{-1}|^{p_n}$ on $\Bz{r_n}$ (which is the inductive
hypothesis after rewriting via $\mathrm{moserExponentSeq}\,d\,p_0\,n = p_n$).

The pre-estimate produces:
(i) integrability of $|u^{-1}|^{\chi p_n}$ on $\Bz{r_{n+1}}$, which by
$\chi p_n = p_{n+1}$ (\leanname{moserExponentSeq\_succ}) is exactly the
integrability assertion at level $n+1$; and
(ii) the one-step bound
$N_{n+1}^{\mathrm{inv}}(u) \le \mathrm{S}_n^{\mathrm{inv}}(A, p_0)\,
N_n^{\mathrm{inv}}(u)$.
Multiplying the inductive bound for $N_n^{\mathrm{inv}}$ by
$\mathrm{S}_n^{\mathrm{inv}}$ and absorbing the additional factor into the
range product
$\prod_{i=0}^{n} \mathrm{S}_i^{\mathrm{inv}} = \bigl(\prod_{i=0}^{n-1} \mathrm{S}_i^{\mathrm{inv}}\bigr)
\cdot \mathrm{S}_n^{\mathrm{inv}}$ closes the induction.
\end{proof}

\begin{lemma}[Geometric majorant, \leanname{supersolution\_geometric\_majorant\_inv}]
\label{lem:geom-maj-inv}
For $d>2$ and $p_0>0$,
\[
  \prod_{i=0}^{n-1}\mathrm{S}_i^{\mathrm{inv}}(A,p_0)
  \;\leq\;
  \big(C_{\mathrm{wH}}^{0}(d)\,(\Lambda p_0^{2}+1)^{d/2}\big)^{1/p_0}.
\]
\end{lemma}

\begin{proof}
By Lemma~\ref{lem:stepInv-le}, the product is bounded by
\[
  \prod_{i=0}^{n-1}\big(K\,L\,(4\chi^{2})^{i}\big)^{1/p_i},
\]
where $K := 16\,C_{\mathrm{MA}}(d)$, $L := \Lambda p_0^{2}+1$, $1/p_i = q^{i}/p_0$.
Taking logs and using $\sum_{i<n} q^{i} \leq d/2$ (\leanname{moserDecayRatio\_sum\_le\_half\_dim})
and $\sum_{i<n} i q^{i} \leq \sum_{i\geq 0} i q^{i}$
(\leanname{moserDecayRatio\_weighted\_sum\_le\_tsum}), the product is at most
\[
  \big(K^{d/2}\,L^{d/2}\,(4\chi^{2})^{\sum_{i\geq 0} i q^{i}}\big)^{1/p_0}.
\]
Since $K^{d/2}\,4^{\sum_{i\geq 0} i q^{i}} \leq C_{\mathrm{Moser}}(d)$ by definition
of $C_{\mathrm{Moser}}$, and $(\chi^{2})^{\sum_{i\geq 0} i q^{i}}$ is exactly the
extra factor included in $C_{\mathrm{wH}}^{0}(d)$, the product is bounded by
$\big(C_{\mathrm{wH}}^{0}(d)\,L^{d/2}\big)^{1/p_0}$.
\end{proof}

\begin{theorem}[Inverse a.e.\ closeout, \leanname{supersolution\_ae\_closeout\_inv}]
\label{thm:ae-closeout-inv}
For $d>2$, $p_0 > 0$, $u>0$ a supersolution and $|u^{-1}|^{p_0}$ integrable on
$\Bz{1}$, for almost every $x \in \Bz{1/2}$,
\[
  |u(x)^{-1}|^{p_0}
  \;\leq\;
  C_{\mathrm{wH}}^{0}(d)\,(\Lambda p_0^{2}+1)^{d/2}\,
  \int_{\Bz{1}} |u^{-1}|^{p_0}\,dx.
\]
\end{theorem}

\begin{proof}
Combine Theorem~\ref{thm:iter-inv} with Lemma~\ref{lem:geom-maj-inv} to get,
for every $n$,
\[
  N_n^{\mathrm{inv}}(u)
  \;\leq\;
  \big(C_{\mathrm{wH}}^{0}(d)\,(\Lambda p_0^{2}+1)^{d/2}\,I_0^{\mathrm{inv}}(u)\big)^{1/p_0}.
\]
The Markov-type closeout argument of Lemma~\ref{lem:moser-closeout}, applied
verbatim with $u_+$ replaced by $u^{-1}$ (and $1<p_0$ replaced by $p_0>0$, which
only affects the choice of test exponent $q_n = 2\chi^{n}$ but not the
convergence) yields the desired a.e.\ inequality.
\end{proof}

\begin{theorem}[Stage one weak Harnack (inverse), \leanname{weak\_harnack\_stage\_one\_inverse}]
\label{thm:wH-stage1-inv}
Let $d>2$, $A$ on $\Bz{1}$, $u>0$ on $\Bz{1}$ a supersolution, and $p_0 > 0$. Then
\begin{multline*}
  (\Lambda p_0^{2}+1)^{-d/2}\,
  \Big(\int_{\Bz{1}} |u^{-1}|^{p_0}\,dx\Big)^{-1}
  \\
  \;\leq\; C_{\mathrm{wH}}^{0}(d)\;(\essinf_{\Bz{1/2}} u)^{p_0}.
\end{multline*}
\end{theorem}

\begin{proof}
If $|u^{-1}|^{p_0}$ is not integrable on $\Bz{1}$, the integral is taken to be
$+\infty$ and the LHS is $0 \leq \mathrm{RHS}$.

Otherwise, set $I := \int_{\Bz{1}} |u^{-1}|^{p_0}\,dx$ and
$c_0 := C_{\mathrm{wH}}^{0}(d)\,(\Lambda p_0^{2}+1)^{d/2}\,I$. By
Theorem~\ref{thm:ae-closeout-inv} we have $|u^{-1}|^{p_0} \leq c_0$ a.e.\ on
$\Bz{1/2}$, equivalently $u^{p_0} \geq c_0^{-1}$ a.e., i.e.\
$u \geq c_0^{-1/p_0}$ a.e.\ on $\Bz{1/2}$. Therefore
$\essinf_{\Bz{1/2}} u \geq c_0^{-1/p_0}$, so
$(\essinf_{\Bz{1/2}} u)^{p_0} \geq c_0^{-1}$. Multiplying by
$C_{\mathrm{wH}}^{0}(d) \geq 0$ and using
$L^{-d/2} I^{-1} = C_{\mathrm{wH}}^{0}(d)\,c_0^{-1}$ (where $L = \Lambda p_0^{2}+1$)
gives the displayed inequality.
\end{proof}

\subsubsection{Stage two (forward low-power iteration)}

The forward iteration is run only finitely many steps: starting from
$p_0 := q\,\chi^{-m}$ for a target exponent $q \in (0,1)$ and a depth $m$
chosen so that $p_0 \leq p$, we iterate $m+1$ times to reach the level $p_m =
q$. The forward step constant
\begin{equation*}
  \mathrm{S}_n^{\mathrm{fwd}}(A,p_0) :=
  \Big(\frac{C_{\mathrm{MA}}(d)}{(r_n-r_{n+1})^{2}}\,
    \Big(\Lambda\Big(\frac{p_n}{1-p_n}\Big)^{2}+1\Big)\Big)^{1/p_n}
\end{equation*}
\leanname{superStepConstFwd} is finite as long as $p_n < 1$, which is
guaranteed by $p_n \leq q < 1$ for $n \leq m$.

\begin{lemma}[Forward depth, \leanname{exists\_forward\_iteration\_depth}]
\label{lem:fwd-depth}
For $0 < p < q < 1$ there exists $m \in \N$ with
\[
  q\,\chi^{-m} < p \quad \text{and} \quad p \leq \chi\,\big(q\,\chi^{-m}\big).
\]
\end{lemma}

\begin{proof}
Since $\chi > 1$, $q\,\chi^{-m} \to 0$ as $m \to \infty$, so any sufficiently
large $m$ satisfies the first inequality; choose the smallest such $m$, then
$q\,\chi^{-(m-1)} \geq p$, equivalently $p \leq \chi\,(q\,\chi^{-m})$.
\end{proof}

\begin{lemma}[Forward step bound, \leanname{superStepConstFwd\_le}]
For $d>2$, $0 < q < 1$, $m \in \N$, $i \leq m$, and $p_0 := q\,\chi^{-m}$,
\[
  \mathrm{S}_i^{\mathrm{fwd}}(A,p_0)
  \;\leq\;
  \Big(16\,C_{\mathrm{MA}}(d)\,\Big(\frac{\Lambda p_0^{2}}{(1-q)^{2}}+1\Big)\,
  (4\chi^{2})^{i}\Big)^{1/p_i}.
\]
\end{lemma}

\begin{proof}
For $i \leq m$, $p_i = p_0\,\chi^{i} = q\,\chi^{i-m} \leq q$, hence
$1 - p_i \geq 1 - q > 0$ and
\[
  \frac{p_i}{1-p_i} \leq \frac{p_i}{1-q}, \qquad
  \Big(\frac{p_i}{1-p_i}\Big)^{2}\leq \frac{p_i^{2}}{(1-q)^{2}}
  \leq \frac{p_0^{2}\chi^{2i}}{(1-q)^{2}}.
\]
Inserting into $\mathrm{S}_i^{\mathrm{fwd}}$ and using
$C_{\mathrm{MA}}/(r_i - r_{i+1})^{2} = 16\,C_{\mathrm{MA}}\,4^{i}$ gives the bound.
\end{proof}

\begin{theorem}[Forward iteration, \leanname{supersolution\_iteration\_forward}]
\label{thm:iter-fwd}
For $d>2$, $u>0$ a supersolution, $0<q<1$, $m \in \N$, set $p_0 := q\,\chi^{-m}$.
If $|u|^{p_0}$ is integrable on $\Bz{1}$, then for every $n \leq m+1$
$|u|^{p_n}$ is integrable on $\Bz{r_n}$ and
\[
  N_n^{\mathrm{fwd}}(u)
  \;\leq\;
  \prod_{i=0}^{n-1} \mathrm{S}_i^{\mathrm{fwd}}(A,p_0)\;\cdot\;
  N_0^{\mathrm{fwd}}(u).
\]
\end{theorem}

\begin{proof}
Induction on $n \leq m+1$. The base case is trivial. For the step, note that
$p_n = p_0\,\chi^{n} = q\,\chi^{n-m}$, so for $n \leq m$ we have $p_n \leq q < 1$
(this is what allows the application of Theorem~\ref{thm:preMoser-fwd}). The
one-step bound gives $N_{n+1}^{\mathrm{fwd}}(u) \leq
\mathrm{S}_n^{\mathrm{fwd}}(A,p_0)\,N_n^{\mathrm{fwd}}(u)$, and combining with
the inductive hypothesis yields the displayed product bound.
\end{proof}

\begin{lemma}[Forward geometric majorant, \leanname{supersolution\_geometric\_majorant\_fwd}]
\label{lem:geom-maj-fwd}
With the notation above, for any $n \leq m+1$,
\[
  \prod_{i=0}^{n-1}\mathrm{S}_i^{\mathrm{fwd}}(A,p_0)
  \;\leq\;
  \Big(C_{\mathrm{wH}}^{0}(d)\,
    \Big(\frac{\Lambda p_0^{2}}{(1-q)^{2}}+1\Big)^{d/2}\Big)^{1/p_0}.
\]
\end{lemma}

\begin{proof}
Identical to the proof of Lemma~\ref{lem:geom-maj-inv}, with $\Lambda p_0^{2}+1$
replaced by $\Lambda p_0^{2}/(1-q)^{2}+1$ throughout.
\end{proof}

\begin{definition}[Forward closeout constant, \leanname{C\_weakHarnack0Forward}]
For $d>2$,
\[
  C_{\mathrm{wH,fwd}}^{0}(d) :=
  \big(C_{\mathrm{wH}}^{0}(d)\,(|\Bz{1}|+1)\big)^{\chi}.
\]
\end{definition}

\begin{theorem}[Stage one weak Harnack (forward), \leanname{weak\_harnack\_stage\_one\_forward}]
\label{thm:wH-stage1-fwd}
Let $d>2$, $A$ on $\Bz{1}$, $u>0$ a supersolution on $\Bz{1}$, and
$0 < p < q < 1$. Then
\begin{multline*}
  \Big(\int_{\Bz{1/2}} |u|^{q d/(d-2)}\,dx\Big)^{p(d-2)/(qd)}
  \\
  \;\leq\;
  C_{\mathrm{wH,fwd}}^{0}(d)\,
  \Big(\frac{\Lambda p^{2}}{(1-q)^{2}}+1\Big)^{d\chi/2}
  \int_{\Bz{1}} |u|^{p}\,dx.
\end{multline*}
\end{theorem}

\begin{proof}
Set $\chi := d/(d-2)$, $q\chi := q d/(d-2)$. By Lemma~\ref{lem:fwd-depth} pick
$m$ with $p_0 := q\chi^{-m} < p$ and $p \leq \chi p_0$, and note $p_0 < p < 1$.
Theorem~\ref{thm:iter-fwd} with this $p_0$ at depth $n = m+1$ gives the iterated
$L^{p_{m+1}} = L^{q\chi}$ bound on $\Bz{r_{m+1}}$, namely
\begin{multline*}
  \Big(\int_{\Bz{r_{m+1}}} |u|^{q\chi}\,dx\Big)^{1/(q\chi)}
  \\ \leq\;
  \prod_{i=0}^{m}\mathrm{S}_i^{\mathrm{fwd}}(A,p_0)\;\cdot\;
  \Big(\int_{\Bz{1}} |u|^{p_0}\,dx\Big)^{1/p_0}.
\end{multline*}
By Lemma~\ref{lem:geom-maj-fwd}, the product is at most
\[
  \Big(C_{\mathrm{wH}}^{0}(d)\,
    \big(\Lambda p_0^{2}/(1-q)^{2}+1\big)^{d/2}\Big)^{1/p_0}.
\]
Use $\Bz{1/2} \subseteq \Bz{r_{m+1}}$ to restrict to $\Bz{1/2}$, raise the
inequality to the $p$-th power, and absorb the change of base from $p_0$ to $p$
via the volume comparison
\leanname{weak\_harnack\_stage\_one\_forward\_power\_upgrade}: by Holder with
exponents $p_0 < p \leq \chi p_0$ on the unit ball,
\[
  \Big(\int_{\Bz{1}} |u|^{p_0}\,dx\Big)^{1/p_0}
  \leq
  \Big(\int_{\Bz{1}} |u|^{p}\,dx\Big)^{1/p}\,
  |\Bz{1}|^{1/p_0 - 1/p},
\]
and consequently
\[
  \Big(\int_{\Bz{1}} |u|^{p_0}\,dx\Big)^{p/p_0}
  \leq (|\Bz{1}|+1)^{\chi}\,\int_{\Bz{1}} |u|^{p}\,dx,
\]
which absorbs the volume into $C_{\mathrm{wH,fwd}}^{0}(d)$. Together with
$(\Lambda p_0^{2}/(1-q)^{2}+1)^{d/(2)\cdot p/p_0} \leq
(\Lambda p^{2}/(1-q)^{2}+1)^{d\chi/2}$
(monotonicity in both base and exponent, since $p/p_0 \leq \chi$ and
$p_0^{2} \leq p^{2}$), the displayed inequality follows. The exponent
$p(d-2)/(qd) = p/(q\chi)$ comes from raising the $L^{q\chi}$-norm to the
$p$-th power.
\end{proof}

\begin{remark}[Reverse-Holder reformulation]
Theorems~\ref{thm:wH-stage1-inv} and~\ref{thm:wH-stage1-fwd} are precisely the
two halves of the weak Harnack inequality at this stage of the development:
the first asserts a lower bound on $u$ in $\Bz{1/2}$ in terms of an inverse
integral on $\Bz{1}$, while the second is a forward $L^{q\chi}$--$L^{p}$
reverse Holder inequality on the same pair of balls. The crossover between
the two stages, completing the weak Harnack inequality, is performed in the
file \texttt{WeakHarnack.lean} using the constants $C_{\mathrm{wH}}^{0}(d)$
and $C_{\mathrm{wH,fwd}}^{0}(d)$ assembled here.
\end{remark}


%% ========================================================================
\section{Crossover, Weak Harnack, Harnack, and H\"older Regularity}\label{sec:crossover}
%% ========================================================================

%% --- Section 9 ---
\subsection{Crossover Estimate}

Let $d \in \N$, $d \geq 3$, and $E = \R^d$ throughout. Write $\Bz r$ for the
open Euclidean ball of radius $r$ centred at the origin. Given a normalized
elliptic coefficient $A$ on $\Bz 1$ with ellipticity constants $\lambda_A,
\Lambda_A$ (with $\lambda_A = 1$ after normalization, so $\Lambda_A \geq 1$),
write $\IsSupersolution{A}{u}$ for the assertion that $u \in W^{1,2}_{\loc}$
satisfies $-\dv(A\nabla u) \geq 0$ in the weak sense, and $\IsSolution{A}{u}$
for the corresponding equality.

The crossover estimate, due in this form to Moser, bridges positive and
negative powers of a positive supersolution. Throughout this section we work
with a fixed positive constant
\[
  \Cpoinc = \Cpoinc(d) > 0
\]
denoting the dimensional Poincar\'e constant on the unit ball, a fixed
constant $\Mst > 1$ used in the smooth cutoff construction, and the
dimensional John--Nirenberg constant $\CJN(d) > 0$.

\begin{definition}[\leanname{c\_crossover\_bmo\_scale}, \leanname{c\_crossover'}]
We define the BMO scale and crossover exponent
\begin{align*}
  c_{\BMO}(d) &:= \bigl(\vol(\Bz 1)\bigr)^{-1/2}\,\Cpoinc(d) \cdot
    \tfrac12^{1 - d/2} \cdot 8\Mst \,\bigl(\vol(\Bz 1)\bigr)^{1/2}, \\
  c'(d)       &:= \frac{1}{2\,\CJN(d)\,c_{\BMO}(d) + 1}.
\end{align*}
\end{definition}

\begin{lemma}[\leanname{crossoverC'\_pos}, \leanname{c\_crossover'\_lt\_one}]
\label{lem:c-crossover-bounds}
$0 < c'(d) < 1$ and $c'(d) \leq 1$.
\end{lemma}

\begin{proof}
The constant $c_{\BMO}(d)$ is a product of strictly positive terms, so
$c_{\BMO}(d) > 0$. Since $\CJN(d) > 0$, the denominator
$2\,\CJN(d)\,c_{\BMO}(d) + 1$ is strictly greater than $1$, hence its
reciprocal lies in $(0,1)$.
\end{proof}

\subsubsection{Logarithmic gradient bound}

The first ingredient is the Caccioppoli-type inequality for $\log u$, obtained
by testing the supersolution against $\varphi^2/(u + \varepsilon)$.

\begin{theorem}[\leanname{log\_gradient\_bound\_eps}]
\label{thm:log-grad-eps}
Let $\Omega \subseteq E$ be open, $A$ an elliptic coefficient on $\Omega$
with ellipticity constants $\lambda, \Lambda$, $u \in W^{1,2}_{\loc}(\Omega)$
positive on $\Omega$ with $\IsSupersolution A u$, and let
$\varphi \in C_c^\infty(\Omega)$. Then for every $\varepsilon > 0$,
\begin{equation*}
  \int_\Omega \frac{\varphi^2 \,|\nabla u|^2}{(u+\varepsilon)^2}\,dx
    \;\leq\; \frac{4\Lambda}{\lambda}\int_\Omega |\nabla\varphi|^2\,dx.
\end{equation*}
\end{theorem}

\begin{proof}
Set $\psi := \varphi^2/(u+\varepsilon)$. By the smooth chain-rule witness
construction \leanname{build\_test\_function\_eps}, $\psi$ lies in
$H^1_0(\Omega)$, $\psi \geq 0$, and its weak gradient satisfies almost
everywhere
\begin{equation*}
  \nabla\psi = \frac{2\varphi\,\nabla\varphi}{u+\varepsilon}
    - \frac{\varphi^2\,\nabla u}{(u+\varepsilon)^2}.
\end{equation*}
Since $u$ is a supersolution and $\psi \geq 0$,
\begin{equation*}
  0 \;\leq\; \int_\Omega \langle A\nabla u, \nabla\psi\rangle\,dx
   \;=\; \int_\Omega \frac{2\varphi}{u+\varepsilon}\langle A\nabla u,
       \nabla\varphi\rangle\,dx
   \;-\; \int_\Omega \frac{\varphi^2}{(u+\varepsilon)^2}
       \langle A\nabla u, \nabla u\rangle\,dx.
\end{equation*}
Set
\begin{align*}
  Q &:= \int_\Omega \frac{\varphi^2}{(u+\varepsilon)^2}\langle A\nabla u,\nabla u\rangle\,dx, \\
  X &:= \int_\Omega \frac{\varphi}{u+\varepsilon}\langle A\nabla u,\nabla\varphi\rangle\,dx.
\end{align*}
The supersolution inequality reads $Q \leq 2X$. By Cauchy--Schwarz applied to
the bilinear form $\langle A\,\cdot,\cdot\rangle$,
\begin{equation*}
  X^2 \;\leq\; Q \cdot \int_\Omega \langle A\nabla\varphi,\nabla\varphi\rangle\,dx
    \;\leq\; Q \cdot \Lambda \int_\Omega |\nabla\varphi|^2\,dx.
\end{equation*}
Combining $Q \leq 2X$ with $X^2 \leq Q\,\Lambda\int|\nabla\varphi|^2$, the
elementary inequality $Q^2 \leq 4QP$ implies $Q \leq 4P$ where
$P := \Lambda \int_\Omega |\nabla\varphi|^2$ (see
\leanname{log\_gradient\_algebraic\_conclusion}). Hence
\begin{equation*}
  \lambda \int_\Omega \frac{\varphi^2|\nabla u|^2}{(u+\varepsilon)^2}\,dx
    \;\leq\; Q \;\leq\; 4\Lambda \int_\Omega |\nabla\varphi|^2\,dx,
\end{equation*}
and dividing by $\lambda > 0$ gives the claim.
\end{proof}

\begin{theorem}[\leanname{log\_gradient\_bound\_of\_supersolution}]
\label{thm:log-grad}
Under the same hypotheses, with $u > 0$ on $\Omega$,
\begin{equation*}
  \int_\Omega \frac{\varphi^2\,|\nabla u|^2}{u^2}\,dx
    \;\leq\; \frac{4\Lambda}{\lambda}\int_\Omega |\nabla\varphi|^2\,dx.
\end{equation*}
\end{theorem}

\begin{proof}
Apply Theorem~\ref{thm:log-grad-eps} with $\varepsilon_n = 1/(n+1)$. The
integrand
\[
  g_n(x) := \varphi(x)^2|\nabla u(x)|^2/(u(x)+\varepsilon_n)^2
\]
increases monotonically as $\varepsilon_n \downarrow 0$ to
$g_0(x) := \varphi(x)^2|\nabla u(x)|^2/u(x)^2$ at every point of $\Omega$
since $u > 0$ there. By Fatou's lemma,
\begin{equation*}
  \int_\Omega g_0\,dx \;\leq\; \liminf_{n\to\infty}\int_\Omega g_n\,dx
    \;\leq\; \frac{4\Lambda}{\lambda}\int_\Omega |\nabla\varphi|^2\,dx.
\end{equation*}
\end{proof}

\subsubsection{BMO control of \texorpdfstring{$\log u$}{log u}}

For $\varepsilon > 0$ define the regularized log
$v_\varepsilon(x) := -\log(u(x)+\varepsilon) + \log\varepsilon$ and its smooth
chain-rule witness $\leanname{regularizedLogFun}$. Combining the log-gradient
bound with the Poincar\'e inequality on $\Bz{1/2}$ gives the BMO estimate.

\begin{theorem}[\leanname{regularizedLog\_unit\_halfBall\_meanOscillation}]
\label{thm:bmo-half}
Let $A$ be normalized elliptic on $\Bz 1$, $u > 0$ a supersolution. Then
for $v := v_\varepsilon$,
\begin{equation*}
  \fint_{\Bz{1/2}} \Bigl|v(x) - \fint_{\Bz{1/2}} v\Bigr|\,dx
    \;\leq\; c_{\BMO}(d)\,\Lambda_A^{1/2}.
\end{equation*}
\end{theorem}

\begin{proof}
Pick a cutoff $\eta \in C_c^\infty(\Bz{3/4})$ with $\eta = 1$ on $\Bz{1/2}$,
$0 \leq \eta \leq 1$, and $|\nabla \eta| \leq 8\Mst$ (existence given by
\leanname{Cutoff.exists}). Theorem~\ref{thm:log-grad-eps} (with
$\varphi = \eta$, $\lambda_A = 1$) yields
\begin{equation*}
  \int_{\Bz{1/2}} |\nabla v_\varepsilon|^2\,dx
    \;\leq\; \int_{\Bz 1} \eta^2 \frac{|\nabla u|^2}{(u+\varepsilon)^2}\,dx
    \;\leq\; 4\Lambda_A \int_{\Bz 1}|\nabla \eta|^2\,dx
    \;\leq\; 4\Lambda_A \cdot (8\Mst)^2 \vol(\Bz 1).
\end{equation*}
Hence $\|\nabla v_\varepsilon\|_{L^2(\Bz{1/2})} \leq 16\Mst\,\Lambda_A^{1/2}\,\vol(\Bz 1)^{1/2}$.
By the Poincar\'e--Wirtinger inequality on $\Bz{1/2}$ with constant
$\Cpoinc(d) \cdot (1/2)^{1-d/2}$ (the dimensional rescaling factor),
\begin{multline*}
  \fint_{\Bz{1/2}}\Bigl|v_\varepsilon - \fint_{\Bz{1/2}} v_\varepsilon\Bigr|\,dx
    \;\leq\; \vol(\Bz{1/2})^{-1/2} \Cpoinc(d) (1/2)^{1-d/2}
        \|\nabla v_\varepsilon\|_{L^2(\Bz{1/2})} \\
    \;\leq\; \vol(\Bz 1)^{-1/2}\Cpoinc(d) (1/2)^{1-d/2}\cdot
        8\Mst\,\Lambda_A^{1/2}\,\vol(\Bz 1)^{1/2}
    \;=\; c_{\BMO}(d)\,\Lambda_A^{1/2}.
\end{multline*}
\end{proof}

The same statement holds on every ball $B(x_0, R/2)$ with $\overline{B(x_0,R)}
\subset \Bz 1$, by affine rescaling $\leanname{rescaleToUnitBall}$
(see \leanname{regularizedLog\_ball\_meanOscillation}).

\subsubsection{John--Nirenberg exponential integrability}

\begin{definition}[\leanname{crossoverJNTheta}, \leanname{crossoverJNKloc}, \leanname{crossoverJNKhalf}]
\label{def:JN-constants}
Set
\begin{align*}
  \theta(d)  &:= 24 \cdot 2^{d+1} \cdot \bigl(c'(d)\,c_{\BMO}(d)\bigr), \\
  K_{\loc}(d)&:= \exp(\theta(d)) \cdot (1 + \CJN(d)),\\
  K_{1/2}(d) &:= \lceil 97^d \rceil \cdot K_{\loc}(d).
\end{align*}
\end{definition}

\begin{theorem}[\leanname{regularizedLog\_halfBall\_exp\_average\_to\_origin\_le}]
\label{thm:exp-half}
With $p := c'(d)/\Lambda_A^{1/2}$, $v := v_\varepsilon$, and
$a := \fint_{\Bz{1/48}} v\,dz$,
\begin{equation*}
  \fint_{\Bz{1/2}} \exp(p\,|v(z) - a|)\,dz \;\leq\; K_{1/2}(d).
\end{equation*}
\end{theorem}

\begin{proof}
Apply the local John--Nirenberg lemma at scale $\rho = 1/48$ on a finite cover
of $\Bz{1/2}$ by balls $B(x_i, 1/48)$ with $\#\{x_i\} \leq \lceil 97^d\rceil$.
On each cover ball, Theorem~\ref{thm:bmo-half} (rescaled to $B(x_i, R)$ via
\leanname{regularizedLog\_ball\_meanOscillation}) gives
$\|v - \langle v\rangle\|_{\BMO} \leq c_{\BMO}(d)\,\Lambda_A^{1/2}$. The
classical John--Nirenberg theorem $\leanname{C\_JN}$ then yields
\[
  \fint_{B(x_i,1/48)} \exp\bigl(p\,|v - \langle v\rangle_{B(x_i,1/48)}|\bigr)\,dz
    \;\leq\; 1 + \CJN(d).
\]
The shift between local and global averages is controlled by $\theta(d) =
24 \cdot 2^{d+1}\,c'(d)\,c_{\BMO}(d)$ via the standard telescoping
$\leanname{regularizedLog\_smallBallAverage\_to\_origin\_le}$. Multiplying
each local exponential bound by $\exp(p\,\theta(d)) = \exp(\theta(d))$ (since
$p\cdot \delta = \theta(d)$ by direct computation, using
$\delta := 24\cdot 2^{d+1}\,(c_{\BMO}(d)\Lambda_A^{1/2})$ and
$p \cdot \Lambda_A^{1/2} = c'(d)$), and summing over the $\lceil 97^d\rceil$
cover balls gives the asserted bound by $K_{1/2}(d)$.
\end{proof}

\subsubsection{Crossover product bound}

\begin{theorem}[\leanname{crossover\_product\_bound\_exists}]
\label{thm:crossover-exists}
There exists a constant $C(d) \geq 1$ such that for every normalized elliptic
$A$ on $\Bz 1$, every positive supersolution $u$, and $p_0 := c'(d)/\Lambda_A^{1/2}$,
\[
  \Bigl(\fint_{\Bz{1/2}} u^{p_0}\,dx\Bigr)\Bigl(\fint_{\Bz{1/2}} u^{-p_0}\,dx\Bigr)
    \;\leq\; C(d).
\]
\end{theorem}

\begin{proof}
By a regularization step $u \mapsto u + \varepsilon$, it suffices to bound
\[
  J_\varepsilon := \Bigl(\fint_{\Bz{1/2}}(u+\varepsilon)^{p_0}\Bigr)
        \Bigl(\fint_{\Bz{1/2}}(u+\varepsilon)^{-p_0}\Bigr)
\]
uniformly in $\varepsilon > 0$ by $K_{1/2}(d)^2$. With $v_\varepsilon =
-\log(u+\varepsilon) + \log\varepsilon$, $a := \fint_{\Bz{1/48}} v_\varepsilon$,
$F_\pm(x) := \exp(\pm p_0(v_\varepsilon - a))$ and
$F_{\mathrm{abs}} := \exp(p_0|v_\varepsilon - a|)$, we have $F_\pm \leq
F_{\mathrm{abs}}$ pointwise. Theorem~\ref{thm:exp-half} together with
$\leanname{setAverage\_mono\_of\_nonneg\_local}$ yields
\[
  \fint_{\Bz{1/2}} F_+\,dx \leq K_{1/2}(d), \qquad
  \fint_{\Bz{1/2}} F_-\,dx \leq K_{1/2}(d).
\]
Now $(u+\varepsilon)^{p_0} = \exp(p_0(\log\varepsilon - a))\,F_+$ and
$(u+\varepsilon)^{-p_0} = \exp(p_0(a - \log\varepsilon))\,F_-$ pointwise on
$\Bz{1/2}$. The two scalar prefactors multiply to $1$ (cancellation
$\leanname{hcancel}$), so
\[
  J_\varepsilon \;=\; \Bigl(\fint F_+\Bigr)\Bigl(\fint F_-\Bigr)
    \;\leq\; K_{1/2}(d)^2.
\]
Letting $\varepsilon \downarrow 0$ along $\varepsilon_n = 1/(n+1)$, the
shifted integrals converge to the original ones by dominated/monotone
convergence (the integrability of $|u|^{p_0}$ is supplied by the
$W^{1,2}$ estimate using $p_0 < 1 \leq 2$, see
\leanname{hpow\_int}). Choose $C(d) := \max\{1, K_{1/2}(d)^2\}$.
\end{proof}

\begin{definition}[\leanname{C\_crossover'}]
\label{def:Ccrossover}
$C_{\mathrm{cross}}(d) := C(d)$ from Theorem~\ref{thm:crossover-exists}, so
$C_{\mathrm{cross}}(d) \geq 1$.
\end{definition}

\begin{theorem}[\leanname{crossover\_estimate}]
\label{thm:crossover-estimate}
Let $d \geq 3$, $A$ normalized elliptic on $\Bz 1$, and $u > 0$ a
supersolution. Then with $c := c'(d)/\Lambda_A^{1/2}$,
\[
  \Bigl(\fint_{\Bz{1/2}} u^c\,dx\Bigr)\Bigl(\fint_{\Bz{1/2}} u^{-c}\,dx\Bigr)
    \;\leq\; C_{\mathrm{cross}}(d).
\]
\end{theorem}

\begin{proof}
Set $v(x) := -\log u(x)$, $v_{\mathrm{avg}} := \fint_{\Bz{1/2}} v\,dx$,
$w(x) := \exp(c(v_{\mathrm{avg}} - v(x)))$, and
$w^{-1}(x) := \exp(c(v(x) - v_{\mathrm{avg}}))$, so $w \cdot w^{-1} \equiv 1$.
On $\Bz{1/2}$ where $u > 0$, $u^c = e^{-c v_{\mathrm{avg}}}\,w$ and
$u^{-c} = e^{c v_{\mathrm{avg}}}\,w^{-1}$, hence
\[
  \Bigl(\fint u^c\Bigr)\Bigl(\fint u^{-c}\Bigr)
    \;=\; e^{-cv_{\mathrm{avg}}}e^{cv_{\mathrm{avg}}}
        \Bigl(\fint w\Bigr)\Bigl(\fint w^{-1}\Bigr)
    \;=\; \Bigl(\fint w\Bigr)\Bigl(\fint w^{-1}\Bigr).
\]
Apply Theorem~\ref{thm:crossover-exists}.
\end{proof}

\begin{corollary}[\leanname{crossover\_estimate\_unaveraged}]
\label{cor:crossover-unaveraged}
Multiplying through by $\vol(\Bz{1/2})^2 > 0$,
\[
  \Bigl(\int_{\Bz{1/2}} u^c\Bigr)\Bigl(\int_{\Bz{1/2}} u^{-c}\Bigr)
    \;\leq\; C_{\mathrm{cross}}(d)\,\vol(\Bz{1/2})^2.
\]
\end{corollary}

%% --- Section 10 ---
\subsection{Weak Harnack Inequality}

\begin{definition}[\leanname{weakHarnackP0}, \leanname{weak\_harnack\_decay\_exp}]
\label{def:p0}
For a normalized elliptic $A$ on $\Bz 1$, set
\[
  p_0 := p_0(A) := \frac{c'(d)}{\Lambda_A^{1/2}}, \qquad
  \beta(d) := \frac{d\,\chi(d)}{c'(d)},
\]
where $\chi(d) := d/(d-2)$ is the Sobolev exponent (\leanname{moserChi}).
\end{definition}

\begin{lemma}[\leanname{p\_0\_pos}, \leanname{p\_0\_lt\_one}, \leanname{Lambda\_mul\_p\_0\_sq}]
$0 < p_0(A) < 1$, $p_0(A) \leq c'(d)$, and $\Lambda_A\,p_0(A)^2 = c'(d)^2$.
\end{lemma}

\begin{proof}
Since $\Lambda_A \geq 1$, $\Lambda_A^{1/2} \geq 1$, so
$0 < p_0 = c'(d)/\Lambda_A^{1/2} \leq c'(d) < 1$. The identity follows from
$(c'(d)/\Lambda_A^{1/2})^2 \cdot \Lambda_A = c'(d)^2$.
\end{proof}

The proof of the weak Harnack inequality chains three pieces at half scale.
We rescale $v(z) := u(z/2)$ from $\Bz{1/2}$ to $\Bz 1$ via the rescaled
coefficient $A' := \leanname{rescaledCoeff}\,A$ which preserves $\Lambda$.
The three steps are:

\begin{itemize}
\item \emph{Forward Moser at level $p_0$}
  (\leanname{weak\_harnack\_stage\_one\_forward\_ball}):
\begin{equation*}
\Bigl(\int_{\Bz{1/4}} u^{q\,d/(d-2)}\Bigr)^{p_0(d-2)/(qd)}
\leq C_{\mathrm{wH},0\to}(d)\,(\Lambda_A p_0^2/(1-q)^2 + 1)^{d\chi/2}
\int_{\Bz{1/2}} u^{p_0}.
\end{equation*}

\item \emph{Crossover} (Corollary~\ref{cor:crossover-unaveraged}) at the level
$p_0$:
\begin{equation*}
\Bigl(\int_{\Bz{1/2}} u^{p_0}\Bigr)\Bigl(\int_{\Bz{1/2}} u^{-p_0}\Bigr)
\leq C_{\mathrm{cross}}(d)\,\vol(\Bz{1/2})^2.
\end{equation*}

\item \emph{Inverse Moser}
  (\leanname{weak\_harnack\_stage\_one\_inverse\_ball}):
\begin{equation*}
(\Lambda_A p_0^2 + 1)^{-d/2}\Bigl(\int_{\Bz{1/2}} u^{-p_0}\Bigr)^{-1}
\leq C_{\mathrm{wH},0}(d)\,\bigl(\essinf_{\Bz{1/4}} u\bigr)^{p_0}.
\end{equation*}
\end{itemize}

Multiplying these three estimates yields the chain.

\begin{theorem}[\leanname{weak\_harnack\_chain}]
\label{thm:wH-chain}
For $d \geq 3$, $A$ normalized on $\Bz 1$, $u > 0$ a supersolution, and
$0 < q < 1$:
\begin{equation*}
\Bigl(\int_{\Bz{1/4}} u^{q d/(d-2)}\,dx\Bigr)^{p_0(d-2)/(qd)}
\leq \frac{C_{\mathrm{chain}}(d)}{(1-q)^{d\chi(d)}}
\Bigl(\essinf_{\Bz{1/4}} u\Bigr)^{p_0}.
\end{equation*}
\end{theorem}

\begin{proof}
We split into the case $p_0 < q$ (the natural case) and the case $q \leq p_0$
(handled by interpolating with $q' := (p_0+1)/2$). Assume $p_0 < q$ first.

By regularization $u \mapsto u + \delta$ (which preserves the supersolution
property by \leanname{isSupersolution\_add\_const\_unitBall}), and using the
identity $\Lambda_A p_0^2 = c'(d)^2$, the product of the forward Moser
estimate, the crossover bound and the inverse Moser estimate gives
\begin{multline*}
  \Bigl(\int_{\Bz{1/4}}(u+\delta)^{q d/(d-2)}\Bigr)^{p_0(d-2)/(q d)}
  \times \bigl(c'(d)^2 + 1\bigr)^{d/2}
  \\[-2pt]
  \leq C_{\mathrm{wH},0\to}\,C_{\mathrm{wH},0}\,C_{\mathrm{cross}}\,
       \vol(\Bz{1/2})^2\,(c'(d)^2/(1-q)^2 + 1)^{d\chi/2}
       \bigl(\essinf_{\Bz{1/4}}(u+\delta)\bigr)^{p_0}.
\end{multline*}
By \leanname{weak\_harnack\_forward\_factor\_bound} the right factor is
absorbed into $(1-q)^{-d\chi}$ times a dimensional constant, so collecting
terms produces a constant $C_{\mathrm{chainMain}}(d)$ with
\[
  \Bigl(\int_{\Bz{1/4}} (u+\delta)^{qd/(d-2)}\Bigr)^{p_0(d-2)/(qd)}
  \;\leq\; \frac{C_{\mathrm{chainMain}}(d)}{(1-q)^{d\chi(d)}}
        \bigl(\essinf_{\Bz{1/4}}(u+\delta)\bigr)^{p_0}.
\]
Letting $\delta = 1/(n+1) \downarrow 0$ and using the dominated/monotone
convergence theorems on each side concludes the case $p_0 < q$.

For $q \leq p_0$ pick $q' := (p_0+1)/2 \in (p_0,1)$ and apply the previous
case with $q'$. Since $q d/(d-2) \leq q' d/(d-2)$, monotonicity of $L^p$ norms
on $\Bz{1/4}$ (\leanname{quarterBall\_lpNorm\_mono}) raises the resulting
$L^{q' d/(d-2)}$ norm to dominate the $L^{q d/(d-2)}$ norm, and a uniform
bound $C_{\mathrm{chainMain}}/(1-q')^{d\chi} \leq C_{\mathrm{chain}}/(1-q)^{d\chi}$
\leanname{(hmain\_to\_public)} then yields the claim.
\end{proof}

\begin{definition}[\leanname{C\_weakHarnack}]
\label{def:CwH}
$
  C_{wH}(d) \;:=\; C_{\mathrm{chain}}(d)^{1/c'(d)}.
$
By Lemma~\ref{lem:c-crossover-bounds} and $C_{\mathrm{chain}} \geq 1$, we have
$C_{wH}(d) \geq 1$ (\leanname{one\_le\_C\_weakHarnack}).
\end{definition}

\begin{theorem}[\leanname{weak\_harnack}]
\label{thm:wH}
For $d \geq 3$, $A$ normalized elliptic on $\Bz 1$, $u > 0$ a supersolution,
and $0 < q < 1$:
\begin{equation*}
  \Bigl(\int_{\Bz{1/4}} u^{q d/(d-2)}\,dx\Bigr)^{(d-2)/(qd)}
    \leq \Bigl(\frac{C_{wH}(d)}{(1-q)^{\beta(d)}}\Bigr)^{\Lambda_A^{1/2}}
        \essinf_{\Bz{1/4}} u.
\end{equation*}
\end{theorem}

\begin{proof}
By Theorem~\ref{thm:wH-chain},
$I^\alpha \leq C_{\mathrm{chain}}(d)/(1-q)^{d\chi}\cdot (\essinf u)^{p_0}$
where $I := \int_{\Bz{1/4}} u^{qd/(d-2)}$ and $\alpha := p_0(d-2)/(qd)$.
Raising both sides to the $1/p_0$ power (\leanname{rpow\_le\_rpow\_of\_exponent}):
\[
  I^{(d-2)/(qd)}
   \;\leq\; \Bigl(\frac{C_{\mathrm{chain}}(d)}{(1-q)^{d\chi(d)}}\Bigr)^{1/p_0}
       \essinf_{\Bz{1/4}} u.
\]
Now use $1/p_0 = \Lambda_A^{1/2}/c'(d)$ and the laws of real powers:
\begin{align*}
  \Bigl(\frac{C_{\mathrm{chain}}(d)}{(1-q)^{d\chi(d)}}\Bigr)^{1/p_0}
   &= \Bigl(\bigl(C_{\mathrm{chain}}(d)/(1-q)^{d\chi(d)}\bigr)^{1/c'(d)}\Bigr)^{\Lambda_A^{1/2}}\\
   &= \Bigl(C_{\mathrm{chain}}(d)^{1/c'(d)}\big/(1-q)^{d\chi(d)/c'(d)}\Bigr)^{\Lambda_A^{1/2}}\\
   &= \Bigl(C_{wH}(d)\big/(1-q)^{\beta(d)}\Bigr)^{\Lambda_A^{1/2}}.
\end{align*}
\end{proof}

\begin{theorem}[\leanname{weak\_harnack\_on\_ball}]
\label{thm:wH-on-ball}
On a general ball $\Ball R{x_0}$ with $R > 0$, the weak Harnack inequality
takes the form
\begin{multline*}
  \Bigl(R^{-d}\int_{\Ball{R/4}{x_0}} u^{qd/(d-2)}\,dx\Bigr)^{(d-2)/(qd)}\\
   \leq \Bigl(\frac{C_{wH}(d)}{(1-q)^{\beta(d)}}\Bigr)^{\Lambda_A^{1/2}}
        \essinf_{z \in \Bz{1/4}}\, u(x_0 + R z).
\end{multline*}
\end{theorem}

\begin{proof}
Apply Theorem~\ref{thm:wH} to the rescaled function
$u_R(z) := u(x_0 + R z)$ on $\Bz 1$ with the rescaled coefficients
$A_R := \leanname{rescaleNormalizedCoeffToUnitBall}\,A$, which preserves
$\Lambda$ (so $p_0(A_R) = p_0(A)$). The change of variables
$\leanname{crossover\_integral\_comp\_affine\_ball}$ yields the
$R^{-d}$ Jacobian factor on the left, while the essential infimum of $u_R$
on $\Bz{1/4}$ equals that of $u$ on $\Ball{R/4}{x_0}$ pulled back through the
affine map.
\end{proof}

\subsubsection{Support constants}

\begin{definition}[\leanname{localMoserBase}, \leanname{localWeakHarnackBase}]
\label{def:local-bases}
\begin{align*}
  M_{\mathrm{loc}}(d) &:= C_{\mathrm{Moser}}(d)\,((d-1))^d\,4^d, \\
  W_{\mathrm{loc}}(d) &:= C_{wH,0}(d)\, d^{\beta(d)}.
\end{align*}
\end{definition}

\begin{lemma}[\leanname{one\_le\_localMoserBase}, \leanname{one\_le\_localWeakHarnackBase}]
For $d \geq 3$, $1 \leq M_{\mathrm{loc}}(d)$ and $1 \leq W_{\mathrm{loc}}(d)$.
Both are strictly positive.
\end{lemma}

\begin{proof}
Each is a product of factors $\geq 1$: $C_{\mathrm{Moser}}(d) \geq 1$,
$(d-1)^d \geq 1$ (since $d-1 \geq 2$), $4^d \geq 1$, $C_{wH,0} \geq 1$,
$d^{\beta(d)} \geq 1$ (since $d \geq 1$ and $\beta(d) \geq 0$ by
\leanname{weak\_harnack\_decay\_exp\_nonneg}).
\end{proof}

\subsubsection{Measure-bound utilities}

The Harnack and Holder layers also use a small toolkit of $\esssup/\essinf$
arithmetic on restricted ball measures (\leanname{Support/MeasureBounds.lean}).

\begin{lemma}[\leanname{le\_essInf\_real\_of\_ae\_le}, \leanname{essSup\_le\_of\_ae\_bdd}, \leanname{essInf\_add\_const\_of\_bdd}, \leanname{essSup\_neg\_of\_bdd}]
\label{lem:meas-bd}
Let $\mu$ be a non-zero (Borel) measure on $E$ and $u : E \to \R$.
\begin{enumerate}
\item If $u \geq c$ $\mu$-a.e., then $c \leq \essinf_\mu u$.
\item If $K \leq u \leq C$ $\mu$-a.e., then $\essinf_\mu u, \esssup_\mu u
  \in [K,C]$.
\item If $u$ is bounded $\mu$-a.e., then
  $\essinf_\mu(u+c) = \essinf_\mu u + c$ and the symmetric statement for
  $\esssup$.
\item $\esssup_\mu(-u) = -\essinf_\mu u$ when $u$ is bounded.
\item For a measurable embedding $T : E\to E$ and measurable $g$,
  $\essinf_\mu(g\circ T) = \essinf_{T_*\mu} g$ \leanname{(essInf\_comp\_measurableEmbedding\_eq)}.
\item For $c \in \ENNReal\setminus\{0\}$, $\essinf_{c\cdot \mu}f =
  \essinf_\mu f$.
\end{enumerate}
Furthermore, $\vol(\Ball r c)\neq 0$ for $r > 0$, so the restriction
$\vol|_{\Ball r c}$ is non-zero \leanname{(restrict\_ball\_ne\_zero)}.
\end{lemma}

\begin{proof}
Each statement is a direct unfolding of $\essinf$/$\esssup$ as a $\sSup$/$\sInf$
over levels at which the corresponding sublevel set has positive/full measure.
For instance, in (1): assume $\mu \neq 0$ and $\mu(\{u < c\}) = 0$. Then
$\{u < a\}$ has measure zero for every $a < c$, so $a$ is dominated by every
upper bound used in the definition of $\essinf_\mu u$, hence $c \leq
\essinf_\mu u$. Statements (3),(4) follow from the corresponding
$\liminf/\limsup$ arithmetic on the filter $\ae \mu$ together with the fact
that the functions are eventually bounded. Statements (5),(6) follow from
the change-of-variables and scaling formulas for the underlying
$\sSup/\sInf$ over level sets.
\end{proof}

%% --- Section 11 ---
\subsection{Harnack Inequality and Holder Regularity}

\begin{definition}[\leanname{C\_harnack}]
\label{def:Charnack}
\[
  C_{\mathrm{Har}}(d) := 17\Bigl(\frac{d-2}{d-1}\,\log\bigl(C_{\mathrm{Moser}}(d)\,(d-1)^d\,4^d\bigr)
        + d + \log\bigl(C_{wH,0}(d)\,d^{\beta(d)}\bigr)\Bigr).
\]
\end{definition}

\begin{lemma}[\leanname{C\_harnack\_nonneg}]
For $d \geq 3$, $0 \leq C_{\mathrm{Har}}(d)$.
\end{lemma}

\begin{proof}
$(d-2)/(d-1) \geq 0$, both logs are nonnegative since
$M_{\mathrm{loc}}(d) \geq 1$ and $W_{\mathrm{loc}}(d) \geq 1$, and $d \geq 0$.
\end{proof}

\begin{theorem}[\leanname{harnack}]
\label{thm:harnack}
Let $d \geq 3$, $A$ normalized elliptic on $\Bz 1$, and $u > 0$ a (local
weak) solution of $-\dv(A\nabla u) = 0$ on $\Bz 1$. Then
\begin{equation*}
  \esssup_{\Bz{1/2}} u \;\leq\; \exp\bigl(C_{\mathrm{Har}}(d)\,\Lambda_A^{1/2}\bigr)
        \essinf_{\Bz{1/2}} u.
\end{equation*}
\end{theorem}

\begin{proof}
Define
\[
K(A) := \bigl(M_{\mathrm{loc}}(d)\Lambda_A^{d/2}\bigr)^{1/p_{\mathrm{Har}}(d)}
        \cdot W_{\mathrm{loc}}(d)^{\Lambda_A^{1/2}},
\]
where $p_{\mathrm{Har}}(d) = (d-1)/(d-2) \cdot$ stuff (the auxiliary
exponent \leanname{harnack\_p}). The forward Moser estimate gives, on each
ball $\Ball{1/16}{c}$ centred at $c \in \Ball{1/2}{0}$,
\[
  \esssup_{\Ball{1/16}{c}} u
   \;\leq\; \bigl(M_{\mathrm{loc}}(d)\Lambda_A^{d/2}\bigr)^{1/p_{\mathrm{Har}}(d)}
        \cdot \Bigl(\fint_{\Ball{1/8}{c}} u^{p_{\mathrm{Har}}(d)}\Bigr)^{1/p_{\mathrm{Har}}(d)},
\]
and Theorem~\ref{thm:wH-on-ball} (weak Harnack on $\Ball{1/8}{c}$) gives
\[
  \Bigl(\fint_{\Ball{1/8}{c}} u^{p_{\mathrm{Har}}(d)}\Bigr)^{1/p_{\mathrm{Har}}(d)}
        \;\leq\; W_{\mathrm{loc}}(d)^{\Lambda_A^{1/2}}\,
        \essinf_{\Ball{1/16}{c}} u.
\]
Concatenating gives the local pointwise (a.e.) bound
$u(x) \leq K(A)\,\essinf_{\Ball{1/8}{c}} u$ for almost every $x \in \Ball{1/16}{c}$
(\leanname{ae\_le\_localHarnack\_on\_eighthBall}).

By compactness of $\overline{\Bz{1/2}}$, there is a finite cover of
$\overline{\Bz{1/2}}$ by balls $\Ball{1/16}{c_i}$ with $c_i = (15/16) x_i$,
$x_i \in \overline{\Bz{1/2}}$. Chaining at most $16$ such local Harnack
inequalities along a path connecting any two centres
(\leanname{quarterBallInf\_chain\_le}) yields, almost everywhere on
$\Bz{1/2}$,
\[
  u(x) \;\leq\; K(A)^{17}\,\essinf_{\Bz{1/2}} u.
\]
Hence $\esssup_{\Bz{1/2}} u \leq K(A)^{17}\,\essinf_{\Bz{1/2}} u$.

Finally, the elementary identity
\[
  17\Bigl[\frac{d-2}{d-1}\log M_{\mathrm{loc}}(d) + \frac{d}{2}\log\Lambda_A
    + \log W_{\mathrm{loc}}(d) \Lambda_A^{1/2}\Bigr]
  \;\leq\; C_{\mathrm{Har}}(d)\,\Lambda_A^{1/2}
\]
(using $\Lambda_A^{1/2} \geq 1$ and $\log\Lambda_A \leq 2\Lambda_A^{1/2}$,
established in \leanname{localHarnackConstant\_pow\_seventeen\_le\_exp})
gives $K(A)^{17} \leq \exp(C_{\mathrm{Har}}(d)\,\Lambda_A^{1/2})$, completing
the proof.
\end{proof}

\begin{corollary}[\leanname{harnack\_of\_homogeneousWeakSolution}]
The same conclusion holds for any positive homogeneous weak solution
$u$ on $\Bz 1$, since such a $u$ is in particular an
$\IsSolution$.
\end{corollary}

\subsubsection{Holder regularity}

\begin{definition}[\leanname{C\_holder\_Moser}, \leanname{moserDecayAlpha}]
\label{def:CholderMoser}
\begin{align*}
  C_{\mathrm{Hol}}(d) &:= 256\bigl(|C_{\mathrm{Har}}(d)| + M_{\mathrm{loc}}(d)
        + C_{\mathrm{Moser}}(d) 2^d + C_{wH,0}(d) d^d + 8\bigr), \\
  \alpha(A) &:= \exp\bigl(-(C_{\mathrm{Har}}(d)\,\Lambda_A^{1/2})\bigr).
\end{align*}
\end{definition}

\begin{lemma}[\leanname{moserDecayAlpha\_pos}, \leanname{moserDecayAlpha\_le\_one}, \leanname{moserLowerBound\_le\_decayAlpha}]
$0 < \alpha(A) \leq 1$, and
$\exp(-(C_{\mathrm{Hol}}(d)\,\Lambda_A^{1/2})) \leq \alpha(A)$.
\end{lemma}

\begin{proof}
$\alpha(A) > 0$ as the exponential of a real, $\alpha(A) \leq 1$ since the
exponent is $\leq 0$ (because $C_{\mathrm{Har}} \geq 0$), and the lower
bound is the monotonicity of $\exp$ together with $C_{\mathrm{Har}}(d) \leq
C_{\mathrm{Hol}}(d)$ (\leanname{C\_harnack\_le\_C\_holder\_Moser}).
\end{proof}

\begin{lemma}[\leanname{one\_sub\_moserDecayAlpha\_le\_two\_rpow\_neg}]
$1 - \alpha(A) \leq (1/2)^{\alpha(A)}$.
\end{lemma}

\begin{proof}
Set $\alpha = \alpha(A) \in (0,1]$. We have $1 - \alpha \leq e^{-\alpha}$
(by $1+x \leq e^x$ at $x = -\alpha$), and $-\alpha \leq -\alpha\log 2$ since
$\log 2 \leq 1$. Hence $e^{-\alpha} \leq e^{-\alpha\log 2} = (1/2)^\alpha$.
\end{proof}

\begin{theorem}[\leanname{holder\_Moser}]
\label{thm:holder-Moser}
Let $d \geq 3$, $A$ normalized elliptic on $\Bz 1$, $p_0 > 1$, and $u$ a
solution of $-\dv(A\nabla u) = 0$ with $|u|^{p_0} \in \Lp{1}(\Bz 1)$.
Then there exists a representative $v$ of $u$ on $\Bz 1$ (i.e.\ $v = u$
almost everywhere on $\Bz 1$) and an exponent
$\alpha = \alpha(A) > 0$ with
$\exp(-(C_{\mathrm{Hol}}(d)\,\Lambda_A^{1/2})) \leq \alpha$ such that for
every $x, y \in \Bz{1/2}$,
\begin{equation*}
  |v(x) - v(y)|
   \;\leq\; C_{\mathrm{Hol}}(d)\,\Lambda_A^{d/(2p_0)}
        \Bigl(\frac{p_0}{p_0 - 1}\Bigr)^{d/p_0}
        \Bigl(\int_{\Bz 1} |u|^{p_0}\,dz\Bigr)^{1/p_0}
        \,\|x - y\|^{\alpha}.
\end{equation*}
\end{theorem}

\begin{proof}
Define $\alpha := \alpha(A)$ and let $v := \leanname{moserRepresentative}\,u$
be the dyadic-average representative supplied by
\leanname{Holder/Representative.lean}. By
\leanname{moserRepresentative\_ae\_eq}, $v = u$ almost everywhere on $\Bz 1$.

Fix $x, y \in \Bz{1/2}$. If $x = y$ the bound is trivial. Otherwise,
distinguish two cases.

\emph{Large distance.} If $\|x - y\| \geq 1/8$, by
\leanname{moserRepresentative\_large\_distance\_bound} there is a direct
control by the $L^{p_0}$ norm of $u$:
\[
  |v(x) - v(y)| \;\leq\; C_{\mathrm{Hol}}(d)\,\Lambda_A^{d/(2p_0)}
    \bigl(p_0/(p_0-1)\bigr)^{d/p_0}\,\|u\|_{L^{p_0}(\Bz 1)}.
\]
Since $\|x-y\|^\alpha \geq (1/8)^\alpha \geq 1/8$ and absorbing constants
into $C_{\mathrm{Hol}}$ via the factor $256$ in
Definition~\ref{def:CholderMoser} gives the claim.

\emph{Small distance.} If $\|x - y\| < 1/8$, choose $n$ with
$2^{-(n+1)} \leq \|x - y\| < 2^{-n}$ \leanname{(exists\_moserDyadicRadius\_near)},
and apply the dyadic oscillation decay
\leanname{moser\_oscillation\_decay\_on\_ball}, which iterates the Harnack
inequality on each ball $\Ball{2^{-k}}{c}$ centred near $x$:
\[
  \osc_{\Ball{2^{-k-1}}{c}} u
    \;\leq\; (1 - \alpha)\,\osc_{\Ball{2^{-k}}{c}} u.
\]
Iterating gives geometric decay
$\osc_{\Ball{2^{-n}}{c}} u \leq (1-\alpha)^n\,\osc_{\Ball 1 0} u$.
By the previous lemma $1 - \alpha \leq (1/2)^\alpha$, hence
$(1-\alpha)^n \leq 2^{-\alpha n} \leq C\,\|x - y\|^\alpha$.

The starting oscillation $\osc_{\Bz 1} u$ is controlled by twice the local
$L^\infty$ bound, which by the forward Moser estimate
\leanname{linfty\_subsolution\_Moser\_on\_ball} is bounded by
$M_{\mathrm{loc}}(d)\Lambda_A^{d/(2p_0)} (p_0/(p_0-1))^{d/p_0}
        \|u\|_{L^{p_0}(\Bz 1)}$.
This yields the desired pointwise Holder bound on the representative $v$.
\end{proof}

\begin{corollary}[\leanname{holder\_Moser\_of\_homogeneousWeakSolution}]
The same conclusion holds verbatim for homogeneous weak solutions
($\IsHomogeneousWeakSolution$), since these satisfy
$\IsSolution$.
\end{corollary}

\begin{remark}
The dependence on the ellipticity bound $\Lambda_A$ enters Holder regularity
in two ways: through the multiplicative prefactor $\Lambda_A^{d/(2p_0)}$
inherited from the forward Moser bound, and through the Holder exponent
$\alpha(A) = \exp(-C_{\mathrm{Har}}(d)\,\Lambda_A^{1/2})$ inherited from the
Harnack constant. As $\Lambda_A \to \infty$, $\alpha(A) \to 0$ exponentially,
quantifying the well-known degeneration of De Giorgi--Nash--Moser regularity
for highly anisotropic elliptic operators.
\end{remark}

\end{document}
